\documentclass[11pt]{article}
\usepackage[margin=2.6cm]{geometry}
\usepackage{booktabs}
\usepackage{graphicx}
\usepackage{siunitx}
\sisetup{group-digits=integer}
\usepackage{amsmath}
\usepackage{hyperref}
\usepackage{seqsplit}
\graphicspath{{figures/}}
\setlength{\emergencystretch}{3em}
\newcommand{\runid}[1]{\texttt{\small\seqsplit{#1}}}
\newcommand{\dsid}[1]{\texttt{\small #1}}
\newcommand{\best}[1]{\textbf{#1}}
\newcommand{\Rp}{\ensuremath{R'}}
\newcommand{\Cp}{\ensuremath{C'}}

\title{Next-Track Session Recommendation: Experiment Synthesis}
\author{}
\date{\today}

\begin{document}
\maketitle

\begin{abstract}
This document is the living synthesis of the \emph{next-track} sequence
recommendation campaigns; the markdown reports in \texttt{experiments/}
(filenames containing \texttt{nexttrack}) are the primary record. The task is
next-item \emph{retrieval}: given a listening \emph{session} (an ordered
sequence of tracks), rank the catalog and recover the next distinct track the
user actually played, scored on a single leak-free chronological
\emph{next-distinct} split (\num{7154} sessions / \num{19402} items;
\num{5723} train / \num{1431} test; cold-item rate \num{0.55}). The headline
metric is Recall@10, reported beside mean reciprocal rank (MRR) and the
graded-relevance diagnostics artist@10 / genre@10 / artist\_mrr. Twenty-one campaigns and one research memo
are consolidated. A session GRU alone only \emph{ties} the first-order Markov
co-listening bar (R@10 \num{0.1069}) on Recall@10 and loses on MRR, and neither
loss-tuning nor any network topology (depth, width, residual, direction) beats
it. The breakthrough is \emph{rank fusion}: an R+M z-blend (GRU $\times$ Markov)
beats the bar by $\sim$$+61\%$ Recall@10 (\num{0.1726}, three-seed confirmed),
because the two legs rank \emph{different} correct tracks; a learned XGB stacker
underperforms the fixed z-blend and adding redundant neural legs does not help.
The one genuinely additive new signal is \emph{content}: an item-representation
sweep found latent \emph{dimension} (not the compressor) to be the lever, with a
peak at PCA-192, and on that rich space a three-leg R+M+C z-blend
(GRU $\times$ Markov $\times$ content-kNN) is the current best-on-record row at
Recall@10 \num{0.1859}, MRR \num{0.1025}. A subsequent literature-fit campaign
then superseded that crown: a supervised \emph{learned content projection}
(\Rp+M+\Cp), which fits a linear metric map of the PCA-192 item space toward
next-track adjacency by InfoNCE and retrieves the GRU and content legs in it,
reaches Recall@10 \num{0.2117}, MRR \num{0.1220} (paired $\Delta$ over R+M+C
$+\num{0.0259}$, CI above zero), is \emph{seed-robust} across three seeds and
\emph{cold-reaching}, and hardens on a second disjoint 0.64-cold split (paired
$\Delta$ $+\num{0.0426}$, CI above zero) --- resolving the two thin caveats
(determinism-only confirm, $+\num{0.0007}$ margin) that had bounded the R+M+C
crown. Unlike every earlier ranking predictor it is registered, promoted and
\emph{servable} (\texttt{blend-gru-markov-content-proj}, best-models rank 1),
closing the productization gap; an item-vector representation exploration and a
continuous ``sounds-alike'' \emph{music@10} metric corroborate the champion
without dethroning it. A final measurement campaign adds a display-only
\emph{session-holisticness} metric suite (artist\_adj@10 / mood\_coh@10 /
ild@10) that fingerprints the champion as the \emph{most} artist-eager model
measured (artist\_adj@10 \num{0.610}) --- with the eagerness sourced from its
Markov bigram leg, not the RNN, and the GRU and LSTM statistically tied --- and
finds an eval-time MMR re-rank ($\lambda\approx\num{0.9}$) a near-free
holisticness lever while a training-time anti-eager regularizer is weak; nothing
is promoted and the primary metric stays Recall@10. Two final architecture
campaigns then close the remaining architecture axes, and both are nulls: a
jointly-trained \emph{dual-tower} GRU over paired causal views of the prefix
fails to beat a parameter-matched single tower --- the fourth independent
learned-combiner null, and the one that answers the ``train the fusion
end-to-end'' rebuttal --- and a learned per-step \emph{pre-encoder} in front of
the recurrence is not merely useless but actively \emph{harmful} in a single
tower ($-\num{0.0140}$ Recall@10, CI below zero), a loss traced by a
counterfactual interpretability read to half-wave rectification of a
zero-centred latent deleting exactly its negative half
(\texttt{frac\_exact\_zero} \num{0.5327} $=$ \texttt{frac\_negative\_coords}
\num{0.5327}). A measured power analysis finally reframes all five architecture
nulls: at $n=\num{1431}$ the paired-bootstrap CI half-width is \num{0.0124} on
Recall@10 while a whole architecture family spans only \num{0.028}, so large
harms were detected reliably while moderate wins were never detectable at all ---
and \emph{holisticness@10} exists in three mutually non-comparable code
generations, which is why this synthesis keeps its board denominated in
Recall@10. A final measurement campaign then changes the
\emph{surface} rather than the model: the shipped product serves an
\emph{autoregressive walk}, while holisticness@10 scores a one-shot top-10 whose
\texttt{mood\_coh} anchors to a \emph{growing} prefix centroid, so journey drift
was invisible by construction (and the composite is $\num{5.5}\times$ more
sensitive to its relevance factor than to its eagerness factor). Four
walk-level metrics calibrated against the user's own listening (\num{73632} real
consecutive transitions, median step cosine \num{0.261}, \SI{27.0}{\percent} of
them negative) show the shipped engine walking $\num{2.5}\times$ tighter than
the listener it imitates (median step \num{0.661}), and three deterministic
retrieval-time controls --- a seed anchor, a stride term, and a hard per-artist
cap that \emph{cannot be outscored} (two disjoint splits:
$\Delta$artist\_conc $-\num{0.527}$ with $\Delta$music@10 $+\num{0.026}$, both
CI clear of zero) --- make a real-listening stride and a held vibe
simultaneously reachable. On \num{40} identical held-out sessions under paired
bootstrap the \emph{dual-tower} family that three campaigns had refuted on
recall is the winner on the stated objective (\texttt{latent/cummean}: stride
error $-\num{0.2545}$, genres $+\num{1.50}$, drift $+\num{0.0351}$, all CI
clear of zero, vibe \emph{held}), while pushing the single GRU or the champion
to the same stride \emph{costs} vibe --- but \texttt{ground} straddles zero for
every arm, so only journey \emph{shape} is established, and the primary metric
is unchanged. A closing architecture campaign then follows that reversal up on its
own stated mechanism and closes the axis: a \emph{bank} of $N$ parallel bare-GRU
towers under one linear head, dosed from $N=2$ to $N=4$ at a matched
\num{3.00}$\times$ parameter budget, wins on neither surface --- every arm ties
the h256 baseline on one-shot recall while \emph{five of eight lose} to a
parameter-matched single GRU, the two arms that earn a walk trade at the primary
cell both flip to \texttt{vibe} below zero at the deployed cell, and the
fixed-reference \texttt{cummean} towers the dose hypothesis was built on are the
\emph{worst} arms in the batch. The decisive datapoint is a zero: four h211
towers and one h512 GRU return the identical \num{186} hits of \num{1431},
paired $\Delta$ \num{0.000000}, so topology and width are interchangeable at
this budget. The batch's one CI-clear result is a by-product --- at
byte-identical parameters dropout \num{0.1} beats dropout \num{0.0} by
$+\num{0.0049}$ recall@10 --- which matters because the
\texttt{fusion\_layers}${}=0$ path constructs \emph{zero} dropout modules, so
every f0 arm on record, including the deployed engine, trained
under-regularized.
Six further campaigns close the score-time and depth questions. An eval-time
\emph{MMR} $\lambda$ frontier, run natively on both anchors, is the \textbf{first
confirmed crown lever on record} --- three single-GRU arms clear the prior all-time
holisticness bar at a recall whose paired CI straddles zero, with three of the four
crown factors (diversity, eagerness, \emph{and} graded relevance) improving at once
--- and it resolves its pre-registered mechanism question unanimously: sonic
diversification de-concentrates \emph{artists} $\numrange{1.59}{3.06}\times$ faster
than it reduces seed-echo in 8 of 8 arms, so an explicit artist penalty is not the
indicated sequel. A purpose-built \emph{rollout} metric (\texttt{hold@}$N$,
calibrated on \num{2135} real train sessions with an anti-gaming term inside the
metric) then returns a sixth architecture null in the deployment regime, and finds
that \emph{per-seed} variance dominates model identity by two orders of magnitude
--- which is a product recommendation (serve an ensemble) rather than a modelling
one. The largest single recall gain in the whole arc then comes from neither:
\textbf{\texttt{markov\_gate}}, a nine-line score-time change that makes the Markov
leg \emph{abstain} on the \SI{55.6}{\percent} of queries where it has no bigram
evidence, because the blend's scale-invariant normalizer had been re-inflating an
evidence-free back-off to a full one-third vote. It buys $+\num{0.02027}$
$[+\num{0.01048}, +\num{0.03075}]$ Recall@10 on the canonical split and
$+\num{0.01957}$ on a disjoint \SI{64}{\percent}-cold one (R@10 \num{0.2320},
\num{332} of \num{1431} --- the best on record), with the two warm evidence buckets
\emph{byte-for-byte identical}, which is a mechanical proof that no global static
leg weight could have found it and which narrows the standing learned-combiner
pitfall to \emph{learned score-level fusion} rather than per-query evidence gating.
Composed with MMR it is also the first crown arm that is recall-neutral against the
champion at $\num{3.6}\times$ its holisticness. Against that, the \emph{training}-time
anti-eager regularizer is refuted with a mechanism --- its margins were calibrated
from the item$\to$item transition tail while the hinge acts on
prediction$\to$current, whose ceiling sits below the quantiles used, so four arms
are \emph{provably inert} (loss traces bit-identical to the control at $\beta$ up to
the schema cap) and the arms that do fire destroy relevance before they help. And a
stage-stack depth campaign closes the serial branch of the architecture axis on
\emph{both} surfaces, with the sharpest instance yet of a dissociation that
should change how future searches are run: its best-fitting arm holds the
\emph{lowest} validation loss in the batch and retrieves \num{55} fewer hits.
\end{abstract}

\section{Introduction}
\label{sec:intro}
This document is the living synthesis of the next-track sequence-recommendation
campaigns; the markdown reports in \texttt{experiments/} are the primary
record. Each section transcribes (never recomputes) the quantitative results of
a campaign, preserves its run and dataset identifiers verbatim for
traceability, and adds the rationale and interpretation the terse reports omit.
This instance was cloned from a sibling taste-fit / rotation lab and retargeted
to a sequence ranking goal; the rotation campaigns are a separate line of work
and are \emph{not} covered here.

\paragraph{Task and split.} The lab predicts, within a listening
\emph{session}, which track the user plays \emph{next} --- a next-item
retrieval task. The model consumes the ordered session prefix and scores every
catalog item; candidates are ranked and the held-out next \emph{distinct} track
is the target (a leave-last-out evaluation that excludes prefix items, matching
the Markov parity convention). All work runs on one sequence artifact,
\dsid{seq-20260713-014226} (\num{7154} sessions / \num{19402} items), split
\emph{chronologically by session start} into \num{5723} train / \num{1431} test
at cut 2024-08-24. The defining property of this split is its \textbf{cold-item
rate of \num{0.55}}: more than half of test-session target tracks never appear
in any training session. This is the leak-free forward reality, and it caps
absolute recall by construction --- every result is read against the incumbent
baseline bar, never an absolute number.

\paragraph{Metrics.} The headline metric is \textbf{Recall@10} (R@10, the
fraction of test sessions whose held-out next track appears in the top ten
retrieved), reported beside \textbf{mean reciprocal rank} (MRR, the average
inverse rank of the true next track) and Recall@20. Since 2026-07-14 the
evaluation also emits first-class \emph{graded-relevance} metrics over the same
prefix-excluded ranking (Section~\ref{sec:graded}): \textbf{artist@$k$} (a
top-$k$ candidate shares the true next track's artist), \textbf{genre@$k$}
(likewise on primary genre), and \textbf{artist\_mrr} (MRR of the first
same-artist candidate). All confidence intervals are bootstrap 95\% CIs over
\num{2000} resamples of the \num{1431} test sessions, and the decisive
statistic is the \emph{paired} per-session $\Delta$ against a named reference.
A board \emph{win} requires beating the reference on \emph{both} R@10 and MRR
beyond non-overlapping marginal CIs (the strict bar), then a reproduction
confirm; the paired-$\Delta$ CI-entirely-above-zero test is the more powerful
alternative statistic (it crowned both champions) and can disagree with the
strict bar. The graded metrics are the honest scorecard: exact recall is modest
by construction, but the champion lands the right \emph{artist} in the top ten
$\sim$32\% of the time and the right \emph{genre} $\sim$40\% --- far more useful
than the $\sim$17\% exact number the cold-item rate imposes. Since 2026-07-20 the
leaderboard also carries \textbf{music@10}, a continuous ``sounds-alike'' graded
metric (Section~\ref{sec:music}): the mean best cosine, in a balanced intrinsic
musical-distance space, between the true next track and the top-ten retrieved
candidates --- a smooth superset of Recall@10 that credits landing something that
\emph{sounds like} the truth on an exact miss. It is display-only; the primary
metric stays Recall@10.

\paragraph{The standing champion.} The current best-on-record \emph{row} ---
and, since 2026-07-18, a \emph{registered, promoted and servable} model --- is a
supervised \emph{learned content projection}, \emph{\Rp+M+\Cp}: a session GRU
$\times$ a first-order Markov co-occurrence model $\times$ a content-kNN leg,
where the GRU (\Rp) and the content leg (\Cp) retrieve in an InfoNCE metric map
of the PCA-192 item space fit toward next-track adjacency, at R@10 \num{0.2117},
MRR \num{0.1220} (Sections~\ref{sec:projection}--\ref{sec:projection-scan}). It
supersedes the three-leg \emph{R+M+C} z-blend on the same PCA-192 space (R@10
\num{0.1859}; Section~\ref{sec:phase2}), which had itself dethroned the two-leg
\emph{R+M} z-blend on the 64-dim song-autoencoder space (AE-64, R@10
\num{0.1726}; Section~\ref{sec:blend}). The projection is \emph{seed-robust} and
\emph{two-split hardened}, and it is deployable via
\texttt{POST /api/models/blend-gru-markov-content-proj/predict} --- closing the
productization gap the earlier off-server fusion drivers left open. The prior
champions remain useful references (R+M is still the leanest two-leg blend), but
they are no longer the deployable anchor. One qualification has since been added:
the best \emph{recall} figure on record is no longer the champion as promoted but
the same model with its Markov leg \emph{gated} to abstain where it has no bigram
evidence (R@10 \num{0.2320}; Section~\ref{sec:markovgate}), a nine-line score-time
change carrying a promotion recommendation that has not yet been executed.

\section{Background: the sequence task, the baselines, and the metrics}
\label{sec:background}
This section gives the minimum theory for a reader who knows machine learning
but not this repo. It is keyed to the project's living facts file and grows only
when a new predictor family or metric enters the reports.

\paragraph{The models and baselines in play.}
\begin{description}
  \item[Session GRU / LSTM (\texttt{seq-nexttrack}).] A recurrent neural network
  that reads the ordered session prefix one item at a time (each item
  represented by its item-latent vector) and emits a predicted next-item latent;
  candidates are ranked by cosine similarity to that prediction. The GRU (gated
  recurrent unit) and LSTM (long short-term memory) are the two recurrent cells
  tried. The recurrence captures \emph{session context} --- the trajectory of
  the listening session --- rather than a single previous item.
  \item[Training loss: InfoNCE versus cosine.] The sequence model can be trained
  to regress the next-item latent by \emph{cosine} loss (pull the prediction
  toward the true next latent) or by \textbf{InfoNCE}, a contrastive objective
  that pulls the prediction toward the true next item and pushes it away from a
  batch of negatives. InfoNCE is the single most important lever in the whole
  arc: it is $\approx 2.3\times$ cosine on Recall@10, because ranking retrieval
  rewards a representation that \emph{separates} the true next item from
  distractors, which a contrastive loss optimizes directly and a pointwise
  regression does not.
  \item[First-order Markov co-occurrence (the bar, \texttt{B2}).] A train-only
  transition model: a \texttt{track\_bigram} counting ``what track follows what''
  with artist/genre back-off for unseen pairs (the transition term the sibling
  app already ships). It directly \emph{memorizes} adjacency, which is exactly
  what exact next-track ranking rewards, and it is the reference bar every arm is
  measured against. A second-order Markov (trigram) back-off was also tried and
  is essentially neutral --- sessions are too short/sparse for a trigram to fire,
  so it falls through to the bigram.
  \item[item2vec.] A skip-gram (SGNS) co-listening embedding trained on the
  \emph{train} sessions only (a leak-safe analogue of word2vec, where ``words''
  are tracks and ``sentences'' are sessions); candidates are ranked by cosine to
  the mean-pooled prefix. On \num{5723} train sessions it is too sparse to
  localize and is the weakest leg on record; the full-history co-listening
  vectors are deliberately \emph{not} used because they leak across the split.
  \item[content-kNN.] A training-free retrieval leg: cosine of each candidate's
  item latent to the prefix latents (max-aggregated). It is weak standalone on
  exact recall but captures \emph{vibe} similarity the counts miss (its genre@10
  already edges the Markov bar), and it is the one genuinely additive third leg
  for the fusion blend.
  \item[Rank-fusion z-blend versus a learned stacker.] Two ways to combine legs.
  The \textbf{z-blend} normalizes each leg's candidate scores (z-score) and takes
  a weighted sum ($\alpha\cdot z(\text{leg})+\dots$) --- a fixed, parameter-light
  fuser. The \textbf{XGB stacker} is a gradient-boosted meta-learner over
  per-candidate features (leg scores $+$ popularity $+$ recency $+$
  artist/genre-match), trained pointwise with sampled negatives. Across two
  campaigns the fixed z-blend \emph{beats} the learned stacker decisively,
  because the stacker's pointwise binary training does not optimize the ranking
  objective as well as directly fusing normalized scores.
  \item[Learned content projection (supervised metric map).] A linear map $W$ of
  the item-latent space fit by \textbf{InfoNCE} over the train consecutive pairs
  $(u\!\to\!v)$ (self-loops dropped, Markov convention), so cosine in $W$-space
  reflects ``likely to follow.'' It is leak-free (train pairs $+$ intrinsic
  content latents only, never the split). Both the content-kNN leg (\Cp) and the
  GRU (\Rp) then retrieve in the projected space; Markov is unchanged. It is the
  one \emph{supervised} alternative to the unsupervised, variance-ordered PCA
  geometry, and --- because content is the only channel with cold reach --- it
  lifts \emph{both} legs that share the space, making it the current champion
  (Section~\ref{sec:projection}).
  \item[Dual-tower recurrent fusion (\texttt{seq-dualgru}).] Two
  \emph{independent} recurrent towers (separate weights) read the prefix in
  parallel over configurable strictly-causal \emph{views} of it ---
  \texttt{latent} (the raw item vector), \texttt{delta} (the step-to-step
  difference) or \texttt{cummean} (the running session centroid) --- and their
  concatenated per-step hidden states are read out by a fusion head
  (\texttt{fusion\_layers}${}=0$ for a pure linear read, ${}=1$ for a 256-wide
  MLP). It reuses the \texttt{seq-nexttrack} objective, batching, early stopping
  and evaluation verbatim, so it isolates architecture as the only variable. The
  motivating idea is that two towers should help only if they see genuinely
  different information; Section~\ref{sec:dualtower} tests that and refutes it.
  \item[Per-step pre-encoder.] A nonlinear re-embedding placed \emph{in front of}
  the recurrence --- $[\,\text{Linear},\ \text{ReLU},\ \text{Dropout}\,]$
  applied at each timestep, exposed as
  \texttt{pre\_hidden}\,/\,\texttt{pre\_layers}. It is an unusually sharp test
  object: a GRU's own input transform is \emph{already} linear, so a linear
  pre-encoder at width $\ge$ the input dimension is exactly as expressive as no
  pre-encoder at all, and the \emph{rectifier} is the only change --- which is
  what makes Section~\ref{sec:preencoder}'s null mechanistically attributable
  rather than merely negative.
\end{description}

\paragraph{Item-latent spaces: song-AE versus PCA, and the wrong-collection
trap.} The sequence model predicts \emph{into}, and retrieves \emph{over}, a
per-item latent space, so the choice of that space is a first-class axis. Three
families appear:
\begin{description}
  \item[AE-64 (\dsid{spotify\_tracks\_song\_ae}).] A 64-dim nonlinear
  \emph{song-autoencoder} latent, the original champion item space. An
  autoencoder compresses each item by reconstructing it through a bottleneck;
  its equal-weight-per-block reconstruction loss spreads capacity evenly across
  content, acoustic and catalog blocks.
  \item[PCA-$k$ (\dsid{spotify\_tracks\_song\_pca$k$}).] A \emph{linear} principal
  components compression at dim $k \in \{32,64,128,192,256\}$, which orders
  dimensions by explained variance and decorrelates them. As
  Sections~\ref{sec:compression}--\ref{sec:dim} show, this variance-ordered,
  decorrelated geometry is \emph{easier} for the GRU to regress into than the
  AE's entangled equal-weight latents, so PCA beats the AE at every dim $\ge 64$.
  \item[The wrong-collection trap.] Both of the above are fit \emph{only} from
  the leak-safe \dsid{spotify\_tracks\_content} source (intrinsic
  content/acoustic/catalog fields; no behavioral or target field). The default
  \dsid{spotify\_tracks} 200-dim collection is a \emph{behavioral} co-listening
  embedding that soft-leaks replay behavior and must \textbf{never} be used as an
  item space.
\end{description}

\section{Next-track leaderboard}
\label{sec:leaderboard}
Table~\ref{tab:lb} is the standing best-on-record board. It is on the single
leak-free next-distinct split and sorted by Recall@10. Since 2026-07-18 the
sequence family carries a ranking serving contract, so the champion
learned-projection blend is a \emph{promoted, servable} model
(\texttt{blend-gru-markov-content-proj}, best-models rank 1), not merely a board
row (Section~\ref{sec:projection-scan}); the earlier blends below it, and the leg
standalones, remain off-server references. The item spaces are all fit from the
leak-safe content source (Section~\ref{sec:background}). Since 2026-07-22 the
evaluation also carries a \emph{display-only session-holisticness} suite
(artist\_adj@10 / mood\_coh@10 / ild@10; Section~\ref{sec:holisticness}); it does
not reorder this board (the primary metric stays Recall@10), but it fingerprints
the champion as the most artist-eager row. Since 2026-07-25 the board also
carries three \emph{refuted} architecture datapoints --- the widest single-GRU
capacity control, the best dual-tower arm and the best single-tower pre-encoder
arm (Sections~\ref{sec:dualtower}--\ref{sec:preencoder}). They are listed
because they \emph{bound} their axes, not because they are candidates: none
beats the single-GRU baseline by the paired statistic and none was registered.
Since 2026-08-02 the board carries two more such bounding rows, the widest
single-GRU point (h512) and the best parallel-tower-bank arm, which return the
\emph{identical} hit count at a matched parameter budget --- neither is
significant against the h256 baseline, and together they close the tower-bank
axis (Section~\ref{sec:towerbank}).
Since 2026-07-27c the board's top row is the \emph{evidence-gated} champion
(Section~\ref{sec:markovgate}): making the Markov leg abstain on the
\SI{55.6}{\percent} of queries where it has no bigram evidence lifts Recall@10 to
\num{0.2320} (\num{332} of \num{1431}), paired CI above zero on two disjoint
splits, with MRR and both graded metrics up. It is the best recall figure on
record, and it is a promotion \emph{recommendation} only --- nothing was registered
or promoted --- so the servable champion below it is unchanged.
On the server-effective primary metric, and on why this board stays denominated
in Recall@10 across three holisticness generations, see
Section~\ref{sec:power-generations}.

\begin{table}[ht]
  \centering
  \caption{Next-track best-on-record board, on the leak-free next-distinct split
  (\num{1431} test sessions; bootstrap 95\% CI over \num{2000} resamples). R@10,
  MRR, artist@10, genre@10, music@10 higher-better; the first-order Markov model is
  the incumbent reference bar. music@10 values are the offline backfill
  (Section~\ref{sec:music}); rows without a mapped run are left blank. Champion
  row and prior champions bolded; best cell per metric bolded. The three refuted
  architecture rows added 2026-07-25/26 are shown to \emph{bound} their axes and
  are excluded from the best-cell bolding; they carry no marginal CI, only a
  paired $\Delta$ against the single-GRU baseline (Tables~\ref{tab:dualtower}
  and~\ref{tab:preencoder}). The top row, added 2026-07-27c, is the evidence-gated
  champion of Section~\ref{sec:markovgate}: it is a promotion \emph{recommendation}
  rather than a registered model, so its R@10 is bolded as the best on record but it
  is otherwise excluded from the best-cell bolding, and it carries a paired $\Delta$
  (Table~\ref{tab:markovgate-paired}) rather than a marginal CI.}
  \label{tab:lb}
  {\footnotesize\setlength{\tabcolsep}{3.2pt}\begin{tabular}{lcccccl}
    \toprule
    Model & R@10 [95\% CI] & MRR & art@10 & gen@10 & mus@10 & Note \\
    \midrule
    \best{\Rp+M+\Cp\ $+$ gate} (best recall) & \best{0.2320} & 0.1335 & 0.367 & 0.463 & 0.493 & not promoted \\
    \best{\Rp+M+\Cp, PCA-192} (champion, servable) & \best{0.2117 [0.191, 0.233]} & \best{0.1220} & 0.333 & 0.401 & 0.454 & new champion \\
    R+M+C, PCA-192 (prior fusion champ.)  & 0.1859 [0.166, 0.207] & 0.1025 & 0.331 & 0.394 & 0.449 & content additive \\
    R+M+C, PCA-128                        & 0.1824 [0.162, 0.203] & 0.0993 & 0.332 & 0.396 & --- & content additive \\
    \best{R+M, AE-64} (prior champion)    & 0.1726 [0.154, 0.192] & 0.0963 & 0.322 & 0.397 & 0.441 & prod.\ anchor; 3-seed \\
    R+M+C, AE-64 (graded-preferred)       & 0.1740 [0.154, 0.194] & 0.1014 & \best{0.330} & 0.400 & --- & content only tied \\
    R+M, PCA-192                          & 0.1712 [0.151, 0.192] & 0.0952 & 0.329 & 0.398 & 0.443 & $\approx$ prior champion \\
    R+M, PCA-128                          & 0.1670 [0.148, 0.187] & 0.0920 & 0.329 & 0.398 & 0.444 & GRU gain not carried \\
    R+A+M, AE-64                          & 0.1684 & 0.0943 & --- & --- & --- & ANN leg redundant \\
    GRU single h512, PCA-192          & 0.1300 & --- & --- & --- & 0.458 & width tie \#7; not sig.\ \\
    tower bank lat$\times$4 h211, PCA-192 & 0.1300 & --- & --- & --- & 0.459 & bank $=$ width; closed \\
    LSTM single, PCA-128                  & 0.1272 [0.110, 0.145] & 0.0495 & 0.335 & \best{0.458} & 0.451 & best single; $\ll$ blend \\
    GRU single h454, PCA-192          & 0.1251 & 0.0539 & 0.338 & 0.454 & 0.457 & width closed \\
    GRU single, PCA-192                   & 0.1230 [0.106, 0.140] & 0.0534 & \best{0.338} & 0.450 & 0.456 & dim point-optimum \\
    dual-tower GRU f0, PCA-192        & 0.1188 & 0.0563 & 0.337 & 0.456 & 0.455 & dual tower: refuted \\
    GRU single, PCA-128 (base)            & 0.1146 [0.099, 0.132] & 0.0489 & 0.333 & 0.452 & \best{0.457} & confirmed base space \\
    GRU single, PCA-256                   & 0.1167 [0.101, 0.134] & 0.0502 & 0.338 & 0.458 & 0.456 & past the dim peak \\
    pre-MLP$\to$GRU, PCA-192          & 0.1090 & 0.0446 & 0.324 & 0.469 & 0.450 & pre-encoder harmful \\
    B2 first-order Markov (bar)           & 0.1069 [0.092, 0.124] & 0.0694 & 0.298 & 0.389 & 0.407 & app incumbent \\
    GRU single, AE-64 (control)           & 0.0957 [0.081, 0.111] & 0.0411 & 0.294 & 0.450 & 0.432 & bake-off control \\
    ANN pooled-MLP, PCA-128               & 0.0901 [0.076, 0.105] & 0.0394 & 0.299 & 0.435 & 0.426 & loses to the GRU \\
    XGB stacker (M+G / M+G+A)             & 0.080--0.087 & 0.042--0.049 & 0.31 & 0.38 & 0.41 & combiner underperforms \\
    content-kNN (song-AE)                 & 0.0636 [0.050, 0.077] & 0.0321 & 0.272 & 0.405 & --- & the additive C leg \\
    item2vec (train-only SGNS)            & 0.0517 & 0.0214 & 0.166 & 0.312 & --- & drags the blend \\
    B0 popularity / B1 recency            & 0.0007 & --- & --- & --- & 0.185 & floors \\
    \bottomrule
  \end{tabular}}
\end{table}

\section{The sequence sweep: a GRU-infonce ties the Markov bar}
\label{sec:sweep}
\emph{Rationale.} The instance's opening campaign asked the founding question:
does a session-context sequence model --- a GRU/LSTM over the ordered prefix,
predicting the next track's 64-dim song-AE latent and retrieving by cosine ---
beat the first-order Markov transition model the sibling app already ships,
\emph{forward}, on the leak-free session split? Four arms crossed the recurrent
cell (GRU, LSTM), capacity (\texttt{hidden} 128, 256) and, decisively, the
training loss (cosine versus InfoNCE), at 60 epochs.

\emph{Results.} No arm beat the bar. The best arm (M3, GRU \texttt{hidden}128
InfoNCE) reached a statistical \emph{tie} with the Markov bar on R@10
(\num{0.0894}, paired-$\Delta$ CI includes zero) but \emph{lost} on MRR
(\num{0.0388} versus \num{0.0694}) and Recall@20. The dominant lever was the
loss objective: InfoNCE was $\approx 2.3\times$ cosine on R@10 (\num{0.0894}
versus \num{0.0391}), while capacity (\texttt{h256}) and the LSTM cell barely
moved the cosine arms (all $\sim$0.038--0.041). The popularity and recency
floors sat near zero (R@10 \num{0.0007}).

\emph{Interpretation.} A latent-plus-single-cosine-hop is a harder route to the
\emph{exact} next track than a co-listening bigram: the bigram directly
memorizes ``what follows what,'' while the GRU must land near the right point in
a 64-d content space and then retrieve, losing on precise ranking (MRR). The
finding is consistent with a taste-drift reality --- \num{0.55} cold items,
identity/adjacency-driven behavior --- and it set the agenda: InfoNCE is the
base loss, and the natural next arm is a \emph{blend} of the GRU's candidate
scores with the Markov bigram.
\emph{Source: \texttt{2026-07-13-nexttrack-sequence-sweep.md}.}

\section{Topology scan: depth, width and direction are a spent lever}
\label{sec:topology}
\emph{Rationale.} With the loss fixed at InfoNCE, this scan isolated network
\emph{topology}: does depth (stacked layers), width (\texttt{hidden} 256), an
inter-layer residual stack or bidirectionality lift the sequence model over the
Markov bar? Because a bidirectional RNN under per-step teacher forcing would leak
future items, the bidirectional arm used a leak-free last-item objective (encode
the full prefix, predict only the held-out next item).

\emph{Results.} Refuted. No arm beat the bar on both metrics. The best arm
(width \texttt{h256}, R@10 \num{0.0957}) is a \emph{tie} on R@10 and loses on
MRR (\num{0.0411} versus \num{0.0694}). Depth is flat 1$\to$2 layers
(\num{0.0894}$\to$\num{0.0922}, in-band) and \emph{hurts} at 3 layers
(\num{0.0825}, out-of-band below the bar); the residual stack over depth-2/h256
added nothing. The bidirectional arm collapsed (\num{0.0482}), but this is an
\emph{objective} artifact, not a directionality verdict: the leak-free last-item
objective supplies $\sim$1 training target per session versus teacher forcing's
$\sim$13 next-step pairs, so it trained on $\sim$13$\times$ less signal.

\emph{Interpretation.} Architecture is not the lever for next-track on this
corpus. The transferable latent-retrieval signal caps at R@10 $\sim$0.09--0.10
across depth, width, residual and direction, below the co-listening bigram ---
consistent with the whole arc that the loss objective moved the needle and
capacity/topology do not. The one remaining unexplored lever was the candidate
framing / fusion, which the next campaign pulled decisively.
\emph{Source: \texttt{2026-07-13-nexttrack-topology-scan.md}.}

\section{The GRU$\times$Markov fusion champion, and why more ensembling fails}
\label{sec:blend}
\emph{Rationale.} The sweep and the topology scan agreed that the best single
GRU only \emph{ties} the Markov bar on R@10 while losing on MRR. This campaign
asked whether \emph{combining} models beats the bar, and how far ensembling
goes. Three legs were defined, each emitting full-vocab candidate scores:
\textbf{M} = the train-only Markov bigram (the bar), \textbf{R} = the
GRU-infonce-h256 recurrence, and \textbf{A} = an ANN bag-of-prefix MLP (a
feed-forward twin of R). They were combined by a weighted z-blend and, as a
learned alternative, an XGB stacker.

\emph{Results.} The \textbf{R+M z-blend} ($\alpha{=}0.5$) \emph{beats the bar
decisively}: R@10 \num{0.1726} $[0.154, 0.192]$, MRR \num{0.0963}, a paired
$\Delta$ over Markov of $+\num{0.0657}$ $[+0.049, +0.083]$ (CI entirely above
zero) --- $\sim$$+61\%$ R@10, $+39\%$ MRR --- on a robust plateau across
$\alpha \in [0.25, 0.75]$ and reproduced across three seeds (\num{0.1705} /
\num{0.1726} / \num{0.1726}). It was registered as the \texttt{seq-blend}
predictor and the \texttt{blend-gru-markov} definition, the first next-track
champion. Further ensembling then \emph{failed} to improve on it
(Table~\ref{tab:blend}): R+A+M (\num{0.1684}) is slightly worse, R+A ($\approx$
GRU alone, \num{0.0894}) is much worse, and both XGB stacker variants (M+R
\num{0.1020}; M+R+A \num{0.0887}) barely clear Markov alone.

\begin{table}[ht]
  \centering
  \caption{The fusion campaign: the R+M z-blend beats the Markov bar
  decisively, and no richer ensemble improves on it. R@10 and MRR
  higher-better; champion bolded.}
  \label{tab:blend}
  {\small\setlength{\tabcolsep}{5pt}\begin{tabular}{lccl}
    \toprule
    Model & R@10 [95\% CI] & MRR & Note \\
    \midrule
    \best{R+M z-blend $\alpha{=}0.5$} (champion) & \best{0.1726 [0.154, 0.192]} & \best{0.0963} & $\Delta$ vs Markov $+0.066$ [$+0.05,+0.08$] \\
    R+A+M z-blend                                & 0.1684 & 0.0943 & ANN leg redundant with GRU \\
    B2 Markov bar (M)                            & 0.1069 [0.092, 0.123] & 0.0694 & the reference \\
    XGB stacker M+R                              & 0.1020 & 0.0561 & learned combiner $\ll$ z-blend \\
    GRU alone (R)                                & 0.0957 & 0.0411 & ties the bar on R@10 only \\
    R+A z-blend                                  & 0.0894 & 0.0377 & $\approx$ GRU alone \\
    XGB stacker M+R+A                            & 0.0887 & 0.0572 & barely clears Markov \\
    \bottomrule
  \end{tabular}}
\end{table}

\emph{Interpretation.} The win is about signal \emph{diversity}, not model
count. The GRU (latent similarity) and Markov (co-occurrence counts) rank
\emph{different} correct next-tracks, so fusing them roughly doubles the hit
rate; piling on more of the \emph{same} kind of signal --- another neural latent
predictor (A), redundant with R --- adds nothing. The learned XGB stacker
\emph{underperforms} the simple z-blend, counter to the usual ``stacking beats
blending'' intuition: its pointwise binary-with-negatives training does not
optimize the ranking objective as well as directly fusing normalized scores. The
remaining headroom is bounded by the \num{0.55} cold-item rate and requires
\emph{genuinely new signal}, not more model fusion --- the hypothesis the next
campaign tests.
\emph{Source: \texttt{2026-07-14-nexttrack-blend-and-ensemble.md}.}

\section{Graded-relevance metrics and new base learners: content is the only additive signal}
\label{sec:graded}
\emph{Rationale.} Two threads in one campaign. (A) Make the evaluation credit
landing the right \emph{band/vibe} even when the exact track is wrong --- the
observation that ``the first choice is usually not the exact track, but the
artist is usually right'' --- by promoting \emph{artist@k}, \emph{genre@k} and
\emph{artist\_mrr} to first-class metrics wired end-to-end
(\texttt{seq\_common.eval\_from\_scores} $\to$ Rust \texttt{Metrics} /
\texttt{MetricsSpec} $\to$ \texttt{domain.toml} $\to$ UI). (B) Try genuinely new
leak-free base learners --- \textbf{C} = content-kNN, \textbf{M2} = a
second-order Markov trigram with back-off, and \textbf{I} = a train-only
item2vec SGNS --- scored on both the exact and the new graded metrics.

\emph{Results.} The exact-metric path was byte-identical (purely additive keys);
the champion reproduced R@10 \num{0.1726} / MRR \num{0.0963} exactly. No new base
learner beat the bar standalone (M2 \num{0.1076} $\approx$ M \num{0.1069} --- the
trigram falls through to the bigram; item2vec \num{0.0517} is the weakest leg).
Blended into the champion, \textbf{content is the only additive signal}:
\textbf{R+M+C} (adding the content leg) is the top blend on MRR
(\num{0.0963}$\to$\num{0.1014}), artist@10 (\num{0.3215}$\to$\num{0.3298}) and
artist\_mrr, and ties the champion on exact R@10 (paired $\Delta$
$+\num{0.0014}$, CI straddles zero). Item2vec \emph{drags} the blend (R+M+I
\num{0.1349}); the trigram is neutral.

\emph{Interpretation.} Content-kNN ranks \emph{different} near-correct tracks
(latent vibe similarity) than the GRU (session recurrence) and Markov
(co-occurrence), exactly the diversity the fusion campaign predicted new signal
would need. On the AE-64 space the content leg only \emph{tied} the champion on
exact recall while lifting the graded metrics, so it was carried forward as the
graded-preferred third leg, not a new exact champion. That ``tie'' becomes a
\emph{win} once the item space is richer (Section~\ref{sec:phase2}): on PCA-192
the same R+M+C construction is the strongest on record. Co-listening item2vec and
the trigram Markov added nothing on this data.
\emph{Source: \texttt{2026-07-14-nexttrack-graded-relevance-and-new-legs.md}.}

\section{The item-representation compression bake-off}
\label{sec:compression}
\emph{Rationale.} The champion's item space --- a 64-dim AE latent --- was one
arbitrary early choice, and the sequence GRU predicts \emph{into} and retrieves
\emph{over} whatever geometry that space imposes. This campaign asks which
\emph{compression} is the best space for the GRU: linear PCA versus the
nonlinear autoencoder, each at latent dim 32 / 64 / 128. Six alternative spaces
were baked from the \emph{same} leak-safe 384-dim content source
(\dsid{spotify\_tracks\_content}), each fit unsupervised over all \num{23529}
items using only intrinsic content/acoustic/catalog fields, then scored under a
\emph{frozen} GRU readout (\texttt{gru-infonce-h256}), so the item latent space
is the only axis that moves. The design bet was that the AE's
equal-weight-per-block reconstruction loss would \emph{preserve} the small
acoustic/catalog blocks that PCA's text-dominated global variance would starve;
the null was that dim, not the compressor, sets the ceiling.

\emph{Results.} The parity gate passed: the AE-64 control reproduced its
on-record R@10 \num{0.0957}. All six spaces trained cleanly
(Table~\ref{tab:compression}; Figure~\ref{fig:compression}). \textbf{PCA-128} is
the single-GRU winner at R@10 \num{0.1146}, a paired $\Delta$ of $+\num{0.0189}$
$[+0.0042, +0.0335]$ over the AE-64 control --- a CI \emph{entirely above zero},
the only arm that clears the control significantly --- and it also tops MRR,
artist@10 and artist\_mrr. PCA-64 is the runner-up (\num{0.1041}, tie). The two
32-dim spaces sit far below the 64-dim control (both CIs below zero). The
training-free content-kNN probe is flat at $\sim$0.06 across all six spaces and
does \emph{not} track the GRU ranking.

\begin{table}[ht]
  \centering
  \caption{The item-representation compression bake-off: six item spaces (linear
  PCA vs the nonlinear AE $\times$ latent dim 32 / 64 / 128), each scored under
  the \emph{same} frozen \texttt{gru-infonce-h256} readout (\num{1431} test
  sessions; paired-$\Delta$ vs the in-batch AE-64 control). R@10, MRR, artist@10,
  genre@10 higher-better. Winner (PCA-128) bolded; best cell per metric bolded.}
  \label{tab:compression}
  {\footnotesize\setlength{\tabcolsep}{3pt}\begin{tabular}{lp{2.05cm}ccccc}
    \toprule
    Space & Run id & R@10 [95\% CI] & $\Delta$ vs AE-64 & MRR & art@10 & gen@10 \\
    \midrule
    \best{PCA-128} (winner) & \runid{run-20260715-042100-a1d9d} & \best{0.1146 [0.099, 0.132]} & \best{$+0.019$ [$+0.004,+0.034$]} & \best{0.0489} & \best{0.3326} & 0.4521 \\
    PCA-64 (runner-up)      & \runid{run-20260715-042100-8851d} & 0.1041 [0.088, 0.121] & $+0.008$ (tie) & 0.0416 & 0.3117 & 0.4598 \\
    AE-64 (control)         & \runid{run-20260715-042059-dd261} & 0.0957 [0.081, 0.111] & --- (control) & 0.0411 & 0.2942 & 0.4500 \\
    AE-128                  & \runid{run-20260715-042059-e96ae} & 0.0929 [0.078, 0.108] & $-0.003$ (tie) & 0.0398 & 0.3026 & \best{0.4640} \\
    AE-32                   & \runid{run-20260715-042059-d3f9e} & 0.0545 [0.043, 0.066] & $-0.041$ & 0.0274 & 0.2376 & 0.3955 \\
    PCA-32                  & \runid{run-20260715-042059-b1242} & 0.0412 [0.031, 0.052] & $-0.055$ & 0.0224 & 0.2187 & 0.3816 \\
    \bottomrule
  \end{tabular}}
\end{table}

\begin{figure}[ht]
  \centering
  \includegraphics[width=0.8\textwidth]{nexttrack_compression}
  \caption{Latent dimension, not the compressor, is the lever. Both series climb
  steeply from dim 32 to 64; PCA (blue) then keeps climbing to 128 and
  \emph{overtakes} the AE, while the AE (grey) plateaus --- its 128-dim point
  actually dips below its own 64-dim control. At dim 32 the AE edges PCA, a
  low-capacity crossover that reverses once dim is generous. The AE-64 control
  (red ring) is the champion's item space; the PCA-128 winner (green star) is the
  only arm whose paired-$\Delta$ CI clears zero. Error bars are bootstrap 95\%
  CIs; the neighbouring CIs overlap, so the win rests on the \emph{paired} test.}
  \label{fig:compression}
\end{figure}

The single-GRU win does \textbf{not} survive into the production blend:
re-fusing the winning PCA-128 GRU leg with the Markov leg scores R@10
\num{0.1670}, \emph{below} the AE-64 champion blend (\num{0.1726}) with
overlapping CIs. Nothing was promoted; PCA-128 was carried forward only as the
sanctioned Phase-2 base space.

\emph{Interpretation.} Latent dimensionality is the dominant lever and the
compressor is secondary. The design's acoustic-preservation hypothesis is
\emph{refuted}: at generous dim the linear PCA beats the nonlinear AE, and
although AE-32 preserves the numeric/acoustic blocks near-perfectly it ranks near
the bottom. What helps is a \emph{variance-ordered, decorrelated} target geometry
that the GRU finds easier to regress into than the AE's entangled equal-weight
latents. The AE's edge at dim 32 is a purely low-capacity phenomenon that
reverses once dim is generous. The flat content-kNN probe that fails to track the
GRU ranking is a flagged divergence: the win comes from the \emph{learnability}
of PCA's target geometry, not a demonstrably richer space. And the load-bearing
caution is the blend result: a clearly stronger GRU leg bought \emph{nothing} at
the blend, because the Markov leg is representation-independent and the extra
tracks the better GRU recovers overlap what Markov already supplies.
\emph{Source:
\texttt{2026-07-15-nexttrack-compression-representation-bakeoff.md}.}

\section{The dimension ceiling: PCA peaks at 192}
\label{sec:dim}
\emph{Rationale.} With dimension established as the lever and PCA still climbing
at 128, this leg extended the PCA sweep upward (PCA-192, PCA-256) to locate the
real optimum before Phase 2 committed a base space, and added AE-192 to test
whether PCA $\ge$ AE holds at higher dimension. The protocol was identical (same
frozen \texttt{gru-infonce-h256} readout, same parity gate); the reference is
PCA-128 (\runid{run-20260715-042100-a1d9d}).

\emph{Results.} The PCA R@10-versus-dimension curve \textbf{peaks at $\sim$192}
(R@10 \num{0.1230}) and \emph{turns over} by 256 (\num{0.1167}); it was not still
climbing. The full trend is $32\!\to\!\num{0.041}$, $64\!\to\!\num{0.104}$,
$128\!\to\!\num{0.115}$, $\mathbf{192\!\to\!\num{0.123}}$, $256\!\to\!\num{0.117}$
(Figure~\ref{fig:dim}). PCA-192 is the point-optimum, but its paired $\Delta$
versus PCA-128 is $+\num{0.0084}$ $[-0.0007, +0.0189]$ --- it misses the ``CI
entirely above zero'' bar by \num{0.0007}, a whisker --- so \textbf{PCA-128
stands as the Phase-2 base}, with PCA-192 flagged as the confirm-and-promote
candidate. PCA $\ge$ AE holds strongly and the gap is widest at 192 (PCA-192
\num{0.123} versus AE-192 \num{0.098}); the AE plateaus flat at $\sim$0.09--0.10
from dim 64 upward.

\begin{figure}[ht]
  \centering
  \includegraphics[width=0.8\textwidth]{nexttrack_dim_curve}
  \caption{The dimension ceiling is found. Extending the compression sweep of
  Figure~\ref{fig:compression} to PCA-192/256 and AE-192, the PCA curve (blue)
  peaks at $\sim$192 (R@10 \num{0.123}) and turns over by 256, while the AE (grey)
  plateaus flat at $\sim$0.09--0.10 and never catches it --- the PCA$\,\ge\,$AE gap
  is widest at 192. PCA-192 is the point-optimum (green star) but its paired
  $\Delta$ over PCA-128 grazes zero, so PCA-128 stays the Phase-2 single-model
  base. Error bars are bootstrap 95\% CIs over the \num{1431} test sessions.}
  \label{fig:dim}
\end{figure}

\emph{Interpretation.} Extra variance past $\sim$192 (PCA-256 captures cumulative
explained-variance ratio \num{0.997}) is low-signal / noise dimensions the GRU
cannot exploit and that slightly dilute the target geometry. The AE's flat
plateau reconfirms the refuted acoustic-preservation hypothesis at higher
dimension: its equal-weight reconstruction objective keeps spending capacity on
the tiny numeric/acoustic blocks (reconstruction $R^2 \approx 0.999$ at 192)
rather than the high-variance text/catalog structure the GRU needs. As before,
the training-free content-kNN probe stays flat and does not reproduce the GRU's
peak, so the dimension optimum is about target \emph{learnability}, not intrinsic
retrieval richness.
\emph{Source: \texttt{2026-07-15b-nexttrack-compression-dim-extension.md}.}

\section{The model-family sweep and the new champion, R+M+C on PCA-192}
\label{sec:phase2}
\emph{Rationale.} With the compression axis settled (PCA-128 base, PCA-192
candidate), this campaign swept the model \emph{family} on the chosen spaces:
does a different recurrent cell (LSTM), a non-recurrent pooled MLP (ANN), or a
\emph{learned} XGB stacker beat the fixed GRU$\times$Markov z-blend --- and does
adding the content-kNN third leg (R+M+C) on the \emph{rich} PCA space finally
beat the champion? A frozen held config ran across all arms; only the model
family and the representation varied, and the server was serialized
(\texttt{-{}-max-runs 1}) to avoid a thread-oversubscription throughput trap.

\emph{Results.} A \textbf{new best-on-record row}: the \textbf{R+M+C z-blend on
PCA-192} (GRU $\times$ train-only Markov $\times$ content-kNN, equal thirds)
reaches R@10 \num{0.1859} $[0.166, 0.207]$, MRR \num{0.1025}, a paired $\Delta$
over the prior champion of $+\num{0.0133}$ $[+0.0007, +0.0259]$ --- entirely
above zero, with no MRR regression (Table~\ref{tab:phase2};
Figure~\ref{fig:phase2}). The content leg is genuinely additive on the rich
space, beating same-space R+M by $+\num{0.0154}$ (PCA-128) / $+\num{0.0147}$
(PCA-192), both CIs above zero --- \emph{stronger} than the AE-64 R+M+C that
only tied (Section~\ref{sec:graded}), because the richer, higher-EVR space makes
the content leg pay off. The learned \textbf{XGB stacker badly underperforms} the
fixed z-blend (\num{0.080}--\num{0.087}, losing even to its own single GRU base
leg); the \textbf{LSTM is the best single model} (PCA-128 \num{0.1272}, paired
$\Delta$ over the GRU $+\num{0.0126}$, CI above zero) but far below the blends;
and the Phase-2 single base \textbf{stays PCA-128}.

\begin{table}[ht]
  \centering
  \caption{Phase-2 model-family sweep on the rich PCA spaces, ranked by R@10
  (higher-better). The fixed z-blend is the combiner of record; the learned XGB
  stacker underperforms it badly. Champion row and prior champion bolded; best
  cell per metric bolded. Paired $\Delta$ is per-session against the named
  reference.}
  \label{tab:phase2}
  {\footnotesize\setlength{\tabcolsep}{4pt}\begin{tabular}{lp{2.4cm}ccp{2.9cm}}
    \toprule
    Model & Run id / source & R@10 [95\% CI] & MRR & paired $\Delta$ \\
    \midrule
    \best{R+M+C, PCA-192} (champ.\ row) & \runid{analyze\_rmc} (A4) & \best{0.1859 [0.166, 0.207]} & \best{0.1025} & vs champ $+0.0133$ [$+0.001,+0.026$] \\
    R+M+C, PCA-128                       & \runid{analyze\_rmc} (A2) & 0.1824 [0.162, 0.203] & 0.0993 & vs champ $+0.0098$ (tie) \\
    \best{R+M, AE-64} (prior champion)   & \runid{run-20260714-170928-645ae} & 0.1726 [0.154, 0.192] & 0.0963 & --- (reference) \\
    R+M, PCA-192                         & \runid{run-20260715-142440-96f53} & 0.1712 [0.151, 0.192] & 0.0952 & vs champ $-0.0014$ (tie) \\
    R+M, PCA-128                         & \runid{run-20260715-044920-36dce} & 0.1670 [0.148, 0.187] & 0.0920 & vs champ $-0.0056$ (tie) \\
    LSTM single, PCA-128                 & \runid{run-20260715-142440-546d0} & 0.1272 [0.110, 0.145] & 0.0495 & vs GRU-128 $+0.0126$ ($>0$) \\
    GRU single, PCA-192                  & \runid{run-20260715-142440-89d6a} & 0.1230 [0.106, 0.140] & 0.0534 & vs GRU-128 $+0.0084$ (tie) \\
    ANN pooled-MLP, PCA-128              & \runid{run-20260715-142440-7e528} & 0.0901 [0.076, 0.105] & 0.0394 & vs GRU-128 $-0.0245$ \\
    XGB stacker M+G, PCA-128             & \runid{run-20260715-142440-371b1} & 0.0867 [0.073, 0.101] & 0.0493 & vs GRU-128 $-0.0280$ \\
    XGB stacker M+G+A, PCA-192           & \runid{run-20260715-142440-f111d} & 0.0853 [0.071, 0.099] & 0.0490 & vs GRU-192 $-0.0377$ \\
    \bottomrule
  \end{tabular}}
\end{table}

\begin{figure}[ht]
  \centering
  \includegraphics[width=0.85\textwidth]{nexttrack_phase2_families}
  \caption{The Phase-2 model-family sweep. The fixed z-blend is the combiner of
  record: R+M+C on PCA-192 (green) is the new champion row, the content leg is
  genuinely additive on the rich space, and the learned XGB stacker (red) badly
  underperforms --- it loses even to its own single GRU base leg. LSTM edges GRU
  as the best single model (grey) but sits far below the blends. The first-order
  Markov bar (\num{0.107}) is the incumbent reference; the champion row is at
  \num{0.186}.}
  \label{fig:phase2}
\end{figure}

\emph{The crown, and its two honest caveats.} The board-event candidate was
confirmed by reproducing the GRU-192 leg five times (the original plus four
reruns) and recomputing R+M+C each time; the paired $\Delta$ over the prior
champion was $+\num{0.0133}$ $[+0.0007, +0.0259]$ \emph{every} rerun, with the
MRR lead intact, so R+M+C on PCA-192 was recorded as the new champion row. Two
caveats travel with the crown and must not be omitted. First, \textbf{the
reproduction spread was exactly zero} because \texttt{seq-nexttrack} training is
deterministic on this CPU --- all five GRU-192 checkpoints were byte-identical
--- so the confirm verified \emph{determinism, not seed-robustness}; the only
uncertainty quantified is the single-split bootstrap. Second, \textbf{the paired
$\Delta$ lower bound is $+\num{0.0007}$}: the crown clears zero by a whisker on a
single split. The crown is real by the pre-agreed criterion, but the margin over
the prior champion is thin, and hardening it needs a second leak-free test split
(or seed-varied legs once the blend is generalized to three legs).

\emph{Interpretation.} The compression-plus-model-family thread is otherwise
closed: PCA-128 is the single-model base, the LSTM is the best single model, the
learned stacker is dead as a combiner, and the fixed z-blend is the combiner of
record. The load-bearing result is that \emph{content is the one leg that adds
genuinely new ranking signal, and it pays off more on the higher-EVR PCA space}
--- the exception to the standing rule that single-model / representation gains
do not carry to the blend (the R+M pair on the rich PCA spaces only ties the
AE-64 champion). The productization gap remains: the champion is train-only and
off-server, so the prior champion R+M on AE-64 stays the deployable anchor until
the three-leg blend is registered with a \texttt{predict} subcommand.
\emph{Source: \texttt{2026-07-15c-nexttrack-phase2-model-sweep.md}.}

\section{The learned content projection: a new champion from the literature-fit campaign}
\label{sec:projection}
\emph{Rationale.} The R+M+C champion (Section~\ref{sec:phase2}) carried two thin
caveats --- a determinism-only reproduction confirm and a paired-$\Delta$ lower
bound of only $+\num{0.0007}$ --- and its headroom was capped by the \num{0.55}
cold-item rate, which only \emph{new} signal can attack. Guided by a survey of the
established session-recommendation literature (the Jannach--Ludewig session-kNN
line, sequential rules / factorized Markov, GRU4Rec training tricks, rank fusion
and learning-to-rank, content / cold-start-aware session rec), this campaign
pre-registered three tiers of classic methods and asked which actually improve the
fit to \emph{this} problem --- a single-user, session-based, next-distinct
\emph{retrieval} task with 55\% cold targets. All arms were measured off-server on
the PCA-192 champion artifact \dsid{seq-20260715-131139}, reusing the sanctioned
driver that reproduces the champion (R+M+C \num{0.1859}) exactly; the decision
statistic is the paired per-session $\Delta$ against that champion.

\emph{Results --- the one win.} A supervised \textbf{learned content projection}
is a new best-on-record row. Fitting a linear metric map $W$ of the PCA-192
content space toward next-track adjacency (InfoNCE over train consecutive pairs,
self-loops dropped) and then retrieving \emph{both} the content leg (\Cp) and the
GRU (\Rp) in that projected space gives \textbf{\Rp+M+\Cp z-blend R@10
\num{0.2117} $[0.191, 0.233]$, MRR \num{0.1220}} --- a paired $\Delta$ over the
champion of $+\num{0.0259}$ $[+0.012, +0.041]$, entirely above zero
(Table~\ref{tab:projection}). Unlike the R+M+C crown it is \textbf{seed-robust}
(three seeds; paired $\Delta$ $+\num{0.026}$ / $+\num{0.025}$ / $+\num{0.029}$,
every CI $>0$ --- a genuine seed-varied confirm, since the projection has random
init and shuffled batches) and \textbf{cold-reaching}: the projected content leg
\Cp\ scores cold targets \num{0.152} $\approx$ warm \num{0.151}, versus Markov's
warm \num{0.226} / cold \num{0.010}, so it reaches never-seen items as well as
seen ones --- genuine content generalization, the one channel that reaches the
55\% cold majority. The symmetric single map beats the directional two-tower, and
rank-192 $=$ rank-64 (the gain is the supervision, not the capacity). As a side
effect the projected single GRU \Rp\ ($\sim$\num{0.176}) nearly matches the
\emph{old} champion blend and crushes the prior best single (LSTM \num{0.127}).

\begin{table}[ht]
  \centering
  \caption{Tier A1 of the literature-fit campaign: the learned content projection
  on PCA-192, versus the R+M+C champion (\num{1431} test sessions; bootstrap 95\%
  CI; paired $\Delta$ per session vs the champion). Winner (\Rp+M+\Cp) bolded; best
  cell per metric bolded.}
  \label{tab:projection}
  {\footnotesize\setlength{\tabcolsep}{4pt}\begin{tabular}{lccccl}
    \toprule
    Model & R@10 [95\% CI] & MRR & art@10 & gen@10 & paired $\Delta$ vs champion \\
    \midrule
    R+M+C champion (thirds)           & 0.1859 [0.166, 0.207] & 0.1025 & 0.331 & 0.394 & --- (reference) \\
    \best{\Rp+M+\Cp (both projected)} & \best{0.2117 [0.191, 0.233]} & \best{0.1220} & 0.333 & 0.401 & \best{$+0.0259$ [$+0.012,+0.041$] $>0$} \\
    R+M+C+\Cp (4-leg)                 & 0.2061 [0.185, 0.227] & 0.1148 & \best{0.345} & \best{0.410} & $+0.0203$ [$+0.013,+0.029$] $>0$ \\
    R+M+\Cp (swap C$\to$\Cp)          & 0.1971 [0.177, 0.218] & 0.1107 & 0.334 & 0.400 & $+0.0112$ [$-0.000,+0.023$] cand. \\
    \Cp\ proj-content (sym, r192)     & 0.1516 & 0.0708 & 0.305 & 0.401 & vs unsup.\ C 0.0685: $+0.083$ \\
    \Cp\ dir.\ two-tower (r192)       & 0.1349 & 0.0552 & 0.291 & 0.393 & symmetric wins \\
    \bottomrule
  \end{tabular}}
\end{table}

\emph{The lanes that closed.} Every other literature-fit direction failed to beat
the fixed z-blend, exactly as the survey predicted for a single-user, cold-heavy
corpus (Table~\ref{tab:projection-lanes}). Principled Markov smoothing
(absolute-discounting $+$ a Kneser--Ney continuation floor) is a real but tiny
standalone estimator gain (\num{0.1069}$\to$\num{0.1104}) that washes out in the
blend ($-\num{0.0021}$, CI straddles 0 --- the content leg already covers the
marginal cold mass it recovers). A ranking-aware combiner has no honest headroom:
even the \emph{oracle} weights, fit directly on the test labels, reach only
\num{0.1901} ($+\num{0.0042}$, straddles 0), and RRF rank fusion \emph{hurts}
recall ($-\num{0.029}$, CI$<$0, discarding the score magnitude the z-blend keeps).
A gated V-SKNN session-kNN is warm-only against 55\% cold targets and \emph{drags}
the blend ($-\num{0.038}$, CI$<$0) --- the single-user setting strips the
cross-user collaborative neighborhood that drives its multi-user benchmark wins.
And GRU4Rec importance negatives (a larger support-weighted negative bank) give a
within-noise standalone tick ($+\num{0.005}$) that also washes out
($-\num{0.0014}$, tie). \textbf{The only additive gain comes from new signal, not
smarter fusion.}

\begin{table}[ht]
  \centering
  \caption{The literature-fit lanes that closed: every borrowed method fails to
  beat the fixed z-blend champion (\num{1431} test sessions; blend paired $\Delta$
  per session vs the champion).}
  \label{tab:projection-lanes}
  {\footnotesize\setlength{\tabcolsep}{5pt}\begin{tabular}{llcl}
    \toprule
    Lane & Standalone R@10 & Blend $\Delta$ vs champion & Verdict \\
    \midrule
    Markov KN-smoothing (Tier A2)         & 0.1069$\to$0.1104 & $-0.0021$ (straddles 0) & null in blend \\
    Ranking combiner, \emph{oracle} (B1)  & 0.1901 (blend) & $+0.0042$ (straddles 0) & no honest headroom \\
    RRF rank fusion (B1)                   & 0.1572 (blend) & $-0.029$ (CI$<$0) & hurts recall \\
    V-SKNN session-kNN (B2)                & 0.0671 & $-0.038$ (CI$<$0) & drags (warm-only) \\
    GRU importance negatives (C)           & 0.1230$\to$0.1279 & $-0.0014$ (tie) & washes out \\
    \bottomrule
  \end{tabular}}
\end{table}

\emph{Interpretation.} The supervised content projection is the first
genuinely-additive new signal since the content leg itself, and it works for the
reason the survey flagged: content is the only channel with cold reach, and
replacing the unsupervised, variance-ordered PCA geometry with a target-fit one
lifts \emph{both} legs that share it. The flashy borrowed methods --- session-kNN,
learned combiners, GRU-side loss tricks --- are dead ends here because they either
re-shuffle the same three legs (no headroom over the z-blend, even for the oracle)
or import a collaborative assumption the single-user corpus violates. The one
movable secondary axis is graded relevance: the oracle / RRF reweights and the
projection nudge artist@10 / genre@10, a lever the exact-recall axis does not
share.
\emph{Source: \texttt{2026-07-18-nexttrack-literature-fit-campaign.md}.}

\section{Registration, the objective / rank / $\tau$ scan, and second-split hardening}
\label{sec:projection-scan}
\emph{Rationale.} The projection win landed as an off-server driver --- the same
productization gap the R+M+C crown had. This follow-through implemented the
projection as a first-class capability of the \texttt{seq-blend} predictor
(hyperparameters \texttt{projection}, \texttt{projection\_rank},
\texttt{projection\_objective}, \texttt{projection\_tau}, \dots: it fits the metric
map from train consecutive pairs, saves it as \texttt{projection.npz} in the run
dir, and bakes it into the promoted model for serving), then (1) reproduced the
win on-server and registered $+$ promoted it, and (2) ran a pre-approved objective
/ rank / $\tau$ ablation scan to confirm the default is the robust config to
promote.

\emph{Registration.} The verify run
(\runid{run-20260718-191539-1b89c-seq-blend}, definition default $=$ InfoNCE /
full-rank / $\tau$ \num{0.07}) reproduced the off-server driver to four decimals
--- R@10 \num{0.21174}, MRR \num{0.12199} --- and was promoted via the API as
\texttt{blend-gru-markov-content-proj}, the projection matrix baked into the model
directory for serving. \texttt{GET /api/best-models} confirms it at \textbf{rank
1} (\num{0.21174}), above the prior servable champion
\texttt{blend-gru-markov-content} (rank 2, \num{0.18588}). With the sequence
family's ranking serving contract this is a genuinely deployable model, servable
via \texttt{POST /api/models/blend-gru-markov-content-proj/predict}.

\emph{The scan.} The decision rule --- re-promote only a config beating the default
on \emph{both} R@10 and MRR by more than the $\sim$$\pm\num{0.015}$ noise band ---
confirms the default (Table~\ref{tab:projection-scan}). \textbf{InfoNCE $\gg$ BPR}
is the sharpest axis: both BPR cells (\num{0.1579}, \num{0.1523}) fall
$\sim$\num{0.055} below the InfoNCE default and land \emph{below even the old R+M+C
champion} (\num{0.1859}), so the pairwise BPR map produces a materially worse
geometry for both shared-latent legs. \textbf{Full-rank $\ge$ rank-64}
($-\num{0.006}$, a tie --- reconfirming ``the gain is the supervision, not the
capacity''). And \textbf{$\tau$ is mildly sensitive and monotone up} over the
tested range (\num{0.05}$\to$\num{0.1929}, \num{0.07}$\to$\num{0.2117},
\num{0.10}$\to$\num{0.2229}) --- but the nominal $\tau$-\num{0.10} top cell beats
the default by only $+\num{0.0112}$ on recall, \emph{inside} the noise band, so it
is a tie and the default \textbf{HELD}. Unlike the reweighting axis of
Section~\ref{sec:projection}, the graded metrics here \emph{track} recall (best and
worst together), because the scan moves the \emph{signal quality} rather than
re-fusing fixed legs.

\begin{table}[ht]
  \centering
  \caption{The on-server objective / rank / $\tau$ ablation scan (6 cells, sorted
  by recall@10). All overlay the promoted definition; hit@10 $=$ recall@10 by
  construction. Promoted champion (default, held) bolded; best cell per metric
  bolded (all held by the sub-noise $\tau$-\num{0.10} cell).}
  \label{tab:projection-scan}
  {\footnotesize\setlength{\tabcolsep}{4pt}\begin{tabular}{lccccp{3.0cm}}
    \toprule
    Config & recall@10 & MRR & art@10 & gen@10 & run id \\
    \midrule
    S2 infonce / full / $\tau$ \num{0.10}          & \best{0.2229} & \best{0.1273} & \best{0.3424} & \best{0.4032} & \runid{run-20260718-193111-a781c} \\
    \best{infonce / full / $\tau$ \num{0.07} (promoted, held)} & \best{0.2117} & \best{0.1220} & 0.3326 & 0.4011 & \runid{run-20260718-191539-1b89c} \\
    S3 infonce / rank 64 / $\tau$ \num{0.07}        & 0.2055 & 0.1169 & 0.3326 & 0.3990 & \runid{run-20260718-193650-66c96} \\
    S1 infonce / full / $\tau$ \num{0.05}          & 0.1929 & 0.1111 & 0.3354 & 0.3997 & \runid{run-20260718-192443-57432} \\
    S4 bpr / full / $\tau$ \num{0.07}              & 0.1579 & 0.0939 & 0.3187 & 0.3913 & \runid{run-20260718-194228-4cc06} \\
    S5 bpr / rank 64 / $\tau$ \num{0.07}            & 0.1523 & 0.0909 & 0.3194 & 0.3906 & \runid{run-20260718-194832-d070f} \\
    \bottomrule
  \end{tabular}}
\end{table}

\subsection{Second-split hardening}
\label{sec:projection-hardening}
\emph{Rationale.} The projection crown was measured on a single leak-free split.
This validation-only pass re-ran it --- plus the two reference blends --- on a
\textbf{second, disjoint, leak-free split} to settle whether the win is
split-specific. Nothing is promoted or re-promoted. The split
\dsid{seq-20260718-211238} is an \emph{earlier-holdout} 60--80\% chronological band
on the same PCA-192 space (train 4,292 [earliest 60\%] / test 1,431 [60--80\%
band]); the item space is byte-identical to the champion split, and the test
window is \textbf{disjoint in both index and time} from the champion split's
80--100\% test. Its \textbf{cold-item rate is \num{0.64}} (versus \num{0.55}), a
stronger, independent cold-reaching stress test --- so absolute R@10 lands
\emph{below} the first-split \num{0.212} by construction, and the verdict is the
paired per-session $\Delta$ on \emph{this} split, not absolute reproduction.

\begin{table}[ht]
  \centering
  \caption{Second-split hardening on the disjoint 0.64-cold earlier-holdout split
  \dsid{seq-20260718-211238} (\num{1431} test sessions; validation-only, nothing
  promoted). Point metrics (top) and the paired per-session $\Delta$ (bottom,
  2000-resample bootstrap, seed 1337). Projection champion bolded.}
  \label{tab:projection-hardening}
  {\footnotesize\setlength{\tabcolsep}{4pt}\begin{tabular}{lccccp{3.0cm}}
    \toprule
    Config & recall@10 & MRR & art@10 & gen@10 & run id \\
    \midrule
    \best{\Rp+M+\Cp projection champion} & \best{0.1705} & \best{0.1050} & \best{0.2991} & \best{0.3417} & \runid{run-20260718-212639-50e1d} \\
    R+M+C crown (primary reference)      & 0.1279 & 0.0816 & 0.2956 & 0.3410 & \runid{run-20260718-213156-cf1dd} \\
    R+M 2-leg (context reference)        & 0.1174 & 0.0791 & 0.2844 & 0.3368 & \runid{run-20260718-214114-93bad} \\
    \midrule
    \multicolumn{6}{l}{\emph{Paired $\Delta$ recall@10 (projection $-$ reference):}} \\
    PROJ vs R+M+C (primary) & \multicolumn{5}{l}{$+\num{0.0426}$ $[+0.0280, +0.0573]$ --- entirely $>0$; MRR $+\num{0.0234}$ (no regression)} \\
    PROJ vs R+M (context)   & \multicolumn{5}{l}{$+\num{0.0531}$ $[+0.0391, +0.0678]$ --- entirely $>0$} \\
    \bottomrule
  \end{tabular}}
\end{table}

\begin{figure}[ht]
  \centering
  \includegraphics[width=0.8\textwidth]{nexttrack_projection_hardening}
  \caption{The projection lead \emph{widens} on the colder split. On the canonical
  0.55-cold split the projection champion beats the R+M+C crown by paired
  $\Delta$ $+\num{0.0259}$; on a disjoint, colder 0.64-cold earlier-holdout split
  the margin \emph{grows} to $+\num{0.0426}$ (both CIs entirely $>0$). Absolute
  R@10 is lower on the colder split by construction (more targets are cold); the
  verdict is the paired $\Delta$, which is unambiguous on both. The mechanism: on
  a colder split the Markov / content-kNN legs of the crown reach fewer targets,
  and the projection's cold-reaching \Cp\ leg carries proportionally more of the
  win.}
  \label{fig:nexttrack-projection-hardening}
\end{figure}

\emph{Verdict --- crown hardened.} The decision rule (paired $\Delta$ vs R+M+C
with 95\% CI entirely $>0$ and MRR not regressing) is met: paired $\Delta$
\textbf{$+\num{0.0426}$ $[+0.0280, +0.0573]$}, MRR $+\num{0.0234}$ (no regression).
The projection's advantage is not split-specific --- in fact the margin is
\emph{wider} here ($+\num{0.0426}$) than on the first split ($+\num{0.0259}$),
consistent with the cold-reaching mechanism
(Figure~\ref{fig:nexttrack-projection-hardening}). Combined with the prior
three-seed robustness and the cold-reach decomposition, the crown is now hardened
on two disjoint leak-free splits --- resolving both caveats that travelled with the
R+M+C crown.
\emph{Source: \texttt{2026-07-18b-nexttrack-projection-registration-and-scan.md}.}

\section{Item-vector representation: is the geometry the bottleneck?}
\label{sec:representation}
\emph{Rationale.} The compression axis (Sections~\ref{sec:compression}--\ref{sec:dim})
found latent \emph{dimension} to dominate but never asked whether the
\emph{composition} of the 192-d item vector --- the equal-weighted 4-block PCA
(text / numeric / acoustic / categorical) --- is itself a handicap. An
explained-variance-ratio (EVR) preflight said it might be: the raw input variance
budget is 44\% acoustic $+$ 43\% numeric versus only 2.5\% text (text is
L2-normalized, so its per-dimension scale is tiny), so full-matrix PCA-192 spends
\textbf{87\% of its variance budget} reconstructing the numeric $+$ acoustic
blocks to $R^2 = \num{1.000}$ while the signal-carrying text block reaches only
$R^2 = \num{0.717}$; per-column standardize (whitening) flips the captured budget
to 73\% text / 21\% categorical (text $R^2$ \num{0.905})
(Figure~\ref{fig:nexttrack-evr-budget}). Eight alternative item-latent spaces
(whitening $\times$ acoustic-drop, text-only, text$+$categorical, and the
musical-distance metric geometry) were built parity-gated to the champion split
(vocab / split byte-identical; only \texttt{item\_latents.f32} differs) and each
evaluated under \emph{two} readouts: \textbf{(A)} the champion projection
\texttt{blend-gru-markov-content-proj} (the InfoNCE map re-fits on that space ---
the space's ceiling under the winning recipe) and \textbf{(C)} the frozen
\texttt{gru-infonce-h256} single (raw representation quality). Validation-only;
nothing promoted.

\begin{figure}[ht]
  \centering
  \includegraphics[width=0.78\textwidth]{nexttrack_evr_budget}
  \caption{Why the item vector is a handicap for a raw model --- and why the
  projection is the fix. The equal-weighted PCA-192 orders dimensions by raw
  variance, so 87\% of its captured budget reconstructs the behaviorally-inert
  numeric $+$ acoustic blocks (popularity / era / loudness) while the
  signal-carrying text block (2.5\% of raw input variance, L2-normalized) is
  under-preserved ($R^2$ \num{0.717}). Per-column whitening flips the captured
  budget to 73\% text / 21\% categorical (text $R^2$ \num{0.905}). The learned
  InfoNCE projection learns this same reweighting from data.}
  \label{fig:nexttrack-evr-budget}
\end{figure}

\begin{table}[ht]
  \centering
  \caption{Item-vector representation exploration: eight alternative 192-d item
  spaces vs the PCA-192 champion, under readout A (champion projection) and readout
  C (frozen GRU). Paired $\Delta$ on exact hit@10 (2000-resample bootstrap, rng
  1337) vs the same-readout baseline; a lift needs the CI entirely $>0$. Baseline
  rows and the best readout-C lift (V5) bolded.}
  \label{tab:representation}
  {\scriptsize\setlength{\tabcolsep}{3pt}\begin{tabular}{llcccc}
    \toprule
    & & \multicolumn{2}{c}{Readout A (projection)} & \multicolumn{2}{c}{Readout C (frozen GRU)} \\
    \cmidrule(lr){3-4}\cmidrule(lr){5-6}
    Variant (space) & dataset id & R@10 & $\Delta$ hit@10 [CI] & R@10 & $\Delta$ hit@10 [CI] \\
    \midrule
    \best{full --- PCA-192 (baseline)} & \dsid{seq-20260715-131139} & \best{0.2117} & (baseline) & \best{0.1230} & (baseline) \\
    V1 noaco               & \dsid{seq-20260718-222742} & 0.2089 & $-0.0028$ [$-0.013,+0.008$] & 0.1481 & $+0.0252$ [$+0.009,+0.041$] \\
    V2 std (whiten all)    & \dsid{seq-20260718-222754} & 0.2082 & $-0.0035$ [$-0.015,+0.008$] & 0.1530 & $+0.0300$ [$+0.011,+0.049$] \\
    V3 std-noaco           & \dsid{seq-20260718-222807} & 0.2166 & $+0.0049$ [$-0.008,+0.017$] & 0.1572 & $+0.0342$ [$+0.015,+0.055$] \\
    V4 textcat             & \dsid{seq-20260718-222820} & 0.1950 & $-0.0168$ [$-0.029,-0.006$] & 0.1565 & $+0.0335$ [$+0.015,+0.051$] \\
    \best{V5 balanced-pca192} & \dsid{seq-20260718-222834} & 0.2159 & $+0.0042$ [$-0.006,+0.015$] & \best{0.1586} & \best{$+0.0356$ [$+0.019,+0.052$]} \\
    V7 std-txtcatup        & \dsid{seq-20260718-222848} & 0.2055 & $-0.0063$ [$-0.019,+0.006$] & 0.1488 & $+0.0259$ [$+0.006,+0.046$] \\
    V10 text-only          & \dsid{seq-20260718-222901} & 0.1999 & $-0.0119$ [$-0.025,+0.001$] & 0.1377 & $+0.0147$ [$-0.005,+0.034$] \\
    V6 sonic-64 (64-d)     & \dsid{seq-20260718-222905} & 0.1936 & $-0.0182$ [$-0.031,-0.005$] & 0.1216 & $-0.0014$ [$-0.017,+0.013$] \\
    \bottomrule
  \end{tabular}}
\end{table}

\begin{figure}[ht]
  \centering
  \includegraphics[width=0.8\textwidth]{nexttrack_representation}
  \caption{The representation handicap is real for the \emph{raw} GRU but
  \emph{absorbed} by the champion projection. Six of eight whitening /
  acoustic-drop / metric spaces lift the frozen GRU (readout C, green; all six with
  paired-$\Delta$ CI entirely $>0$, $+\num{0.025}$ to $+\num{0.036}$), but none
  lifts the champion projection (readout A, blue; the two nominally-positive spaces
  straddle 0, and V4/V6 significantly regress). The learned InfoNCE map already
  re-weights away the loud acoustic / numeric directions, so whitening the input
  first is redundant at the champion ceiling: whitening and the projection are the
  same fix.}
  \label{fig:nexttrack-representation}
\end{figure}

\emph{Verdict --- exact recall.} \textbf{Readout A: no variant beats the PCA-192
champion with paired $\Delta$ CI entirely $>0$.} The two nominally-positive spaces
(V3 std-noaco $+\num{0.0049}$, V5 balanced-pca192 $+\num{0.0042}$) straddle zero;
two spaces are significantly \emph{worse} (V4 textcat $-\num{0.0168}$, dropping
numeric$+$acoustic costs the projection real signal; V6 sonic-64 $-\num{0.0182}$,
the expected 64-d handicap). \textbf{Readout C: six of eight variants beat the
frozen GRU with paired $\Delta$ CI entirely $>0$} ($+\num{0.025}$ to
$+\num{0.036}$; best V5 balanced $+\num{0.0356}$, V3 std-noaco $+\num{0.0342}$). So
the raw full-matrix PCA-192 geometry \emph{is} a handicap --- suppressing the loud
acoustic / numeric directions lifts the raw GRU by $\sim$20--29\% relative ---
\textbf{but the gain is fully absorbed by the projection}: whitening lifts readout
C, not readout A, because the learned InfoNCE map already re-weights away those
directions from data. Whitening and the learned projection are the \emph{same
fix}, and the champion has already pulled it
(Figure~\ref{fig:nexttrack-representation}).

\emph{Graded and music@10.} The representation levers move the \emph{graded} axis
more cleanly than exact recall: whitening lifts the frozen GRU's genre@10 to
$\sim$\num{0.47} (V7 std-txtcatup \num{0.4717}, V3 std-noaco \num{0.4710} vs
baseline \num{0.4500}), and under the projection a text-only space (V10) posts the
readout-A maxima on artist@10 (\num{0.3466}) and genre@10 (\num{0.4144}) while
trading $-\num{0.0119}$ exact recall. The continuous \textbf{music@10} metric
(Section~\ref{sec:music}) \emph{corroborates} the exact-recall verdict rather than
overturning it: the projection champion leads on it (\num{0.4538}) over the R+M+C
crown (\num{0.4491}), and under the champion projection it is flat across the eight
representations ($\sim$\num{0.447}--\num{0.460}). Re-basing the geometry moves band
/ vibe matching, but exact next-track adjacency at the champion ceiling is
immovable --- which is precisely the axis the supervised projection owns.

\emph{Interpretation.} The equal-weighted PCA-192 item vector is a genuine
handicap for a \emph{raw} model, and the learned projection is exactly the right
response to it --- this campaign independently validates \emph{why} the projection
is the champion (it is the representation fix, learned) and confirms there is no
free lunch left in re-basing the input geometry. The one caution: dropping
numeric$+$acoustic \emph{entirely} (V4 textcat) hurts the champion even though it
helps the raw GRU, so the numeric block carries something the projection uses ---
keep numeric, whiten / drop only acoustic. The most promising open direction is
composing the two fixes: fitting the projection \emph{on} a whitened / std-noaco /
metric space to see whether it reaches further or is redundant.
\emph{Source: \texttt{2026-07-18c-nexttrack-representation-exploration.md}.}

\section{The music@10 metric: the recall champion is not the musically-closest}
\label{sec:music}
\emph{Rationale.} Exact Recall@10 is unforgiving on a 55\%-cold single-user corpus,
and the artist@10 / genre@10 diagnostics are coarse category hits. This campaign
introduces a \emph{continuous} ``sounds-alike'' graded metric, \textbf{music@10},
and back-fills it across every next-track run. It is defined as the mean over test
cases of the best (max) cosine, in a \emph{balanced} intrinsic musical-distance
space (\dsid{spotify\_tracks\_content\_metric}: a block-weighted, partial-whitened
fusion of the whitened 64-d song-AE latent, the 384-d name/tag text embedding, five
metadata numerics, and a per-genre AE prototype), between the true next track and
the top-ten retrieved candidates. A top-ten that contains the exact truth scores
\num{1.0}, so music@10 is a \emph{smooth superset} of Recall@10 --- it credits
landing something that \emph{sounds like} the truth even on an exact / artist /
genre miss. It is display-only; the primary metric stays Recall@10.

\emph{Results.} music@10 was recomputed offline and non-destructively across all
53 next-track runs (from each run's stored \texttt{predictions.json} and its
dataset \texttt{items.json}, on the shared \num{1431}-case split). It
\textbf{correlates} with Recall@10 ($r = \num{0.609}$) --- a better ranker also
lands closer misses --- but it is far from redundant, and it crowns a
\emph{different} winner (Table~\ref{tab:music-families};
Figure~\ref{fig:nexttrack-music-scatter}). The Recall@10 champion is a
\texttt{seq-blend} run (\num{0.223}), but the \emph{musically-closest} predictions
come from \texttt{seq-nexttrack} --- which regresses the next track's content
latent directly and retrieves by cosine --- at best music@10 \num{0.482} and the
highest family mean (\num{0.448}). When a next-track model misses the exact track,
it still tends to surface something that \emph{sounds} right; the
\texttt{seq-popularity} floor is correctly worst (music@10 \num{0.185}), confirming
the metric is not trivially saturated.

\begin{table}[ht]
  \centering
  \caption{music@10 by predictor family, from the offline backfill over all 53
  next-track runs (one shared \num{1431}-case split). The recall@10 champion
  (\texttt{seq-blend}) is \emph{not} the music@10 champion (\texttt{seq-nexttrack}).
  Best cell per column bolded.}
  \label{tab:music-families}
  {\footnotesize\setlength{\tabcolsep}{6pt}\begin{tabular}{lcccc}
    \toprule
    Predictor family & runs & best recall@10 & best music@10 & mean music@10 \\
    \midrule
    seq-blend      & 23 & \best{0.223} & 0.461 & 0.441 \\
    seq-nexttrack  & 24 & 0.159 & \best{0.482} & \best{0.448} \\
    seq-markov     & 1  & 0.107 & 0.407 & 0.407 \\
    seq-ann        & 1  & 0.090 & 0.426 & 0.426 \\
    seq-stacker    & 3  & 0.087 & 0.418 & 0.411 \\
    seq-popularity & 1  & 0.001 & 0.185 & 0.185 \\
    \bottomrule
  \end{tabular}}
\end{table}

\begin{figure}[ht]
  \centering
  \includegraphics[width=0.82\textwidth]{nexttrack_music_scatter}
  \caption{music@10 versus Recall@10 across all 53 next-track runs, one point per
  run, coloured by predictor family. The two axes correlate ($r = \num{0.609}$: a
  better ranker also lands closer misses) but crown different winners --- the
  Recall@10 champion (a \texttt{seq-blend} run at \num{0.223}) is not the
  musically-closest, which is a \texttt{seq-nexttrack} run (music@10 \num{0.482})
  that regresses the content latent and retrieves by cosine. The
  \texttt{seq-popularity} floor is correctly worst on both. music@10 is
  display-only; Recall@10 stays the primary metric.}
  \label{fig:nexttrack-music-scatter}
\end{figure}

\emph{Interpretation.} music@10 is a genuinely distinct, non-saturated view of
next-track quality: it rewards \emph{sonic} proximity, which the exact-hit metric
ignores. The recall-versus-music decoupling has a clean mechanism --- a model that
regresses the content latent (\texttt{seq-nexttrack}) lands \emph{near} the truth
in content space even when it misses the exact next track, so its misses
``sound right'' more often than the blend's, which optimizes exact rank. This is
why the metric is display-only and must not silently become a selection objective:
promotion still ranks by Recall@10. The metric inherits a mild circularity from the
\emph{balanced} dial (the text block embeds the artist name; the genre block is a
genre prototype), so it partly rewards same-artist / genre proximity --- intended
for ``sounds / feels alike,'' not a perceptual ground truth. A companion
\texttt{GET /api/pathfinder/similar} endpoint serves ``songs like this'' over the
same index.
\emph{Source: \texttt{2026-07-20-musical-distance-metric.md}.}

\section{Session holisticness metrics and the anti-eager levers (MMR $\lambda$ / eager $\beta$)}
\label{sec:holisticness}
\emph{Rationale.} A recurring qualitative complaint about the top blends is that
their predictions feel ``too album-eager / bigram-like'' --- they parrot the
current artist rather than curating a holistic continuation. This campaign
quantifies that symptom with a newly-added \emph{display-only session-holisticness}
metric suite, tests the user's hypothesis that \emph{the GRU beats the LSTM because
it is more eager}, and measures whether two new Phase-2 levers reduce eagerness at
an acceptable recall cost: an eval-time MMR re-rank (\texttt{mmr\_lambda}) and an
anti-eager training regularizer (\texttt{eager\_beta}). It is a measurement /
exploratory pass: nothing is promoted, the primary metric stays Recall@10, and the
holisticness objective decision is explicitly deferred to the user. Every arm runs
on the \emph{canonical} PCA-192 champion split \dsid{seq-20260715-131139}
(\num{1431} test sessions; cold-item rate \num{0.55}), and the champion reproduced
its on-record Recall@10 \num{0.2117} exactly (\num{0.21174}).

\paragraph{The new display-only holisticness metrics.} Wired end-to-end alongside
music@10 (\texttt{seq\_common.eval\_from\_scores} $\to$ Rust $\to$
\texttt{domain.toml} $\to$ UI); the primary metric is unchanged.
\begin{description}
  \item[artist\_adj@10 ($\downarrow$).] Fraction of the top-10 sharing the
  \emph{seed} track's artist --- the ``parrots the current artist'' eagerness rate.
  \textbf{Lower is better}.
  \item[album\_adj@10.] The same on (artist, album); \textbf{omitted on this
  dataset} because its \texttt{items.json} carries no album field (expected, not an
  error --- covered by a separate album-carrying dataset).
  \item[mood\_coh@10 ($\uparrow$).] Mean cosine of the top-10 to the prefix mood
  centroid in the content-metric space; higher is more mood-coherent.
  \item[ild@10 ($\uparrow$).] Intra-list diversity (mean pairwise sonic distance) of
  the top-10; higher is less duplicative.
  \item[suffix\_recall@10 / cont\_prec@10.] Multi-step continuation coverage,
  produced only by the on-demand harness \texttt{seq\_continuation\_eval.py}
  (Phase E below), not by normal runs.
\end{description}
The sanity gate passed: music@10 / mood\_coh@10 / ild@10 populate on all blend and
GRU runs, so the content-metric vectors loaded and the MMR / mood / ild levers are
live, not no-ops. Both new levers are OFF by default and verified byte-identical
when off.

\emph{Results --- the holisticness fingerprint (Phase A).} Table~\ref{tab:holisticness-fingerprint}
fingerprints the model family. \textbf{Eagerness tracks Recall@10 inside the blend
family} --- the champion (Recall@10 \num{0.2117}) is the \emph{most} artist-eager
model measured (artist\_adj@10 \num{0.610}), with R+M next (\num{0.579}) --- but the
eagerness lives in the \emph{Markov} leg, not the RNN: the standalone first-order
Markov bar is already at \num{0.543}, while the single GRU / LSTM sit far lower at
artist\_adj@10 $\approx$ \num{0.345}. So the blends buy their Recall@10 lead by being
artist-adjacent (bigram-like), which is exactly what the holisticness suite flags.
The single RNNs are conversely the \emph{mood-coherent, low-eager} models (highest
mood\_coh@10, lowest real-model artist\_adj@10). Crucially the metrics are \emph{not
trivially gamed}: the popularity floor maxes ild@10 (\num{0.898}) and minimizes
artist\_adj@10 (\num{0.003}) yet scores $\approx 0$ recall and \emph{negative}
mood\_coh@10 --- you maximize them by being random, so the meaningful low-eager /
diverse models are the real single RNNs.

\begin{table}[ht]
  \centering
  \caption{Phase A holisticness fingerprint across the model family, sorted by
  Recall@10 (\num{1431} test sessions). All runs share the prefix
  \protect\runid{run-20260722-165345-}; the id column lists the distinguishing suffix.
  Recall@10 (raw hits of \num{1431}), MRR, artist@10, genre@10, music@10,
  mood\_coh@10, ild@10 higher-better; \textbf{artist\_adj@10 lower-better}. Champion
  row and best cell per metric bolded. \textsuperscript{\dag}~The popularity floor's
  extreme artist\_adj@10 / ild@10 are \emph{degenerate} (it ignores the seed): it
  lands $\approx 0$ recall and negative mood\_coh@10, proving the metrics are not
  gamed low/high by a good model.}
  \label{tab:holisticness-fingerprint}
  {\scriptsize\setlength{\tabcolsep}{3pt}\begin{tabular}{llcccccccc}
    \toprule
    Config & id & R@10 (raw) & MRR & art@10 & gen@10 & mus@10 & artadj@10\,$\downarrow$ & mood@10\,$\uparrow$ & ild@10\,$\uparrow$ \\
    \midrule
    \best{A1 champion \Rp+M+\Cp} & \texttt{82a8f} & \best{0.2117 (303)} & \best{0.122} & 0.333 & 0.401 & 0.454 & 0.610 & 0.398 & 0.406 \\
    A2 R+M (blend)               & \texttt{28dbb} & 0.1712 (245) & 0.095 & 0.329 & 0.398 & 0.443 & 0.579 & 0.399 & 0.403 \\
    A3 GRU alone ($\beta{=}0$)   & \texttt{e39c9} & 0.1230 (176) & 0.053 & \best{0.338} & 0.450 & \best{0.456} & 0.347 & \best{0.479} & 0.454 \\
    B LSTM (else matched to A3)  & \texttt{bb387} & 0.1188 (170) & 0.049 & 0.328 & \best{0.455} & 0.449 & 0.344 & 0.458 & 0.463 \\
    A4 markov (first-order)      & \texttt{047cc} & 0.1069 (153) & 0.069 & 0.298 & 0.389 & 0.407 & 0.543 & 0.310 & 0.462 \\
    A5 popularity                & \texttt{34614} & 0.0007 (1)   & 0.001 & 0.014 & 0.085 & 0.185 & \best{0.003}\textsuperscript{\dag} & $-0.041$ & \best{0.898}\textsuperscript{\dag} \\
    \bottomrule
  \end{tabular}}
\end{table}

\emph{GRU versus LSTM (Phase B) --- the hypothesis is not supported.} With
hyperparameters matched (h256, InfoNCE, 40 epochs, lr \num{1e-3}), the GRU and LSTM
are statistically \emph{tied}: Recall@10 \num{0.1230} versus \num{0.1188} ($\Delta$
$+\num{0.0042}$, inside the $\pm\num{0.015}$ band), and artist\_adj@10 \num{0.347}
versus \num{0.344} (a negligible $+\num{0.003}$). The GRU is not measurably more
eager than the LSTM on this split, so the ``GRU wins because it parrots the current
artist harder'' hypothesis is \textbf{refuted here}.

\emph{The MMR $\lambda$ lever (Phase C) --- the strong knob.} MMR is a pure
eval-time re-rank, so the four originally-launched per-$\lambda$ training runs were
\emph{cancelled} (not failed experiments); the sweep was computed offline from the
single champion checkpoint (the $\lambda{=}$off cell reproduced the champion's
Recall@10 \num{0.2117} exactly). As $\lambda$ drops from off toward \num{0.3}
(Table~\ref{tab:mmr-sweep}; Figure~\ref{fig:nexttrack-mmr-tradeoff}), Recall@10 falls
--- the \emph{only} quantity traded --- while ild@10 rises and artist\_adj@10 falls
(de-eagering). \textbf{$\lambda{=}\num{0.9}$ is essentially free}: the recall
$\Delta$ is $-\num{0.0014}$, well inside the $\pm\num{0.015}$ band (a tie), yet it
already lifts ild@10 \num{0.406}$\to$\num{0.436} and trims eagerness
\num{0.610}$\to$\num{0.603}. \textbf{$\lambda\approx\num{0.7}$} trades $\sim$2 recall
points ($\Delta$ $-\num{0.0181}$, just outside the band) for a large ild@10 gain
(\num{0.406}$\to$\num{0.520}) and a real eagerness drop
(\num{0.610}$\to$\num{0.571}). A notable side effect: \textbf{music@10 rises
monotonically} as you diversify (\num{0.454}$\to$\num{0.492}) --- sonic closeness and
list diversity are not in tension; only exact recall is traded.

\begin{table}[ht]
  \centering
  \caption{Phase C MMR $\lambda$ sweep on the champion, computed offline from its
  checkpoint (\texttt{mmr\_pool} \num{200}; \num{1431} test sessions). Recall@10 (raw
  hits), MRR, artist@10, music@10, mood\_coh@10, ild@10 higher-better;
  \textbf{artist\_adj@10 lower-better}. Best cell per metric bolded; the champion
  default ($\lambda{=}$off) recall / MRR bolded.}
  \label{tab:mmr-sweep}
  {\footnotesize\setlength{\tabcolsep}{4pt}\begin{tabular}{lcccccccc}
    \toprule
    $\lambda$ & R@10 (raw) & $\Delta$ vs off & MRR & art@10 & mus@10 & artadj\,$\downarrow$ & mood\,$\uparrow$ & ild\,$\uparrow$ \\
    \midrule
    off (1.0) & \best{0.2117 (303)} & --- & \best{0.122} & 0.333 & 0.454 & 0.610 & \best{0.398} & 0.406 \\
    0.9 & 0.2103 (301) & $-0.0014$ (TIE) & 0.121 & 0.341 & 0.462 & 0.603 & 0.391 & 0.436 \\
    0.7 & 0.1936 (277) & $-0.0181$ & 0.116 & 0.347 & 0.476 & 0.571 & 0.364 & 0.520 \\
    0.5 & 0.1642 (235) & $-0.0475$ & 0.108 & \best{0.359} & 0.487 & 0.486 & 0.315 & 0.653 \\
    0.3 & 0.1097 (157) & $-0.1020$ & 0.100 & 0.358 & \best{0.492} & \best{0.293} & 0.238 & \best{0.815} \\
    \bottomrule
  \end{tabular}}
\end{table}

\begin{figure}[ht]
  \centering
  \includegraphics[width=0.8\textwidth]{nexttrack_mmr_tradeoff}
  \caption{The MMR $\lambda$ trade-off on the next-track champion. As $\lambda$ drops
  from off toward \num{0.3}, Recall@10 (red, left axis) falls --- the only quantity
  traded --- while ild@10 rises and artist\_adj@10 falls (right axis; more diverse,
  less artist-eager) and music@10 even \emph{rises}. $\lambda{=}\num{0.9}$ sits inside
  the $\pm\num{0.015}$ recall tie band yet already lifts ild@10 by \num{0.030};
  $\lambda\approx\num{0.7}$ buys a large ild@10 gain and real de-eagering for
  $\sim$2 recall points. MMR is the strong holisticness lever; $\beta$
  (Table~\ref{tab:eager-beta}) is weak.}
  \label{fig:nexttrack-mmr-tradeoff}
\end{figure}

\emph{The anti-eager regularizer $\beta$ (Phase D) --- the weak knob.} A
training-time penalty on near-copy prediction of the current track's latent
(\texttt{eager\_margin} \num{0.0}) barely moves the needle
(Table~\ref{tab:eager-beta}): across $\beta$ \numrange{0}{0.2}, artist\_adj@10 shifts
only \num{0.347}$\to$\num{0.340} ($\Delta$ $-\num{0.0074}$) and ild@10 only
\num{0.454}$\to$\num{0.458}, with Recall@10 flat (all four inside $\pm\num{0.015}$).
The single GRU is already low-eager (\num{0.347}), so there is little copy behaviour
left to penalize --- $\beta$ is not a useful holisticness knob at these settings.

\begin{table}[ht]
  \centering
  \caption{Phase D anti-eager regularizer $\beta$ sweep on the single GRU
  (\texttt{eager\_margin} \num{0.0}; \num{1431} test sessions). Recall@10, artist@10,
  mood\_coh@10, ild@10 higher-better; \textbf{artist\_adj@10 lower-better}. Best cell
  per metric bolded.}
  \label{tab:eager-beta}
  {\footnotesize\setlength{\tabcolsep}{4pt}\begin{tabular}{lccccccc}
    \toprule
    run id & $\beta$ & R@10 (raw) & MRR & art@10 & artadj\,$\downarrow$ & mood\,$\uparrow$ & ild\,$\uparrow$ \\
    \midrule
    \runid{run-20260722-165345-e39c9} & 0.0  & 0.1230 (176) & 0.053 & \best{0.338} & 0.347 & \best{0.479} & 0.454 \\
    \runid{run-20260722-165345-353b5} & 0.05 & \best{0.1237 (177)} & 0.053 & 0.337 & 0.346 & 0.478 & 0.455 \\
    \runid{run-20260723-002044-af11a} & 0.1  & 0.1202 (172) & 0.053 & 0.335 & 0.344 & 0.477 & 0.456 \\
    \runid{run-20260723-002044-17cf8} & 0.2  & 0.1202 (172) & 0.051 & 0.336 & \best{0.340} & 0.475 & \best{0.458} \\
    \bottomrule
  \end{tabular}}
\end{table}

\emph{Multi-step continuation (Phase E).} The on-demand harness
(\texttt{seq\_continuation\_eval.py} \texttt{-m 3 -k 10}, \num{1172} continuation
cases) shows the GRU \emph{beats} the Markov bigram on both multi-step metrics
(Table~\ref{tab:continuation}): the RNN's advantage \emph{widens} past the immediate
next track, because the bigram's single-step co-occurrence does not compound into
good 3-step continuations whereas the GRU's learned sequence state does --- the
mirror image of the single-step board, where the bigram-carrying blend wins raw
Recall@10.

\begin{table}[ht]
  \centering
  \caption{Phase E multi-step continuation (harness \texttt{-m 3 -k 10},
  \num{1172} cases) on \dsid{seq-20260715-131139}: the GRU (A3 checkpoint) beats the
  first-order Markov bigram on both metrics (higher-better). Winner bolded.}
  \label{tab:continuation}
  {\footnotesize\setlength{\tabcolsep}{6pt}\begin{tabular}{lcc}
    \toprule
    Model & suffix\_recall@10 & cont\_prec@10 \\
    \midrule
    markov (first-order bigram) & 0.0869 & 0.0236 \\
    \best{GRU (gru-infonce-h256, A3)} & \best{0.0963} & \best{0.0273} \\
    \bottomrule
  \end{tabular}}
\end{table}

\emph{Interpretation and the deferred decision.} The campaign delivers a clean
mechanistic account of the eagerness symptom: it is sourced from the Markov bigram
leg (which parrots the current artist via raw co-occurrence), not the neural leg, so
any future de-eagering should target the bigram / the fusion, not the RNN. Of the two
new levers, the eval-time MMR re-rank is the \emph{strong} knob and the training-time
$\beta$ regularizer is \emph{weak}. Nothing is promoted --- the primary metric stays
Recall@10 and the champion is unchanged --- but \textbf{$\lambda\approx\num{0.9}$ is
flagged as a near-free holisticness upgrade for a future user promotion decision}
(Recall@10 tie within noise, ild@10 \num{0.406}$\to$\num{0.436}, eagerness
\num{0.610}$\to$\num{0.603}), with $\lambda\approx\num{0.7}$ the ``meaningful
de-eagering for $\sim$2 recall points'' operating point; adopting either requires a
user decision plus a 3-seed / native-run confirm. The $\beta$ regularizer is not
recommended as a lever at these settings.
\emph{Source: \texttt{2026-07-22-session-holisticness-and-antieager-levers.md}.}

\section{Dual-tower fusion refuted: the fourth combiner null}
\label{sec:dualtower}
\emph{Rationale.} By 2026-07-25 the record held three independent
\emph{learned-combiner} nulls --- the XGB stacker (\numrange{0.080}{0.087},
\emph{below} its own GRU base leg; Section~\ref{sec:phase2}), reciprocal-rank
fusion ($-\num{0.029}$) and even \emph{oracle} test-fit blend weights
($+\num{0.004}$, CI straddling zero; Section~\ref{sec:projection}). All three
were \textbf{post-hoc, score-level} combiners, which left one obvious rebuttal
standing: \emph{of course} a bolted-on combiner fails --- train the fusion
end-to-end so the parts co-adapt. This campaign runs exactly that rebuttal, at
the \emph{representation} level, jointly, from scratch.

The vehicle is a new predictor, \texttt{seq-dualgru}: two \emph{independent}
recurrent towers (separate weights) read the prefix in parallel over
configurable \emph{causal views}, their per-step hidden states are concatenated,
and one MLP head fuses them into the predicted next-item latent. Three views are
exposed, all strictly causal at step $t$: \texttt{latent} (the raw item vector),
\texttt{delta} (the step-to-step difference, $x_{-1}=0$), and \texttt{cummean}
(the running session centroid --- the quantity the \texttt{seq-ann} pooled-MLP
leg was built from). Crucially it reuses \texttt{seq\_nexttrack.make\_batches}
and \texttt{step\_loss} verbatim, so objective, batching, early stopping and
evaluation are byte-for-byte the family's, and it carries no bidirectional mode
(a backward pass would leak the teacher-forcing target).

The hypothesis was sharp: two towers should beat one \emph{only if they see
genuinely different views}; two towers on the same \texttt{latent} view should
reduce to a capacity / implicit-ensemble effect a wider single GRU can buy more
cheaply. The design therefore crosses the \emph{complete upper triangle} of the
$\{\texttt{latent},\texttt{delta},\texttt{cummean}\}^2$ view matrix (six cells at
\texttt{fusion\_layers}${}=1$) with a \emph{fusion-depth} axis (two cells at
\texttt{fusion\_layers}${}=0$, a pure linear read of the concatenated states) and
--- the axis the earlier topology scan lacked (Section~\ref{sec:topology}) --- a
\textbf{parameter-matched single-GRU capacity control}. Parameter counts were
re-measured by instantiating the real classes at \texttt{latent\_dim}${}=192$:
the control C1 (h425, \num{871017} parameters) matches the six
\texttt{fusion\_layers}${}=1$ dual arms (\num{871872}) to
\SI{0.098}{\percent}, and is \SI{10}{\percent} \emph{larger} than the
\texttt{fusion\_layers}${}=0$ arms (\num{789696}) --- deliberately generous to
the single tower.

\paragraph{The honest bar.} Two boards were pre-registered \emph{separately}, so
that the scientific verdict and the board event could not be traded off after the
fact. On \textbf{Recall@10} the bar for a standalone single model is the
single-model board \num{0.123}, \emph{not} the three-leg blend champion
\num{0.212} --- comparing one trained model against a fusion of three legs would
be a category error, and dual-tower-as-the-\Rp-leg is a \texttt{seq-blend} code
change explicitly out of scope. On \textbf{holisticness@10} the crown rule
carries an anti-degenerate recall floor (see below). All ten runs are on the
canonical split \dsid{seq-20260715-131139} at seed 1337, executed one at a time
on a single remote worker, and the baseline C0 reproduced its on-record
Recall@10 \num{0.12299} (\num{176} of \num{1431}) \emph{exactly}.

\emph{Results.} Table~\ref{tab:dualtower} is the batch. \textbf{No arm is bolded
as a winner because there is no winner}: the single-GRU baseline holds the
batch's best Recall@10, and not one of the eight dual arms beats it by
paired-$\Delta$ CI above zero --- six lose outright and two (the
\texttt{fusion\_layers}${}=0$ pair) tie. The pre-registered refutation rule
fired.

\begin{table}[ht]
  \centering
  \caption{Dual-tower view-pairing $\times$ fusion-depth scan, ten runs on
  \dsid{seq-20260715-131139} (\num{1431} test sessions), sorted by Recall@10.
  All arms share the run-id prefix \protect\runid{run-20260725-143502-}; C0 is
  \protect\runid{run-20260725-143143-97c1e}. Paired $\Delta$ is a per-session
  paired bootstrap against C0 (\num{2000} resamples, rng 1337). Recall@10 raw =
  successes of \num{1431}; CI bounds are given to three decimals, and a CI
  straddling zero is a \emph{tie}. Arm codes carry the view pair
  (lat\,$=$\,\texttt{latent}, del\,$=$\,\texttt{delta},
  cum\,$=$\,\texttt{cummean}) and the fusion depth (\texttt{-f0} $=$ linear
  read, otherwise a 256-wide MLP head). R@10, MRR, mus@10, ild@10
  higher-better; \textbf{artadj@10 lower-better}; H is the \emph{3-factor}
  holisticness generation (Section~\ref{sec:power}). Best cell per metric
  bolded.
  \textsuperscript{\dag}~D-CC's extreme artadj@10 / H are \emph{not} a real
  crown: they are bought at \num{95} of \num{1431} hits, \num{0.0566} below the
  baseline, failing the $\pm\num{0.015}$ recall floor by $3.8\times$.}
  \label{tab:dualtower}
  {\scriptsize\setlength{\tabcolsep}{1.7pt}\begin{tabular}{llrrlccccc}
    \toprule
    Arm & id & params & R@10 (raw) & $\Delta$ vs C0 [95\% CI] & MRR & mus@10 & artadj\,$\downarrow$ & ild\,$\uparrow$ & H\,$\uparrow$ \\
    \midrule
    C0 h256 \emph{base} & \texttt{97c1e} & \num{394944} & \best{0.12299 (176)} & --- & 0.05341 & 0.4558 & 0.3474 & 0.4536 & 0.14172 \\
    C1 h425 \emph{ctrl} & \texttt{be11b} & \num{871017} & 0.12089 (173) & $-0.0021$ [$-.011$, $+.006$] & 0.05215 & \best{0.4597} & 0.3640 & 0.4506 & 0.13850 \\
    D-LL-f0 \emph{lat/lat} & \texttt{18b4a} & \num{789696} & 0.11880 (170) & $-0.0042$ [$-.013$, $+.004$] & \best{0.05631} & 0.4553 & 0.3630 & 0.4506 & 0.13774 \\
    D-LC-f0 \emph{lat/cum} & \texttt{3e5d4} & \num{789696} & 0.11740 (168) & $-0.0056$ [$-.017$, $+.005$] & 0.05087 & 0.4458 & 0.3557 & 0.4630 & 0.13648 \\
    D-LL \emph{lat/lat} & \texttt{12fe6} & \num{871872} & 0.10273 (147) & $-0.0203$ [$-.035$, $-.007$] & 0.04235 & 0.4446 & 0.3092 & 0.4873 & 0.14635 \\
    D-LC \emph{lat/cum} & \texttt{a95c5} & \num{871872} & 0.09713 (139) & $-0.0259$ [$-.039$, $-.013$] & 0.03901 & 0.4380 & 0.2997 & 0.4913 & 0.15000 \\
    D-LD \emph{lat/del} & \texttt{8d214} & \num{871872} & 0.09154 (131) & $-0.0315$ [$-.045$, $-.018$] & 0.04029 & 0.4385 & 0.2999 & 0.4988 & 0.14673 \\
    D-DC \emph{del/cum} & \texttt{e9cfb} & \num{871872} & 0.08665 (124) & $-0.0363$ [$-.050$, $-.022$] & 0.03379 & 0.4125 & 0.2319 & 0.5413 & 0.16340 \\
    D-CC \emph{cum/cum} & \texttt{88a4c} & \num{871872} & 0.06639 (95) & $-0.0566$ [$-.072$, $-.041$] & 0.02664 & 0.3935 & \best{0.1899}\textsuperscript{\dag} & 0.5235 & \best{0.19715}\textsuperscript{\dag} \\
    D-DD \emph{del/del} & \texttt{4324f} & \num{871872} & 0.05800 (83) & $-0.0650$ [$-.082$, $-.049$] & 0.02819 & 0.4016 & 0.2363 & \best{0.5457} & 0.13335 \\
    \bottomrule
  \end{tabular}}
\end{table}

\emph{Capacity is not the missed lever either.} The param-matched control C1
(h425, \num{871017} parameters --- $2.2\times$ C0) scores \num{0.12089} against
C0's \num{0.12299}, a paired $\Delta$ of $-\num{0.0021}$ straddling zero. This
fills a real gap: single-GRU \emph{width} had only ever been scanned on the
64-dim AE space (Section~\ref{sec:topology}), and the rich PCA-192 space was
untested. It behaves the same way, so \textbf{single-GRU width is now closed on
PCA-192 as well as AE-64}. Because C1 is a genuine \SI{0.098}{\percent}
parameter match to the \texttt{fusion\_layers}${}=1$ arms, the null cannot be
re-attributed to ``capacity would have helped'': it demonstrably would not.

\emph{View attribution --- complementarity is real, but lands on the wrong
pair.} The complete view triangle is what makes this mechanically decidable
rather than a judgment call, and the answer is genuinely surprising
(Table~\ref{tab:dualtower-views}, Figure~\ref{fig:nexttrack-dualtower-matrix}).
The pair the design nominated as its best hope --- \texttt{latent} $+$
\texttt{cummean} --- is \emph{not} complementary: it ties its latent/latent
parent and beats only the centroid-only one, i.e.\ the raw-latent tower carries
the cell and the centroid tower is dead weight. Same for \texttt{latent} $+$
\texttt{delta}. But \texttt{delta} $+$ \texttt{cummean} \emph{is} complementary
--- it beats \emph{both} parents with CI above zero. Read together: the
co-adaptation mechanism the hypothesis proposed genuinely \textbf{works}; it is
simply that \textbf{the raw \texttt{latent} view already subsumes everything the
other two carry}. A GRU over raw latents integrates step-to-step movement and
the running mean internally, so pairing it with an explicit computation of
either is redundant --- while two \emph{impoverished} views do have
non-overlapping information to pool, at a hopeless absolute level
(\num{0.08665}). That is a clean mechanistic explanation for the null, not
merely an absence of effect.

\begin{table}[ht]
  \centering
  \caption{View attribution: each mixed cell against \emph{both} of its
  single-view parents (paired $\Delta$ Recall@10, \num{1431} sessions). Only one
  of three mixed pairs is genuinely complementary --- and it is the one with no
  raw-latent tower.}
  \label{tab:dualtower-views}
  {\scriptsize\setlength{\tabcolsep}{3.4pt}\begin{tabular}{llll}
    \toprule
    Mixed cell & vs parent 1 & vs parent 2 & Verdict \\
    \midrule
    D-LC \emph{lat/cum} & D-LL: $-0.0056$ [$-.018$, $+.007$] & D-CC: $+0.0308$ [$+.018$, $+.045$] & \emph{not} complementary \\
    D-LD \emph{lat/del} & D-LL: $-0.0112$ [$-.023$, $+.001$] & D-DD: $+0.0335$ [$+.020$, $+.047$] & \emph{not} complementary \\
    \best{D-DC \emph{del/cum}} & D-DD: $+0.0287$ [$+.015$, $+.043$] & D-CC: $+0.0203$ [$+.008$, $+.034$] & \best{complementary: both CI$>$0} \\
    \bottomrule
  \end{tabular}}
\end{table}

\begin{figure}[ht]
  \centering
  \includegraphics[width=0.8\textwidth]{nexttrack_dualtower_matrix}
  \caption{The complete view-pairing matrix at \texttt{fusion\_layers}${}=1$,
  one run per cell. Every dual cell sits \emph{below} the single-GRU baseline
  (\num{0.12299}), and the ordering is monotone in how much raw \texttt{latent}
  the pair contains: both-latent is best, both-delta worst. The lower triangle
  was never run (the pairing is symmetric by construction) and is masked in grey
  rather than mirrored, so the colour scale reflects only measured cells. The
  gradient is the finding: a recurrent tower over raw latents already computes
  what \texttt{delta} and \texttt{cummean} make explicit.}
  \label{fig:nexttrack-dualtower-matrix}
\end{figure}

\emph{The MLP fusion head is actively harmful, not merely useless.} Both
\texttt{fusion\_layers}${}=0$ arms beat their \texttt{fusion\_layers}${}=1$ twins
by paired-$\Delta$ CI above zero ($+\num{0.0161}$ on latent/latent,
$+\num{0.0203}$ on latent/cummean) at \SI{10}{\percent} \emph{fewer} parameters
--- so removing the head makes the model measurably better in \emph{both} view
regimes. Mechanistically this fits the whole family arc: the objective is
cosine/InfoNCE retrieval \emph{in} the item-latent space, so the head's job is
to land \emph{in} that geometry, and an extra 256-wide non-linear layer between
the recurrent state and a metric read-out just adds a warp the loss must then
undo. The trade-off runs the other way on holisticness (both f0 arms have
\emph{lower} H, $\Delta$H CI below zero): the f1 arms' extra warp diffuses the
top-ten (higher ild, lower artist\_adj) at the cost of exact hits. The rule
adopted is \textbf{default \texttt{fusion\_layers}${}=0$}, later confirmed a
third time (Table~\ref{tab:fusion-depth}).

\emph{Holisticness and recall are anti-correlated, and the recall floor earned
its place.} Five arms lift holisticness@10 by paired-$\Delta$H CI above zero,
topping out at D-CC's \num{0.19715} ($+\num{0.055}$) --- which on the naked crown
metric would have ranked second on the entire board. But D-CC lands \num{95} of
\num{1431} exact hits against the baseline's \num{176}. The ordering of the batch
by H is almost the \emph{reverse} of its ordering by recall
(Figure~\ref{fig:nexttrack-holisticness-antiwin}), so \textbf{all five
$\Delta$H winners fail the pre-registered $\pm\num{0.015}$ recall floor} and the
crown rule never fired. Without that floor the campaign would have ``won the
crown'' with a model \SI{46}{\percent} worse at its actual job --- precisely the
degenerate \texttt{mood-session} failure mode (H \num{0.2332} at Recall@10
\num{0.0238}) the floor was written to block. On the present evidence,
\textbf{unconstrained holisticness is maximized by predicting less well}, so any
future holisticness objective on this corpus must be a \emph{constrained} one.

\begin{figure}[ht]
  \centering
  \includegraphics[width=0.8\textwidth]{nexttrack_holisticness_antiwin}
  \caption{Why the crown rule needs a hard recall floor. Across the ten arms,
  holisticness@10 runs almost exactly \emph{backwards} along Recall@10: the five
  arms that clear the $\Delta$H CI-above-zero leg (green) are the five furthest
  \emph{below} the baseline on exact recall, and all five fall outside the
  $\pm\num{0.015}$ floor (shaded). The open square is the incumbent
  \texttt{mood-session} model that motivated the floor --- the highest
  holisticness on record at \num{34} of \num{1431} correct next tracks. H values
  are the 3-factor generation (Section~\ref{sec:power}).}
  \label{fig:nexttrack-holisticness-antiwin}
\end{figure}

\emph{Interpretation.} Nothing was registered and nothing promoted; the champion,
the leaderboard and the pins are untouched. What the campaign buys is a
\emph{closure}: the combiner null is now robust across the whole combiner design
space on this corpus --- post-hoc score-level (three ways) \emph{and}
representation-level jointly-trained with a parameter-matched control --- which
materially consolidates ``the fixed equal-thirds z-blend is the combiner of
record.'' Two genuine sub-findings survive the null (the harmful MLP fusion head,
and \texttt{delta}$+$\texttt{cummean} complementarity), and the machinery itself
is sound: 10 of 10 runs succeeded, the new predictor trains, converges and
evaluates at exact parity with the family, and D-LL-f0 even posts the batch's
best MRR (\num{0.05631} against C0's \num{0.05341}) --- the one metric on which a
dual arm leads. Per the noise-band discipline that is ``at least equal, likely
neither better nor worse,'' not a win.

\paragraph{A provenance finding that revises two earlier numbers.} Before the
batch, the live remote worker directory was found \emph{stale}: its
\texttt{seq\_common.py} had no \texttt{holisticness\_at\_k} at all,
\texttt{seq\_nexttrack.py} differed from the hub by md5, and three
\texttt{seq\_*.py} modules were missing entirely. After refreshing and verifying
the whole family md5-identical to the hub, C0 returned \num{0.12299},
bit-identical to the hub's on-record value --- \emph{not} the \num{0.1201957}
that two earlier remote runs (\runid{run-20260723-002044-af11a},
\texttt{-17cf8}) had returned. The ``cross-host reduction-order offset of
$-\num{0.0028}$'' those runs had suggested is therefore \textbf{withdrawn: it was
a code difference, not a host effect}. Their recall numbers belong to a
superseded code state and are not code-comparable to later batches; they remain
valid as the $\beta$-eagerness measurements they were taken for
(Table~\ref{tab:eager-beta}), since those conclusions rest on within-batch
comparisons. The general lesson --- a divergent remote checkout produces a
\emph{wrong number}, not a crash --- is now a standing pre-flight gate.
\emph{Source: \texttt{2026-07-25-nexttrack-dual-tower-fusion-scan.md}.}

\section{The per-step pre-encoder refuted, and the rectifier mechanism}
\label{sec:preencoder}
\emph{Rationale.} The dual-tower null closed \emph{deterministic} view
asymmetry: \texttt{delta} and \texttt{cummean} are lossy functions of what a
raw-latent tower already holds, so pairing them buys nothing. This campaign asks
the strictly stronger version --- does a \textbf{learned} per-step
re-embedding, which is \emph{not} a function of information the tower already
has, enlarge the hypothesis space usefully? The axis is a nonlinear
re-embedding in front of the recurrence,
$[\,\text{Linear}(d_{\text{in}},w),\ \text{ReLU},\ \text{Dropout}\,]\times
\texttt{pre\_layers}$, applied per timestep (strictly causal), with the GRU's
\texttt{input\_size} becoming $w$.

The hypothesis has an unusually sharp analytic edge. A GRU's own input transform
is \emph{already linear} ($W_i x$ per gate), and a \emph{linear} pre-encoder at
$w \ge d_{\text{in}}$ is exactly as expressive as a bare GRU, since
$W_i(Vx) = (W_iV)x$ with no rank constraint. \textbf{The ReLU is therefore the
only expressivity change}, which means no linear-pre control run is needed and a
null cannot be explained away as ``the pre-encoder was redundant by
construction.'' The design's decisive move is the \textbf{single-tower ablation}
S1 --- a \texttt{MLP256}$\to$\texttt{GRU256} against a parameter-matched plain
GRU --- because without it any dual-tower result is uninterpretable: pre-encoder,
or second tower?

Ten runs, same canonical split, same single remote worker, zero failures. Two
reproduction gates both passed \emph{bit-identically}: C0 returned Recall@10
\num{0.12299091544374564} (\num{176} of \num{1431}, identical to the last digit,
a seventh reproduction), confirming the \texttt{pre\_hidden} edit did not perturb
the family's default code path; and C3, the bare dual-tower arm, returned
\num{0.10272536687631029}, exactly its on-record value, validating every dual
verdict.

\emph{Results.} Table~\ref{tab:preencoder} is the batch and
Figure~\ref{fig:nexttrack-preencoder-forest} is the verdict at a glance. Not one
of the six pre-MLP arms beats the baseline by paired-$\Delta$ CI above zero:
five lose outright, one ties. \textbf{And the decisive ablation does not merely
tie --- it loses}: S1 at \num{0.10901} against C0's \num{0.12299}, $\Delta$
$-\num{0.0140}$ with the CI entirely below zero, \num{156} hits against
\num{176}. Since a linear pre-encoder here is provably expressivity-neutral,
\textbf{the ReLU re-embedding is destroying prefix information that the GRU's own
gated linear input transform preserved}.

\begin{table}[ht]
  \centering
  \caption{Pre-encoder presence $\times$ placement $\times$ dose, ten runs on
  \dsid{seq-20260715-131139} (\num{1431} test sessions), sorted by Recall@10.
  All arms share the prefix \protect\runid{run-20260726-161801-}; C0 is
  \protect\runid{run-20260726-161429-bdcbc}. Paired $\Delta$ is a per-session
  paired bootstrap against C0 (\num{2000} resamples, rng 1337); CI bounds are
  given to three decimals, and a CI straddling zero is a \emph{tie}. Dual-arm
  codes give the two towers' pre-encoder widths (A0/B256 $=$ tower A bare, tower
  B rectified at width 256) and the fusion depth. H here is the
  \emph{4-factor} (grounded) holisticness generation and is \textbf{not}
  comparable to Table~\ref{tab:dualtower}'s 3-factor column
  (Section~\ref{sec:power}). Best cell per metric bolded; \textbf{artadj@10
  lower-better}. \textsuperscript{\dag}~T4's minimal artadj and T1's maximal ild
  are the familiar de-eagering-by-predicting-worse trade: both sit
  $\sim$\num{0.025} below C0 on recall.}
  \label{tab:preencoder}
  {\scriptsize\setlength{\tabcolsep}{1.7pt}\begin{tabular}{llrrlccccc}
    \toprule
    Arm & id & params & R@10 (raw) & $\Delta$ vs C0 [95\% CI] & MRR & mus@10 & artadj\,$\downarrow$ & ild\,$\uparrow$ & H\,$\uparrow$ \\
    \midrule
    C2 h454 \emph{ctrl} & \texttt{7646d} & \num{969936} & \best{0.12509 (179)} & $+0.0021$ [$-.007$, $+.011$] & 0.05389 & 0.4572 & 0.3668 & 0.4466 & 0.045765 \\
    T1-f0 \emph{A0/B256, f0} & \texttt{a6423} & \num{888256} & 0.12369 (177) & $+0.0007$ [$-.012$, $+.013$] & \best{0.05470} & \best{0.4575} & 0.3627 & 0.4616 & 0.044541 \\
    C0 h256 \emph{base} & \texttt{bdcbc} & \num{394944} & 0.12299 (176) & --- & 0.05341 & 0.4558 & 0.3474 & 0.4536 & 0.047086 \\
    C1 h297 \emph{ctrl} & \texttt{7bab3} & \num{494697} & 0.11321 (162) & $-0.0098$ [$-.017$, $-.003$] & 0.05199 & 0.4513 & 0.3602 & 0.4474 & 0.045285 \\
    \best{S1 \emph{preMLP256}} & \texttt{ad53f} & \num{493504} & 0.10901 (156) & $\mathbf{-0.0140}$ [$-.027$, $-.001$] & 0.04463 & 0.4499 & 0.3240 & 0.4815 & \best{0.047132} \\
    C3 \emph{bare dual, f1} & \texttt{0df6f} & \num{871872} & 0.10273 (147) & $-0.0203$ [$-.035$, $-.007$] & 0.04235 & 0.4446 & 0.3092 & 0.4873 & 0.046612 \\
    T2 \emph{A0/B128, f1} & \texttt{dc7d2} & \num{847424} & 0.09993 (143) & $-0.0231$ [$-.037$, $-.009$] & 0.04490 & 0.4415 & 0.3039 & 0.4939 & 0.046431 \\
    T3 \emph{A0/B384, f1} & \texttt{ca0b8} & \num{1093440} & 0.09993 (143) & $-0.0231$ [$-.037$, $-.009$] & 0.04333 & 0.4441 & 0.3068 & 0.4961 & 0.045895 \\
    T4 \emph{A256/B256, f1} & \texttt{bd5b5} & \num{1068992} & 0.09783 (140) & $-0.0252$ [$-.039$, $-.011$] & 0.03829 & 0.4343 & \best{0.2845}\textsuperscript{\dag} & 0.4986 & 0.043985 \\
    T1 \emph{A0/B256, f1} & \texttt{9450b} & \num{970432} & 0.09713 (139) & $-0.0259$ [$-.040$, $-.013$] & 0.04176 & 0.4369 & 0.2962 & \best{0.5016}\textsuperscript{\dag} & 0.045769 \\
    \bottomrule
  \end{tabular}}
\end{table}

\begin{figure}[ht]
  \centering
  \includegraphics[width=0.8\textwidth]{nexttrack_preencoder_forest}
  \caption{The pre-encoder verdict. Paired per-session bootstrap $\Delta$
  Recall@10 against the reproduced baseline, one row per arm; the shaded strip is
  the $\pm\num{0.015}$ practical tie band. Not one pre-encoder arm reaches a CI
  entirely above zero. The decisive row is S1 --- the single tower, the only arm
  whose result is unambiguously attributable to the pre-encoder --- which loses
  outright. Because a \emph{linear} pre-encoder at this width is provably
  expressivity-neutral, the rectifier is the only candidate cause.}
  \label{fig:nexttrack-preencoder-forest}
\end{figure}

\emph{The dose curve is flat, and placement is irrelevant.} Sweeping
\texttt{pre\_hidden\_b} through $\{0,128,256,384\}$ at
\texttt{fusion\_layers}${}=1$ (Table~\ref{tab:preencoder-dose}), every adjacent
pair straddles zero and so do the endpoints --- a \textbf{flat} axis, which by
the pre-registered clause would have blocked registration even had a win rule
fired. But note the \emph{nominal} ordering is downward-then-flat, never upward:
every dosed arm sits below the bare dual. Taken with S1's significant loss, the
sign of this axis is consistently negative across both towers and all four
widths, so there is no interior peak to chase. Symmetry is likewise a
non-factor: T1 (asymmetric, one bare tower) against T4 (both towers rectified)
straddles zero on recall, so keeping tower A bare is not load-bearing at this
dose.

\begin{table}[ht]
  \centering
  \caption{The dose curve, as measured: paired $\Delta$ Recall@10 against the
  bare dual-tower arm C3 at \texttt{fusion\_layers}${}=1$. Flat by the
  pre-registered test --- but nominally downward, never upward.}
  \label{tab:preencoder-dose}
  {\footnotesize\setlength{\tabcolsep}{6pt}\begin{tabular}{rlll}
    \toprule
    \texttt{pre\_hidden\_b} & Arm & R@10 (raw) & $\Delta$ vs C3 (bare) \\
    \midrule
    0 & C3 & 0.10273 (147) & --- \\
    128 & T2 & 0.09993 (143) & $-0.0028$ [$-0.0147$, $+0.0091$] straddles \\
    256 & T1 & 0.09713 (139) & $-0.0056$ [$-0.0189$, $+0.0077$] straddles \\
    384 & T3 & 0.09993 (143) & $-0.0028$ [$-0.0154$, $+0.0105$] straddles \\
    \bottomrule
  \end{tabular}}
\end{table}

\emph{The one solid positive: the fusion-depth regression, measured a third
time and bigger.} T1-f0 beats T1 by $+\num{0.0266}$ with the CI above zero,
against the $+\num{0.016}$ and $+\num{0.020}$ measured in the two dual-tower view
regimes (Table~\ref{tab:fusion-depth}). The mechanism holds and sharpens: every
extra nonlinear layer between the recurrent state and the metric read-out adds a
warp the InfoNCE loss must undo, and \emph{the more nonlinearity already added on
the input side, the more the output-side warp costs}. Note too that T1-f0 --- the
only pre arm that ties the baseline --- is exactly the arm with \emph{both} a bare
tower A \emph{and} \texttt{fusion\_layers}${}=0$: the shortest un-warped path from
raw latents to the read-out. \texttt{fusion\_layers}${}=0$ is the family default,
full stop.

\begin{table}[ht]
  \centering
  \caption{The \texttt{fusion\_layers} regression, three independent
  measurements. Paired $\Delta$ Recall@10 of the linear-read arm minus its
  MLP-head twin; positive means removing the fusion MLP \emph{helps}. The
  $\Delta$H column is available only for the two same-batch 3-factor pairs.
  \emph{Sources: \texttt{2026-07-25-nexttrack-dual-tower-fusion-scan.md} (rows
  1--2), \texttt{2026-07-26-nexttrack-pre-encoder-scan.md} (row 3).}}
  \label{tab:fusion-depth}
  {\scriptsize\setlength{\tabcolsep}{4pt}\begin{tabular}{llll}
    \toprule
    Pair (f0 $-$ f1) & Regime & $\Delta$ R@10 [95\% CI] & $\Delta$H [95\% CI] \\
    \midrule
    D-LL-f0 $-$ D-LL & redundant views, no pre-encoder & $+0.0161$ [$+0.0021$, $+0.0301$] & $-0.0086$ [$-0.0116$, $-0.0055$] \\
    D-LC-f0 $-$ D-LC & complementary views, no pre-encoder & $+0.0203$ [$+0.0063$, $+0.0335$] & $-0.0135$ [$-0.0171$, $-0.0100$] \\
    \best{T1-f0 $-$ T1} & \best{with a per-step pre-encoder} & \best{$+0.0266$ [$+0.0140$, $+0.0391$]} & --- \\
    \bottomrule
  \end{tabular}}
\end{table}

\emph{The mechanism, measured rather than inferred.} The analytic argument named
the rectifier as the only candidate; a follow-up interpretability read then
\emph{measured} it (Figure~\ref{fig:nexttrack-rectifier}). The decisive control
was a counterfactual tap the sparse-autoencoder engine cannot express --- the
model's \emph{own} affine map with the ReLU deleted. Reading next-item genre
decodability at the same step (\num{11184} probe test rows, 25 classes, paired
standard error $\le\num{0.0065}$): the raw latent scores \num{0.4033}, the
affine map \num{0.4084} (\textbf{neutral, ${\sim}0.8\sigma$}, exactly as
$W_i(Vx)=(W_iV)x$ predicts), and the rectified tap \num{0.3700}
(\textbf{${\sim}5.9\sigma$ --- the ReLU is the whole loss}). The smoking gun is
arithmetic: \texttt{frac\_exact\_zero} on the rectified tap is \num{0.5327},
matching the affine tap's \texttt{frac\_negative\_coords} \num{0.5327} to four
digits. The rectifier deletes \emph{precisely} the negative half, leaving
${\sim}120$ live units for 192 signed directions --- \emph{fewer live coordinates
than the input has dimensions}. PCA-192 latents are zero-centred and
near-symmetric by construction, and at $w=256$ (only $1.33\times$ the input
dimension) there is nowhere near enough over-completeness to re-encode the
discarded half-space, unlike the wide embeddings where ReLU MLPs normally live.
The damage survives the recurrence: within one model, \texttt{tower\_b} (fed
rectified) scores \num{0.3929} against its \emph{identical twin}
\texttt{tower\_a} (fed raw) at \num{0.4183}, ${\sim}3.9\sigma$ --- and
\texttt{tower\_a} reproduces the standalone plain GRU (\num{0.4235}), which is
exactly what ``the bare tower \emph{rescues}'' requires.

\begin{figure}[ht]
  \centering
  \includegraphics[width=0.8\textwidth]{nexttrack_rectifier}
  \caption{The rectifier mechanism, measured. Next-item genre decodability at
  each representation tap, including one counterfactual the interpretability
  engine cannot express: the model's own affine map with the ReLU deleted. On the
  input side the affine map is \emph{free} (${\sim}0.8\sigma$) and the ReLU is
  the \emph{entire} cost (${\sim}5.9\sigma$); after the recurrence the damage
  persists (${\sim}3.9\sigma$ between two identical towers that differ only in
  whether their input was rectified), while the raw-fed tower reproduces the
  standalone plain GRU. \texttt{frac\_exact\_zero} on the rectified tap
  (\num{0.5327}) matches the affine tap's \texttt{frac\_negative\_coords}
  (\num{0.5327}) to four digits: half-wave rectification of a zero-centred
  latent deletes exactly the negative half, leaving ${\sim}120$ live units for
  192 signed directions.}
  \label{fig:nexttrack-rectifier}
\end{figure}

\emph{Interpretation, and one honest caveat about the controls.} This is the
fifth consecutive architecture null, and it closes the \emph{input side} as the
view null closed the input-view direction: together with the four
learned-combiner nulls, adding representation machinery \emph{around} a single
GRU over a good latent space does not help on this corpus, whereas every real
gain on record came from improving the \emph{latent space itself} (the InfoNCE
content projection, PCA-192) or from the fixed z-blend. The axis is now
quantitatively bounded --- best pre arm \num{0.12369}, a tie; dose flat to twice
the input dimension; placement irrelevant; the fusion-depth interaction
resolved. The caveat is that the \emph{parameter-matched width control was a
weaker instrument than assumed}: C1 (h297) came in at \num{0.11321}, itself
losing to C0, and the batch's width points run \num{176} $\to$ \num{162} $\to$
\num{173} $\to$ \num{179} hits at h256 / h297 / h425 / h454 --- non-monotone,
spanning \num{0.0119} in recall with no trend. ``Single-GRU width is flat on
PCA-192'' survives as a \emph{trend} claim (h454 ties h256), but per-width jitter
is ${\sim}\pm\num{0.01}$ on a single split, roughly two-thirds of the noise band.
\textbf{A width-matched control that is itself a single unreplicated run is not a
reliable capacity anchor at this effect size} --- which is why it matters that
S1's verdict rests on the C0 comparison, reproduced to the last digit. A second
caveat is un-disentangled: the pre-MLP ends in \texttt{Dropout}, so pre arms
receive GRU-\emph{input} dropout that the bare arms lack, and S1's loss could be
partly over-regularization rather than the rectifier --- separating them needs a
\texttt{pre\_dropout} parameter that does not yet exist.
\emph{Source: \texttt{2026-07-26-nexttrack-pre-encoder-scan.md}.}

\section{Reading the record: statistical power and the holisticness generations}
\label{sec:power}
Two cross-cutting facts emerged from the 2026-07-25/26 batches that change how
every number above should be read. Neither is a result about a model; both are
results about the \emph{measuring instrument}, and they belong here rather than
buried in one campaign.

\subsection{The nulls are power-limited, and that is a different claim}
\label{sec:power-stats}
Every verdict in this document rests on a paired per-session bootstrap over the
same \num{1431} test sessions. That $n$ has a measurable resolution. Averaged
over the nine paired comparisons of the pre-encoder batch the 95\% CI half-width
is \textbf{\num{0.0124}} on Recall@10 (\num{0.0128} over the dual-tower batch's
nine, so the figure is stable), and \textbf{${\sim}\num{0.0018}$} on
holisticness@10. Now compare those to the \emph{entire spread} of an architecture
family: \num{0.028} on Recall@10 and \num{0.0031} on holisticness@10
(Table~\ref{tab:power}). The ratio is the whole story --- a family's full range
is only ${\sim}2.3$ resolution widths wide on recall and ${\sim}1.9$ on
holisticness, so of eight adjacent arms roughly \textbf{two tiers} are separable
on Recall@10 and \textbf{one} on holisticness@10. Resolving a \num{0.005}
Recall@10 difference at this precision would require
$n \approx 1431\cdot(\num{0.0124}/\num{0.005})^2 \approx \num{8800}$ test
sessions, six times what the split provides
(Figure~\ref{fig:nexttrack-power}).

\begin{table}[ht]
  \centering
  \caption{Measured resolution of the test split. Half-widths are means over the
  nine paired-bootstrap comparisons of the pre-encoder batch (\num{2000}
  resamples, $n=\num{1431}$); the family spread is that batch's full range on
  each metric. ``Tiers'' is spread\,$/$\,half-width, i.e.\ how many genuinely
  distinguishable levels eight adjacent arms can occupy.
  \emph{Sources: \texttt{2026-07-25-nexttrack-dual-tower-fusion-scan.md},
  \texttt{2026-07-26-nexttrack-pre-encoder-scan.md}.}}
  \label{tab:power}
  {\footnotesize\setlength{\tabcolsep}{6pt}\begin{tabular}{lcccc}
    \toprule
    Metric & CI half-width at $n=\num{1431}$ & family spread & tiers & $n$ for $\Delta=\num{0.005}$ \\
    \midrule
    Recall@10 & 0.0124 & 0.0280 & $\sim$2 of 8 & $\sim$\num{8800} \\
    holisticness@10 & 0.0018 & 0.0031 & $\sim$1 of 8 & --- \\
    \bottomrule
  \end{tabular}}
\end{table}

\begin{figure}[ht]
  \centering
  \includegraphics[width=0.8\textwidth]{nexttrack_power}
  \caption{The architecture nulls are power-limited. The paired-bootstrap CI
  half-width falls as $1/\sqrt{n}$; at the split's $n=\num{1431}$ it is
  \num{0.0124} on Recall@10, while a whole architecture family spans only
  \num{0.028}. Resolving a \num{0.005} difference would need
  ${\sim}\num{8800}$ test sessions. The honest reading of the five architecture
  nulls follows directly: \emph{large harms were detected reliably, moderate wins
  were never detectable at all}.}
  \label{fig:nexttrack-power}
\end{figure}

\textbf{The honest reading of the five architecture nulls follows directly, and
it is not the reading a casual glance gives.} What these campaigns established
firmly is that the tested architectures do not deliver \emph{large} gains, and
that several of them do real \emph{harm} --- the harms are the well-powered
results, because they are big (S1 $-\num{0.0140}$; the \texttt{fusion\_layers}
regression $+\num{0.0266}$; six of eight dual arms below the baseline). What the
campaigns \emph{cannot} establish is the absence of a moderate gain: a genuine
$+\num{0.005}$ improvement --- which on a \num{0.123} baseline would be a
respectable \SI{4}{\percent} relative lift --- would have been invisible by
construction. So ``dual-tower fusion is refuted'' and ``the pre-encoder is
refuted'' should be read as \emph{no large effect, and no effect of the
hypothesized magnitude}, not as ``these architectures do not work.'' The
practical consequence is a prioritization rule rather than a caveat: on this
split, further architecture arms are a poor use of runs no matter how they come
out, because the design space's whole dynamic range is barely twice the
measurement error. Either enlarge $n$ or move to axes whose effects are large ---
which is, empirically, exactly what the representation and supervision axes have
delivered (the content projection's $+\num{0.0259}$ is twice the half-width, and
the second-split hardening's $+\num{0.0426}$ is three times it).

\subsection{Three generations of holisticness@10, and why the board is denominated in recall}
\label{sec:power-generations}
The composite \emph{holisticness@10} was revised twice in four days, and stored
values are \textbf{not comparable across the cuts}:
\begin{description}
  \item[Gen.\ 1 --- 3-factor (to 2026-07-24).] The bare product
  \[\mathrm{clamp}(\text{mood\_coh},0)\times\text{ild}\times(1-\text{artist\_adj}),\]
  carrying \emph{no correctness term at all}. Values land in
  \numrange{0.063}{0.233} --- the column in Table~\ref{tab:dualtower}.
  \item[Gen.\ 2 --- 4-factor, grounded (2026-07-25).] Multiplies in a
  \texttt{music@10} grounding factor,
  \[\mathrm{clamp}\bigl((\text{music@10}-\num{0.185})/\num{0.815},\,0,\,1\bigr).\]
  Because \texttt{music@10} falls in lockstep with recall, this cancels the
  de-eagering ``wins'' the 3-factor form rewarded; values land in
  \numrange{0.044}{0.050} --- the column in Table~\ref{tab:preencoder}.
  \item[Gen.\ 3 --- concentration-aware (2026-07-26).] Eagerness becomes
  \texttt{max(artist\_adj,} \texttt{artist\_conc)} and the grounding floor
  rises \num{0.185}$\to$\num{0.34}. Values land near \numrange{0.009}{0.025}.
\end{description}
Every facet is stored per row, so any generation is recomputable after the fact
(the helper \texttt{rescore\_holisticness.py} does this for the runs it can
reach),
and Table~\ref{tab:hgen} re-derives all three for a representative slice. Two
consequences matter for reading this document. First, \textbf{a holisticness
column must always be labelled with its generation}, because the generation
effect dwarfs the model effect: every 3-factor value mechanically outranks every
4-factor value, since the grounding factor multiplies by ${\sim}0.3$. Second,
and more importantly, \textbf{the recall and holisticness boards point in
opposite directions at every generation}. The champion blend --- \num{0.21174}
Recall@10, \num{303} of \num{1431} --- ranks \emph{last} in
Table~\ref{tab:hgen} under all three formulas, because it is the most
artist-eager model on record; the degenerate \texttt{cummean/cummean} arm and
\texttt{mood-session} lead. Generation 3 demoted \texttt{mood-session} from
first to a genuine zero, but it did \emph{not} fix the dual-tower degenerate arm,
so the $\pm\num{0.015}$ recall floor remains mandatory on any holisticness claim.

\begin{table}[ht]
  \centering
  \caption{The same models under all three holisticness generations, sorted by
  Recall@10. The generation effect dwarfs the model effect, and the ordering is
  roughly \emph{inverse} to Recall@10 throughout --- which is why this document's
  leaderboard stays denominated in Recall@10. Dashes are runs whose facets could
  not be re-derived under that generation (remote-only runs with no materialized
  predictions file). \emph{Sources:
  \texttt{2026-07-25-nexttrack-dual-tower-fusion-scan.md},
  \texttt{2026-07-26-nexttrack-pre-encoder-scan.md}.}}
  \label{tab:hgen}
  {\footnotesize\setlength{\tabcolsep}{2.4pt}\begin{tabular}{llccc}
    \toprule
    Model (run) & R@10 (raw) & gen-1 (3-f) & gen-2 (4-f) & gen-3 (current) \\
    \midrule
    \best{\Rp+M+\Cp blend} (\texttt{1c1a3}) & \best{0.21174 (303)} & 0.0632 & 0.020830 & 0.00860 \\
    single GRU h454 (\texttt{7646d}) & 0.12509 (179) & --- & 0.045765 & 0.02399 \\
    single GRU h256 (\texttt{97c1e}, \texttt{bdcbc}) & 0.12299 (176) & 0.14172 & 0.047086 & \best{0.02454} \\
    single GRU h425 (\texttt{be11b}) & 0.12089 (173) & 0.13850 & 0.046677 & 0.02487 \\
    S1 pre-MLP$\to$GRU (\texttt{ad53f}) & 0.10901 (156) & 0.14499 & 0.047132 & 0.02410 \\
    D-CC \emph{cum/cum} (\texttt{88a4c}) & 0.06639 (95) & 0.19715 & \best{0.050436} & --- \\
    \texttt{mood-session} (degen.) & 0.02380 (34) & \best{0.23319} & 0.043881 & 0.00000 \\
    \bottomrule
  \end{tabular}}
\end{table}

\paragraph{What is primary, and what this document sorts on.} For completeness:
the \texttt{[metrics].primary} key in \texttt{domain.toml}, and the best-models
group's \texttt{primary\_metric}, have both read \texttt{holisticness@10} since
2026-07-23, so the \emph{server-effective} primary metric is no longer
Recall@10. Statements
in the earlier sections that ``the primary metric stays Recall@10'' were true
when written and are superseded \emph{for the server}. This document
nevertheless continues to denominate its leaderboard
(Table~\ref{tab:lb}) and every campaign verdict in \textbf{Recall@10},
deliberately: it is the record's comparability currency, it is stable across all
three metric generations, and --- as Table~\ref{tab:hgen} shows --- a
holisticness ordering would place the servable champion last. Holisticness is
reported throughout as a labelled, generation-tagged diagnostic column.
\emph{Sources: \texttt{2026-07-25-nexttrack-dual-tower-fusion-scan.md},
\texttt{2026-07-26-nexttrack-pre-encoder-scan.md}.}

\section{The MMR $\lambda$ frontier --- the first confirmed crown lever}
\label{sec:mmr}
\emph{Rationale.} Section~\ref{sec:holisticness} left one standing flag: an
eval-time MMR re-rank at $\lambda\approx\num{0.9}$ was ``the only near-free
holisticness lever on record'', measured in an \emph{offline} sweep over a stored
checkpoint and explicitly deferred to a user decision. This campaign is the
native confirm, and it was designed to settle a second, pre-registered
\emph{mechanism} question rather than merely to re-measure the knob: does sonic
diversification de-concentrate \textbf{artists} (H1, ``responsive''), or are
same-artist tracks far enough apart in the content-metric space that the
diversity penalty leaves eagerness untouched (H2, ``sticky'')? The answer decides
whether an explicit artist-diversity penalty is the indicated sequel or a
redundant one.

\paragraph{What is being maximized.} All ten arms are scored under the
\emph{live} four-factor crown, taken verbatim from
\texttt{predictors/seq\_common.py}:
\[
  \begin{aligned}
    H = {}& \max(\texttt{mood\_coh},0)\;\times\;\texttt{ild}\\
          &\times\;\bigl(1-\mathrm{clamp}(\max(\texttt{artist\_adj},
             \texttt{artist\_conc}),0,1)\bigr)\\
          &\times\;\mathrm{clamp}\bigl((\texttt{music@10}-\num{0.34})/\num{0.66},
             \,0,\,1\bigr),
  \end{aligned}
\]
each facet the \emph{mean over rows} and the product taken \emph{of the
aggregates} --- generation 3 of Section~\ref{sec:power-generations}. Two anchors
frame the batch: \textbf{A}, the recall champion \Rp+M+\Cp\ (H \num{0.008603}, its
eagerness factor bound by \texttt{artist\_conc} \num{0.691653}), and \textbf{B},
the best \emph{single} model on the live crown, \texttt{gru-infonce-h256} (H
\num{0.024538}). The bar to beat is the all-time healthy-arm crown maximum
$H_{\mathrm{bar}} = \num{0.024872}$, held by an h425 capacity control
(\runid{run-20260725-143502-be11b-seq-nexttrack}). All arms ride the canonical
split \dsid{seq-20260715-131139}; no dataset work was needed.

\paragraph{Why this batch is unusually sharp.} $\lambda$ is an \emph{eval-only}
parameter: the trained weights are bit-identical across the arms of one model, so
$\lambda=\num{1.0}$ is a bit-identity gate and every $\Delta$ is a maximally
paired comparison. Both controls reproduced their anchors exactly (\num{303} and
\num{176} hits of \num{1431}, MRR agreeing to \num{3.1e-11}), and every arm's $H$
was \emph{recomputed} from the stored \texttt{predictions.json} under one code
path, so the per-row rounding bias cancels in every $\Delta$ (recomputed versus
stored agreed to \num{7.4e-8} and \num{2.1e-7}). A third gate --- MMR-active ---
required each $\lambda<1$ arm's \texttt{ild} to differ from its own control by far
more than \num{1e-6}, which is the registered detector for MMR silently no-opping
when the content-metric index is unreachable. No arm was a no-op.

\emph{Results.} Table~\ref{tab:mmr-arms} is the frontier, and
Figure~\ref{fig:nexttrack-mmr-frontier} plots it. Arms are tiered by a
pre-registered rule: \textbf{TIER-1 FREE} means $\Delta H$ above zero with the
recall $\Delta$ straddling zero and $|\Delta\text{recall}| \ll \num{0.015}$;
\textbf{TIER-2 PRICED} means a real $H$ gain at a recall cost the CI can see;
\textbf{TIER-3 REJECTED} means $\Delta$recall below $-\num{0.03}$.

\begin{table}[ht]
  \centering
  \caption{The MMR $\lambda$ frontier on both anchors, canonical split
  \dsid{seq-20260715-131139} ($n=\num{1431}$), sorted by holisticness@10. Every
  run id is \protect\runid{run-20260726-} plus a launch time and hex suffix, in
  arm order: \texttt{184115-bf395} (B3), \texttt{184115-1dc77} (A4),
  \texttt{184116-20d9b} (B4), \texttt{184115-147a6} (B2),
  \texttt{184115-e1539} (B1), \texttt{183333-9a0b2} (B0),
  \texttt{184115-7a2f7} (A3), \texttt{184115-36edb} (A2),
  \texttt{184115-6cb38} (A1), \texttt{183333-1c361} (A0); the A arms are
  \texttt{seq-blend}, the B arms \texttt{seq-nexttrack}. Every $\Delta$ is a
  paired bootstrap (\num{2000} resamples, rng \num{1337}) against the arm's
  \emph{own} $\lambda=\num{1.0}$ control, with all five facets re-aggregated on
  the same resampled index set before the product; CI bounds are given to three
  decimals for width. Tier codes: \emph{t1} TIER-1 FREE, \emph{t2} TIER-2 PRICED,
  \emph{t3} TIER-3 REJECTED, \emph{ctrl} the arm's own $\lambda=\num{1.0}$
  control. The binding facet --- the one entering the eagerness factor through the
  $\max()$ --- is \texttt{conc} on A0/A1/A2, B0 and \texttt{adj} everywhere else.
  H@10, R@10, \texttt{mus} ($=$ music@10) and \texttt{ild}
  higher-better; \texttt{adj}, \texttt{conc} lower-better. Best cell per metric
  bolded; the promotable arm B4 is bolded in full. \texttt{mood\_coh} is omitted
  for width and read in the interpretation: B0$\to$B3 \num{0.4788}$\to$\num{0.4585}
  ($-\SI{4}{\percent}$) against A0$\to$A4 \num{0.3982}$\to$\num{0.2378}
  ($-\SI{40}{\percent}$), the champion's collapse under diversification.
  \emph{Source: \texttt{2026-07-26-nexttrack-mmr-lambda-frontier.md}.}}
  \label{tab:mmr-arms}
  {\scriptsize\setlength{\tabcolsep}{1.9pt}\begin{tabular}{lllllllll}
    \toprule
    Arm & $\lambda$/pool & H@10 & $\Delta H$ [95\% CI] & R@10 (hits) & $\Delta$R@10 [95\% CI] & mus & ild & adj / conc \\
    \midrule
    B3 \emph{t2} & 0.7 / 200 & \best{0.034171} & $+.0096$ [$+.008$, $+.011$] & 0.11461 (164) & $-.0085$ [$-.016$, $-.001$] & 0.4747 & 0.5314 & 0.313 / 0.271 \\
    A4 \emph{t3} & 0.3 / 200 & 0.031626 & $+.0230$ [$+.021$, $+.026$] & 0.10971 (157) & $-.1019$ [$-.120$, $-.084$] & \best{0.4925} & \best{0.8146} & 0.293 / \best{0.190} \\
    \best{B4} \emph{t1} & \best{0.7 / 50} & \best{0.030355} & \best{$+.0058$ [$+.005$, $+.007$]} & \best{0.12020 (172)} & \best{$-.0029$ [$-.008$, $+.003$]} & 0.4670 & 0.5017 & 0.327 / 0.302 \\
    B2 \emph{t1} & 0.8 / 200 & 0.029458 & $+.0049$ [$+.004$, $+.006$] & 0.11950 (171) & $-.0036$ [$-.009$, $+.001$] & 0.4648 & 0.4961 & 0.330 / 0.307 \\
    B1 \emph{t1} & 0.9 / 200 & 0.026761 & $+.0022$ [$+.002$, $+.003$] & 0.12089 (173) & $-.0021$ [$-.006$, $+.002$] & 0.4597 & 0.4716 & 0.341 / 0.337 \\
    B0 \emph{ctrl} & 1.0 / 200 & 0.024539 & --- & 0.12299 (176) & --- & 0.4558 & 0.4536 & 0.347 / 0.356 \\
    A3 \emph{t3} & 0.5 / 200 & 0.023576 & $+.0150$ [$+.013$, $+.017$] & 0.16422 (235) & $-.0473$ [$-.061$, $-.033$] & 0.4871 & 0.6527 & 0.486 / 0.461 \\
    A2 \emph{t2} & 0.7 / 200 & 0.015289 & $+.0067$ [$+.006$, $+.008$] & 0.19357 (277) & $-.0181$ [$-.029$, $-.008$] & 0.4759 & 0.5201 & 0.571 / 0.608 \\
    A1 \emph{t1} & 0.9 / 200 & 0.010435 & $+.0018$ [$+.001$, $+.002$] & 0.21034 (301) & $-.0013$ [$-.008$, $+.005$] & 0.4625 & 0.4362 & 0.603 / 0.671 \\
    A0 \emph{ctrl} & 1.0 / 200 & 0.008603 & --- & \best{0.21174 (303)} & --- & 0.4538 & 0.4064 & 0.610 / 0.692 \\
    \bottomrule
  \end{tabular}}
\end{table}

\begin{figure}[ht]
  \centering
  \includegraphics[width=0.95\textwidth]{nexttrack_mmr_frontier}
  \caption{\textbf{Left:} the $\lambda$ frontier on the recall $\times$ crown
  plane. Each engine traces a path as $\lambda$ falls from \num{1.0}; colour is
  the pre-registered tier (green TIER-1 FREE, blue TIER-2 PRICED, red TIER-3
  REJECTED, grey control). Three single-GRU arms clear the prior all-time crown
  bar \num{0.024872} at a recall whose paired CI straddles zero --- the first
  genuinely free crown movement on record. The champion's path reaches comparable
  $H$ only after paying recall the CI can see, and its two best-$H$ points are
  rejected as degenerate. \textbf{Right:} the mechanism question. H1 required
  $|\Delta\texttt{artist\_conc}|\ge|\Delta\texttt{artist\_adj}|$; H2 required the
  ratio below \num{0.5} (shaded). Measured \numrange{1.59}{3.06} in \emph{8 of 8}
  arms, so H2 is refuted by a factor of \numrange{3}{6} in the wrong direction ---
  and the ratio shrinks monotonically as $\lambda$ falls, i.e.\ the easy
  de-concentration is bought first and the mechanism saturates.}
  \label{fig:nexttrack-mmr-frontier}
\end{figure}

\emph{Interpretation.} \textbf{MMR is a real crown lever, and the first one on
record.} Every one of the eight $\lambda<1$ arms lifts $H$ with the CI entirely
above zero. What makes the gain legitimate rather than an artefact of the
composite is that \emph{three} of the four factors improve at once: \texttt{ild}
rises (\num{0.4536}$\to$\num{0.5314} on the GRU, \num{0.4064}$\to$\num{0.8146} on
the champion), eagerness falls, and --- the load-bearing part ---
\texttt{music@10} \emph{also} rises (\num{0.4558}$\to$\num{0.4747}). A lever that
raised $H$ by degrading relevance would be cancelled by the grounding factor; this
one is not, which is precisely the distinction the floor was written to enforce
after the \texttt{cummean/cummean} and \texttt{seq-mood} pathologies.

\textbf{H1 is confirmed unanimously and H2 is refuted, which \emph{cancels} the
sequel the previous campaign had pre-registered.} In 8 of 8 arms
$\Delta\texttt{artist\_conc} \le \Delta\texttt{artist\_adj}$ (both negative) at a
ratio of \numrange{1.59}{3.06}. Same-artist tracks \emph{are} near-neighbours in
the content-metric space, so a max-similarity penalty de-concentrates artists as a
side effect, with no explicit notion of artist anywhere in it. Sonic MMR
therefore \emph{is} a de-eagering lever, and porting \texttt{seq-mood}'s explicit
\texttt{artist\_penalty} to the blend --- the pre-registered H2 follow-up --- is no
longer the indicated next move. It stays on file as a \emph{complement}, because
\texttt{artist\_adj} falls ${\sim}\num{2.5}\times$ slower than
\texttt{artist\_conc} and becomes the binding facet at low $\lambda$.

\textbf{The binding facet switches, and that locates the end of the lever.} On the
champion the switch from \texttt{artist\_conc} to \texttt{artist\_adj} happens
between $\lambda\num{0.7}$ and $\lambda\num{0.5}$; on the GRU it has already
happened by $\lambda\num{0.9}$. Once \texttt{artist\_adj} binds, the
$(1-\text{eager})$ factor is the stiff one and each further recall point buys
progressively less crown --- the structural reason the GRU's frontier flattens
where it does, and a prediction that pushing below ${\sim}\lambda\num{0.7}$ on the
GRU is poor value. B3 is the confirmation: the highest $H$ measured anywhere
(\num{0.034171}) at the cost of losing the recall tie.

\textbf{The promotable operating point is the pool probe.} \textbf{B4}
($\lambda\num{0.7}$, \texttt{mmr\_pool}~\num{50}) is the single best TIER-1 arm at
$H$~\num{0.030355} --- $+\SI{22.1}{\percent}$ over $H_{\mathrm{bar}}$ and
\num{5.4} measured half-widths clear of it --- and it beats B2 on $H$ by
$+\num{0.000889}$ with the CI above zero at a recall that straddles, so its rank
is \emph{resolved} rather than nominal. B2 ($\lambda\num{0.8}$) was the design's
pre-registered promotion hypothesis and fired almost exactly as projected
($\num{0.0321}$ predicted against \num{0.029458} measured, with a \emph{smaller}
recall cost than projected). Nothing was promoted: six definitions were
registered, each carrying the note ``NOT PROMOTED --- robustness confirm $+$ user
approval pending''. The recall board is untouched --- \textbf{no arm improves
recall}; MMR only ever costs it.

\textbf{The second axis, \texttt{mmr\_pool}, closes as a fine-tuning knob rather
than an axis.} Capping the candidate pool at \num{50} recovers recall
($\Delta$ vs pool~\num{200} $+\num{0.005584}$, CI above zero, as hypothesized ---
a cap bounds how far MMR can reach for diversity) but gives $H$ back
($-\num{0.003838}$, CI below zero). It does not win at fixed $\lambda$; it moves
along the \emph{same} frontier. What earns it a line in the record is the
cross-comparison above: the cap reaches a slightly \emph{better} point than the
nearest pure-$\lambda$ setting. Prefer $\lambda$ for coarse movement, pool for
fine-tuning at a chosen $\lambda$.

\textbf{The champion/GRU contrast resolves in favour of factor health, not
de-eagering headroom.} The design framed this as the campaign-carrying question:
the champion is $\num{1.9}\times$ more eager, so a given absolute de-eagering is
worth ${\sim}\num{2.1}\times$ more relative $H$ there --- yet it starts
$\num{2.85}\times$ behind, and the gap \emph{widens} under diversification
($H$~\num{0.015289} against \num{0.034171} at $\lambda\num{0.7}$). The facets say
why: the champion's \texttt{mood\_coh} \textbf{collapses} as it diversifies
(\num{0.3982}$\to$\num{0.2378} at $\lambda\num{0.3}$, $-\SI{40}{\percent}$) while
the GRU's barely moves ($-\SI{4}{\percent}$). The champion's recall comes from
artist-adjacency --- the Markov bigram leg, per
Section~\ref{sec:holisticness} --- so diversifying away from the seed artist also
walks it out of the prefix's mood neighbourhood. It pays twice. \textbf{MMR's
value scales with the health of the other three factors, not with room to
de-eager}, and no $\lambda$ closes the champion's crown gap.

\paragraph{Two instrument findings, one of them a warning about the record
itself.} First, \textbf{the shared-weights instrument is
$\numrange{3}{33}\times$ sharper than the architecture-family noise floor}:
measured half-widths run \numrange{0.000382}{0.002476} on $\Delta H$ against the
family's conservative \num{0.0018}, with the number of sessions whose hit@10
changed at all as low as \textbf{\num{9}}. Arm A1 had been pre-registered as
``likely UNRESOLVABLE'' and resolved at \num{4.8} measured half-widths.
Eval-time levers must therefore budget their minimum detectable effect from the
\emph{pairing}, not from the architecture figure of
Section~\ref{sec:power-stats} --- the two regimes differ by more than an order of
magnitude, and conflating them under-powers a whole class of cheap experiments.
The corollary is reassuring in the other direction: the 2026-07-22 \emph{offline}
$\lambda$ sweep projected recall \num{0.2103} / \num{0.1936} / \num{0.1642} /
\num{0.1097} and the native runs returned \num{0.21034} / \num{0.19357} /
\num{0.16422} / \num{0.10971} --- four for four to four decimals, which validates
the offline driver as a \emph{screening} instrument.

Second, a new infrastructure pitfall was found the hard way and it retroactively
qualifies part of the record: \textbf{a remote worker running a
\texttt{(deleted)} binary image silently strips metric keys its struct does not
know}. The metric-completeness gate failed on the first attempt with
\texttt{artist\_conc\_at\_k} missing, although the predictors were md5-identical
to the hub and the worker's own \texttt{metrics.json} carried the field: the
worker \emph{daemon} had started a day before the field existed, \texttt{rsync}
had replaced the file on disk but not the running image, and the old image
deserialized the predictor's metrics into a struct lacking the field, dropped it,
and posted the filtered set \textbf{without erroring}. \texttt{holisticness@10}
was nevertheless correct throughout, because the predictor computes $H$ itself
before the server sees it --- which is exactly what made the defect invisible.
Every remote run in a four-hour window on 2026-07-26 lost
\texttt{artist\_conc\_at\_k}, which is why anchor B's quoted concentration was a
\emph{reconstruction} from \texttt{top\_k\_ids} and the dataset artist map. That
reconstruction is now independently validated: the natively emitted values
(\num{0.6916530786551751} and \num{0.35591272614333413}) match the rescoring
tool's reconstructions to the last digit. Remedy: cycle the daemon whenever a
metric field is added; \texttt{rsync} of the binary is not sufficient.
\emph{Source: \texttt{2026-07-26-nexttrack-mmr-lambda-frontier.md}.}

\section{Autoregressive rollouts: a generation metric, and five sweeps that say no}
\label{sec:rollout}
\emph{Rationale.} This campaign was provoked by a \emph{user preference}, not by a
metric. On a real 135-track playlist the user strongly preferred one engine's
20-step journey over another's, while \texttt{holisticness@10} scored those two
journeys \num{0.003} apart. That gap is the point: every metric on record ---
Recall@10, holisticness@10 and its four facets, \texttt{music@10}, even the
multi-step continuation evaluator --- scores \emph{one} ranked list produced from a
\emph{real} prefix. None of them ever feeds a model its own output. A 20-step
journey does exactly that, so a trajectory failure --- a rollout whose energy
decays until a party seed has become a comedown --- is structurally invisible to
all of them. This is the first campaign in this document to build an instrument
for the \emph{deployment} regime, and it is the direct ancestor of the walk
surface of Section~\ref{sec:walk}.

\paragraph{What was built.} Four things, all off the training path.
\begin{description}
  \item[\texttt{hold@N}, a generation metric.] Per legible axis $a$ in
  \{energy, valence, acousticness, tempo\},
  $\texttt{fade}_a = |\mathrm{mean}(\text{2nd half}) -
  \mathrm{mean}(\text{1st half})|/\texttt{band}_a$, and then
  \[
    \texttt{hold@}N = 1 - \mathrm{clamp}\bigl(\mathrm{mean}_a(\texttt{fade}_a)
      + \texttt{stasis},\,0,\,1\bigr).
  \]
  The per-axis bands are calibrated on \num{2135} real
  \emph{train} sessions (energy fade p10 \num{0.0141}, p50 \num{0.0769}, p90
  \num{0.1963}; tempo \num{1.28}, \num{8.05}, \num{22.43}), so axes with
  incompatible units combine and ``a big fade'' means big \emph{relative to how
  this listener actually listens}. Real sessions are not flat --- the median energy
  fade is \num{0.077} --- so a metric rewarding zero fade would rank an
  unnaturally rigid generator above the listener it imitates.
  \item[\texttt{stasis}, the anti-gaming term.] A pure fade metric is
  degenerate: twenty near-identical tracks fade by \num{0} and score \num{1.0},
  which is the
  same pathology that forces \texttt{holisticness@10} to carry an \emph{external}
  recall floor. \texttt{stasis} puts the floor \emph{inside} the metric by
  penalizing steps smaller than the listener's own p10 jump, and it was verified
  on a synthetic flatline (moves-naturally \num{1.000}, at-your-floor \num{1.000},
  half-floor \num{0.500}, barely-moves \num{0.048}).
  \item[An offline rollout evaluator.] \texttt{seq\_continuation\_eval.py
  --rollout} scores sampled \emph{test} sessions by \textbf{importing} the
  generation loop from the serving path rather than reimplementing it, so an
  offline score describes the journeys the lab actually serves; the refactor was
  verified byte-identical. It is on-demand and sampled because a rollout costs
  \texttt{steps} full-vocabulary passes per session against the primary path's one
  (${\sim}20\times$).
  \item[A playlist-lab readout.] A mood report card, a comparative journey-shape
  chart and per-step \texttt{passed\_over} diagnostics, with the holisticness
  facets demoted to a collapsed detail block.
\end{description}

\emph{Results.} Table~\ref{tab:rollout} is five sweeps, \num{120} paired
rollouts and \num{240} journeys, under a paired bootstrap over \emph{seeds} plus a
leave-one-seed-out jackknife. Rollout policy was frozen throughout at the lab
defaults (\num{20} steps, artist penalty \num{2}, anchor \num{1}, temperature
\num{0} --- deterministic, so one seed is one sample and there is no rollout
variance).

\begin{table}[ht]
  \centering
  \caption{Five \texttt{hold@20} sweeps, paired bootstrap over seeds (\num{2000}
  resamples, rng \num{1337}) with a leave-one-seed-out jackknife (``LOO'' is the
  number of refits preserving the verdict). Positive $\Delta$ favours the
  dual-tower arm. The pre-registered rule was CONFIRMED iff the CI lies entirely
  above zero \emph{and} the verdict survives $\ge\SI{90}{\percent}$ of LOO refits;
  TIE iff the CI straddles. S2 and S3 met that rule and were nonetheless
  \emph{retracted} --- see the interpretation.
  \emph{Source: \texttt{2026-07-27-rollout-hold-metric.md}.}}
  \label{tab:rollout}
  {\footnotesize\setlength{\tabcolsep}{4pt}\begin{tabular}{llrlll}
    \toprule
    \# & Sweep & $n$ & $\Delta$\,\texttt{hold@20} [95\% CI] & LOO & Verdict \\
    \midrule
    S1 & pilot, mood-extreme seeds, 10-track prefix & 8 & $-0.034$ [$-0.138$, $+0.079$] & 0/8 & TIE \\
    S2 & pilot, mood-extreme seeds, 30-track prefix & 8 & $+0.314$ [$+0.107$, $+0.529$] & 8/8 & \emph{retracted} \\
    S3 & replication, \textbf{random} seeds & 39 & $+0.122$ [$+0.044$, $+0.201$] & 39/39 & \emph{confounded} \\
    S4 & \textbf{same-split control} & 39 & $+0.033$ [$-0.030$, $+0.093$] & 0/39 & \best{TIE} \\
    S4b & same-split, other dual checkpoint & 39 & $+0.051$ [$-0.012$, $+0.118$] & 0/39 & \best{TIE} \\
    S4c & \textbf{within-architecture} & 39 & $-0.072$ [$-0.141$, $-0.002$] & --- & significant \\
    S5 & \textbf{the user's own \num{11} playlists} & 11 & $+0.004$ [$-0.088$, $+0.104$] & 0/11 & \best{TIE} \\
    --- & the single judged playlist & 1 & $+0.379$ (\num{0.623} / \num{0.245}) & --- & outlier \\
    \bottomrule
  \end{tabular}}
\end{table}

\emph{Interpretation.} \textbf{Rollout fade does not separate this family.} At
matched training split the two architectures tie \emph{twice}; on the user's own
playlists the pair they preferred ties; and the only surviving effect is
\textbf{checkpoint-level} within one architecture (S4c), which does not transfer to
the playlist regime. Mean \texttt{hold@20} over the identical \num{39} random
seeds runs \num{0.6488} / \num{0.6309} for the two dual-tower checkpoints against
\num{0.5983} / \num{0.5266} for the two single-GRU ones --- so S3's apparent
architecture win decomposes into a non-significant architecture component
($+\num{0.05}$) and a significant \emph{checkpoint} component ($+\num{0.072}$).
Two \texttt{seq-nexttrack} models differ from each other more than dual-tower
differs from single-GRU at matched split, and the checkpoint the playlist lab
pre-selected by default happens to be the worst holder of the four. The user's
original comparison was the best holder against the worst.

Combined with the five prior architecture nulls this is a sixth, and it is the
strongest form available, because the standing rebuttal to those nulls was
``Recall@10 cannot see session quality''. It cannot; and when an instrument built
for the deployment regime looks, there is still nothing there.

\textbf{Per-seed variance dominates model identity, and that is the operational
finding.} Individual gaps on \emph{one} model pair span $+\num{0.39}$ to
$-\num{0.27}$, while the model-level mean difference on the user's playlists is
$+\num{0.004}$. Which seed you hand the generator matters far more than which of
these models generates. Two consequences: prefer serving an \emph{ensemble} and
letting the listener choose over hunting a single better generator, and budget
$\ge\num{40}$ seeds for any future rollout comparison --- \numrange{8}{12} cannot
resolve anything at this variance.

\textbf{\texttt{hold@N} is a working instrument but not a model ranker.} It fired
on a real checkpoint difference, it fired on the judged playlist (\num{0.62}
against \num{0.25}) where \texttt{holisticness@10} saw a \num{0.003} gap, and it
correctly returned TIE three times. The recommendation is therefore to keep it as
a per-journey \emph{lab readout} and \textbf{not} to wire it as a leaderboard
column: its ground truth is listener preference, which stands at $n=1$ blind vote.
Promoting an unvalidated composite to primary is how the crown ended up needing an
external recall floor, and replacing one unvalidated composite with another repeats
the mistake.

\paragraph{Four methodological pitfalls, each of which cost a retracted
positive.} They are worth more than the null. \textbf{(1) A cross-split comparison
masquerading as an architecture result.} The lab's two default pre-selects are
trained on \emph{different} splits with identical vocab and item latents, so
comparing them attributes checkpoint and split differences to architecture; pair
from the same batch when the claim is architectural. \textbf{(2) Mood-extreme seed
selection inflates effects ${\sim}\num{2.5}\times$} (S2's $+\num{0.314}$ became
S3's $+\num{0.122}$ on a random sample). \textbf{(3) An anti-gaming term can be
too weak to do its job:} \texttt{stasis} first normalized the shortfall by the
band \emph{width}, and because this listener's jump band is wide a flatlined
rollout scored \num{0.740} --- \emph{above} a real generator's \num{0.477}.
Normalize by the \emph{floor}; a guard beatable by the degeneracy it targets
launders that degeneracy as a good score. \textbf{(4) A metric that returns
nothing must drop the \emph{pair}} and report the smaller $n$, or the pairing
silently breaks. The jackknife earned its keep at $n<40$ throughout: it killed two
of three per-axis positives at $n=8$ and flagged S4 and S5 as unsupported at 0/39
and 0/11.

\paragraph{What the per-axis read adds, including a failed negative control.} At
$n=\num{39}$ only \emph{energy} fade separates the arms ($-\num{0.039}$
[$-\num{0.069}$, $-\num{0.009}$], favouring dual-tower); valence, acousticness and
tempo all straddle. But one pre-registered \emph{negative} control moved: the
energy \emph{level} deviation from the seed favours the single GRU
($+\num{0.018}$ [$+\num{0.005}$, $+\num{0.032}$]). Those are different properties,
so the claim had to be re-scoped from ``holds the seed better'' to ``is more
stationary'' --- and it shows that \texttt{hold@N}'s deliberate exclusion of a level
term is a choice that \emph{changes the ranking}, not a neutral simplification.
Two other limits travel with the batch: \SI{74}{\percent} of the corpus carries
the acoustic fields the metric needs, so seeds were pre-filtered for
$\ge\SI{60}{\percent}$ prefix coverage (a mild selection effect), and nothing here
has been listened to.

Nothing was promoted, registered or wired; the champion, the crown board and the
best-models pins are unchanged. For the record, the live-crown standing relevant
to the dual-tower question is a \emph{tie} rather than a loss --- the best
legitimate dual arm \num{0.02378} against the single GRU's \num{0.02454}, a
\num{0.0008} gap against a ${\sim}\num{0.0018}$ half-width. The family was refuted
on Recall@10, ties on the live primary, and shows nothing on the rollout axis.
The indicated next axis is therefore not topology but the \emph{training regime}:
all five architecture nulls held one-step teacher forcing byte-identical while the
product rolls twenty steps on its own output, and that exposure-bias mismatch is
the one thing those nulls could not have detected.
\emph{Source: \texttt{2026-07-27-rollout-hold-metric.md}.}

\section{A research memo, and what its provenance permits}
\label{sec:memo}
\emph{Status, first.} This section consolidates a document that is \textbf{not a
campaign report}. Nothing was run on the server, nothing was promoted, nothing was
registered; the champion, the leaderboard, the crown board and the best-models
pins are all unchanged by it. It was produced by a research workflow (four scouts,
five proposal lenses, five adversarial screeners, then synthesis) plus
\emph{offline} scoring against frozen checkpoints. Where it reports a number, that
number came from an offline screen over stored artifacts, not from a native run.
It is included here because it is the direct parent of the campaign in
Section~\ref{sec:markovgate} and because its negative results --- the graveyard
below --- are cheaper to reuse than to rediscover. It is reported at a lower
register than the measured campaigns, and its figures are deliberately kept out of
every board in this document.

\paragraph{The provenance partition, kept visible.} The memo itself distinguishes
two classes of number, and that distinction survives into this synthesis.
\begin{description}
  \item[Independently re-derived by the orchestrator.] The share of test queries
  whose last prefix item has \emph{zero} train-bigram support:
  \textbf{\num{795} of \num{1431} $=$ \num{0.5556}}, reproduced by an independent
  script. The ungated control, which is self-validating because it must equal the
  champion: \textbf{\num{0.21174}}, exact. The gated arm: \textbf{\num{0.23201}},
  a paired $\Delta$ of $+\num{0.02027}$ $[+\num{0.01048}, +\num{0.03075}]$ with the
  CI above zero. And its bucket decomposition: the zero-support bucket moves
  \num{0.1660}$\to$\num{0.2025} over \num{795} queries while the two warm buckets
  are \textbf{bit-identical} (\num{0.3116} and \num{0.2326}, unchanged).
  \item[\textbf{Not} verified by the orchestrator --- agent-reported only, treated
  as provisional.] The $+\num{0.01957}$ second-split figure, the replay channel's
  $+\num{0.02865}$, the \texttt{latent\_gamma} numbers, and the recomputed
  \num{0.0706} discordance with its ${\sim}\num{0.0145}$ noise floor. None of these
  appears in Table~\ref{tab:lb} or on any crown board, and the first of them is
  precisely what the follow-up campaign was commissioned to establish natively.
\end{description}

\paragraph{One caveat the memo states about its own headline.} The ten gating
variants were screened \textbf{on the test set}, so quoting the best arm's CI is
optimistic by construction and the screen cannot itself be evidence. Three things
mitigate it --- four mechanistically related variants all land at
\numrange{0.019}{0.020}, the bucket decomposition is a mechanical explanation
rather than a fit, and the warm buckets do not move by a single hit --- but the
mitigation is an argument, not a measurement. The native two-split confirm is the
real gate, and it is Section~\ref{sec:markovgate}.

\paragraph{The honest headline: ``a better topology'' is the wrong question.} The
record had by then answered it six times --- cell and loss on the AE-64 space;
depth, width and direction; width on PCA-192; dual-tower fusion
(Section~\ref{sec:dualtower}); the MLP fusion head, an \emph{active} regression
recovered by deleting it; the per-step pre-encoder, a measured loss with the
mechanism identified down to \texttt{frac\_exact\_zero} $=$
\texttt{frac\_negative\_coords} (Section~\ref{sec:preencoder}); and rollout fade
(Section~\ref{sec:rollout}). Five lenses were pointed at the topology question
and \emph{every} architecture-shaped proposal came back weakened, sub-noise, or
pre-killed. Two came back with measured negative results: a split Markov leg lost
at every point of its own pre-registered grid, and a factored artist softmax head
was killed by its own pre-kill run, which found that top-5 artist gating on the
real checkpoint changes recall by \textbf{exactly \num{0.0}} because the trained
GRU's top-10 already lives inside its own top-5 artists.

\textbf{Where the evidence points instead: the model is not the constraint. What
it is asked to rank against, and which legs are allowed to speak on a given
query, are.} Every proposal that survived adversarial screening is a candidate-set
or evidence-conditioning change --- an evidence-gated Markov leg, a causal
cross-session replay channel, per-leg reliability scalars, a prefix-only negative
mask, the never-swept softmax temperature. None is a topology. This is the same
conclusion the record had reached from the representation side, arrived at from
five independent directions, and it comes with a measurement rationale: eval-time
and score-time levers are $\numrange{3}{33}\times$ cheaper to resolve than
training-time ones (Section~\ref{sec:mmr}), so a shortlist of score-time
proposals is the rational response to the instrument this lab has.

\paragraph{The two survivors worth naming.} \textbf{(1) \texttt{markov\_gate}} ---
gate the Markov leg to abstain where it has no bigram evidence, restoring the
\texttt{trust} term that the blend's z-normalization erases. Twelve lines,
zero learned parameters, and confirmed natively in the next section.
\textbf{(2) A causal cross-session replay channel} --- \SI{76.17}{\percent} of the
held-out targets were played in a \emph{strictly earlier} session and are not in
the current prefix, and nothing in the family can see this, because the Markov leg
is a train-frozen bigram and the recency scorer only ever scored prefix items,
which the evaluator masks to $-\infty$ (its recorded \num{0.001} is dead code, not
a measurement). Its hit-Jaccard with the GRU is \num{0.060} against \num{0.201}
for the GRU/Markov pair that produced a working blend, i.e.\ it is genuinely
different signal. Its catch is \emph{governance}, not code: it is an
information-set amendment that moves the rankable ceiling from \num{0.4493} to
\num{0.7617}, which needs its own board column and an explicit decision before any
number from it is comparable with anything in this document. A third proposal,
\texttt{latent\_gamma} (a one-scalar spectral rescale of the item latents), is
worth recording for its \emph{reframe} rather than its delta: on the offline
screen the champion's entire recall contribution is reproducible by a
\num{192}-number diagonal rescale, which recasts ``the single biggest win on
record'' as a conditioning fix that supervision happened to find, and predicts
that further linear metric learning is near-exhausted. All of these figures are in
the provisional class.

\paragraph{The graveyard, which is the reusable part.} Each entry is a lane that
should not be re-opened without new information. Full-catalogue or
support-restricted cross-entropy is pre-empted by an existing arm at
${\sim}\SI{53}{\percent}$ support coverage (\num{0.1223} against \num{0.1230}) ---
the literature's top candidate does not transfer because the item side is frozen.
A factored artist softmax head is refuted by its own pre-kill (exactly \num{0.0}).
Splitting the Markov leg into track and back-off components was measured at every
grid point and lost at all five ($-\num{0.0217}$ to $-\num{0.0384}$, all CIs below
zero), because separate z-normalization structurally destroys the \texttt{trust}
term inside the warm regime. A pool-matched stage-2 re-ranker is unobtainable: the
train pool's positive density is \num{0.5975} against test's \num{0.3368}, so its
own go-gate passes \emph{because of} the confound. A mixture-of-vMF head loses to
its own oracle decomposition. Diffusion, DreamRec and flow-matching are
externally falsified for this setting, and generative retrieval with semantic IDs
is strictly dominated when exact cosine over \SI{15}{\mega\byte} is already
affordable and the corpus is \SIrange{55}{64}{\percent} cold. Popularity
debiasing has nothing to debias (the popularity baseline is \num{0.0007}).
Hard-negative mining is harmful at low temperature here, because the hardest
negatives are same-album and same-artist tracks --- i.e.\ mining false positives.
And the GRU's pre-normalization output norm is not a confidence signal
(correlation with correctness $+\num{0.0073}$, and the loss is exactly
norm-invariant); the top-1-minus-next-10 score margin is, at $+\num{0.294}$.

\paragraph{The one durable theoretical result.} The memo's cross-cutting
observation that \textbf{the learned-combiner anomaly is mechanically explained}
is the part that changes how the record should be read. A prior oracle search over
\emph{global static scalar} leg weights returned $+\num{0.004}$ straddling zero,
and it is now clear why that is compatible with a real $+\num{0.020}$ from a
gate: dropping the Markov leg \emph{globally} gains in the zero-support bucket and
loses in the warm bucket, and the two cancel almost exactly. No constant can
express an intervention whose sign flips with a per-query data statistic, and
z-normalization is scale-invariant so it provably cannot see one. The registered
``four independent learned-combiner nulls'' pitfall should therefore forbid
\emph{learned score-level fusion} --- fitted weights, stackers, rank fusion,
end-to-end representation fusion --- and \textbf{not} per-query evidence gating,
which is a third thing: nothing is learned, nothing co-adapts, and the condition
is a train-side count.
\emph{Source: \texttt{2026-07-27b-nexttrack-model-topology-research-memo.md}
(research memo; see the provenance partition above).}

\section{\texttt{markov\_gate}: a leg that abstains where it has no evidence}
\label{sec:markovgate}
\emph{Rationale.} This is the native two-split confirm of the memo's headline
proposal (Section~\ref{sec:memo}), and it produces the largest single recall gain
in the whole next-track arc from a nine-line change with no new parameters.

\paragraph{The mechanism, and why it is not the refuted combiner lane.} The Markov
leg mixes its bigram row with an artist/genre back-off through a trust weight
$\texttt{trust} = n_u/(n_u+8)$, where $n_u$ is the train-level bigram support of
the query's last prefix item $u$. At $n_u = 0$ the trust is zero and the leg is
\textbf{pure back-off}: it carries no track-level evidence at all. But the blend's
candidate normalizer is \emph{scale-invariant} by construction --- it maps each
leg's score vector to unit variance so that three legs on wildly different scales
are commensurable --- so the evidence-free back-off is re-inflated and speaks at a
full one-third weight anyway. On the canonical split that happens on \textbf{\num{795}
of \num{1431} test queries}, \SI{55.6}{\percent}. The intervention is to make the
leg \emph{abstain} at $n_u = 0$ and renormalize the divisor from three legs to two.
Nothing is learned; the condition is a train-side count. That is exactly the third
category the registered learned-combiner pitfall did not cover, and the bucket
decomposition below is a mechanical proof that no global static weight could find
the effect: dropping the leg globally is a $-\num{0.0070}$ tie, because the
zero-support gain and the warm-bucket loss cancel.

\paragraph{Design and the preflight gate.} Six runs on the hub, all
\texttt{seq-blend} from the champion definition, with the controls
\textbf{re-trained in this batch} rather than substituted from the archive so that
every paired comparison is code-matched to one generation. Two splits: the
canonical \dsid{seq-20260715-131139} and the disjoint earlier-holdout
\dsid{seq-20260718-211238} (\SI{64}{\percent} cold, $n_{\mathrm{test}}$
\num{1431}). The hard pre-registered stop was that the ungated control had to
reproduce Recall@10 \num{0.21174}, exactly \num{303} of \num{1431}, proving the new
flag left the champion bit-identical. \textbf{It did, to the last digit}
(\num{0.21174004192872117}).

\emph{Results.} Table~\ref{tab:markovgate} is the batch,
Table~\ref{tab:markovgate-paired} the paired statistics, and
Table~\ref{tab:markovgate-buckets} the decomposition that carries the argument.
Figure~\ref{fig:nexttrack-markov-gate} plots both.

\begin{table}[ht]
  \centering
  \caption{The \texttt{markov\_gate} batch, six runs from the champion definition
  \texttt{blend-gru-markov-content-proj}. \dsid{ds1} $=$
  \dsid{seq-20260715-131139} (canonical, \SI{55}{\percent} cold), \dsid{ds2} $=$
  \dsid{seq-20260718-211238} (disjoint, \SI{64}{\percent} cold);
  \dsid{ds2}'s absolute recall is lower \emph{by construction}, so read it only
  through the paired $\Delta$. Every run id is \protect\runid{run-20260728-} plus
  the value in the \emph{run} column. The graded metrics track the gate too:
  artist@10 \num{0.333}$\to$\num{0.367} and genre@10 \num{0.401}$\to$\num{0.463}
  from C1 to A1, and \num{0.382} / \num{0.497} on A3.
  Rows sorted by Recall@10. H@10 is generation~3
  (Section~\ref{sec:power-generations}). R@10, MRR, H@10, mus@10, \texttt{ild},
  \texttt{mood}, art@10, gen@10 higher-better; \texttt{adj}, \texttt{conc}
  lower-better. Best cell per metric bolded among the canonical-split arms; the
  new best-recall row A1 is bolded in full.
  \emph{Source: \texttt{2026-07-27c-nexttrack-markov-gate-confirm.md}.}}
  \label{tab:markovgate}
  {\scriptsize\setlength{\tabcolsep}{1.9pt}\begin{tabular}{llllllllllll}
    \toprule
    id & run & ds & overrides & R@10 (hits) & MRR & H@10 & mus@10 & ild & adj & conc & mood \\
    \midrule
    \best{A1} & \texttt{005349-507bb} & ds1 & \best{gate} & \best{0.23201 (332)} & \best{0.13352} & 0.018699 & 0.49319 & 0.43377 & 0.57477 & 0.59821 & \best{0.46226} \\
    A3 & \texttt{012030-6d1b7} & ds1 & gate $+$ $\lambda$0.7 & 0.21314 (305) & 0.12852 & \best{0.030980} & \best{0.51649} & \best{0.56770} & \best{0.51279} & \best{0.47325} & 0.41885 \\
    C1 & \texttt{004444-b32ea} & ds1 & --- (champion) & 0.21174 (303) & 0.12199 & 0.008603 & 0.45379 & 0.40637 & 0.60971 & 0.69165 & 0.39822 \\
    C3 & \texttt{011245-bfc67} & ds1 & $\lambda$0.7 only & 0.19357 (277) & 0.11560 & 0.015289 & 0.47592 & 0.52010 & 0.57135 & 0.60823 & 0.36433 \\
    A2 & \texttt{010658-39e20} & ds2 & gate & 0.19008 (272) & 0.11801 & 0.008414 & 0.42953 & 0.41401 & 0.58672 & 0.64847 & 0.42618 \\
    C2 & \texttt{010113-e2599} & ds2 & --- (champion) & 0.17051 (244) & 0.10498 & 0.003392 & 0.39616 & 0.40177 & 0.61013 & 0.71945 & 0.35364 \\
    \bottomrule
  \end{tabular}}
\end{table}

\begin{table}[ht]
  \centering
  \caption{Paired bootstrap over the \num{1431} test sessions (\num{2000}
  resamples, rng \num{1337}, paired by row id). $\Delta H$ is computed
  \emph{product-of-means per resample} --- all five facets re-aggregated on the same
  resampled index set for both arms, then multiplied --- not as a mean of per-row
  products. $n_{disc}$ counts sessions whose hit@10 changed at all, split by sign.
  A verdict of WIN means the recall CI lies entirely above zero, TIE that it
  straddles; bounds are abbreviated to four decimals for width.
  \emph{Source: \texttt{2026-07-27c-nexttrack-markov-gate-confirm.md}.}}
  \label{tab:markovgate-paired}
  {\scriptsize\setlength{\tabcolsep}{4pt}\begin{tabular}{llllr}
    \toprule
    Comparison & $\Delta$R@10 [95\% CI] & Verdict & $\Delta H$ [95\% CI] & $n_{disc}$ ($+$/$-$) \\
    \midrule
    A1 $-$ C1 (gate, ds1) & \best{$+.02027$ [$+.0105$, $+.0308$]} & \best{WIN} & $+.010096$ [$+.0085$, $+.0117$] & 57 (43/14) \\
    A2 $-$ C2 (gate, ds2) & \best{$+.01957$ [$+.0105$, $+.0301$]} & \best{WIN} & $+.005022$ [$+.0039$, $+.0062$] & 56 (42/14) \\
    A3 $-$ C3 (gate, $\lambda$0.7) & \best{$+.01957$ [$+.0091$, $+.0301$]} & \best{WIN} & $+.015691$ [$+.0136$, $+.0179$] & 64 (46/18) \\
    A3 $-$ C1 (gate$+$MMR / champ.) & $+.00140$ [$-.0119$, $+.0147$] & TIE & \best{$+.022377$ [$+.0199$, $+.0249$]} & 96 (49/47) \\
    C3 $-$ C1 (MMR $\lambda$0.7 alone) & $-.01817$ [$-.0294$, $-.0077$] & costs recall & $+.006686$ [$+.0057$, $+.0077$] & 60 (17/43) \\
    \bottomrule
  \end{tabular}}
\end{table}

\begin{table}[ht]
  \centering
  \caption{The bucket decomposition, reproduced \emph{natively} through the real
  transition builder and bucketed by the train-level bigram support $n_u$ of each
  query's last prefix item --- the exact quantity the Markov scorer computes.
  Recall@10 within bucket, hits in parentheses. In both warm buckets the two arms'
  top-$k$ lists are \textbf{byte-for-byte identical}, so the gate is a surgically
  pure $n_u=0$ intervention with literally zero warm-side cost, on three
  independent control/arm pairs.
  \emph{Source: \texttt{2026-07-27c-nexttrack-markov-gate-confirm.md}.}}
  \label{tab:markovgate-buckets}
  {\footnotesize\setlength{\tabcolsep}{5pt}\begin{tabular}{llrllll}
    \toprule
    Pair & Bucket & $n$ & control & gated & $\Delta$ & top-$k$ identical? \\
    \midrule
    \multicolumn{7}{l}{\emph{Canonical split, A1 vs C1}} \\
    & $n_u = 0$ & 795 & 0.1660 (132) & 0.2025 (161) & \best{$+0.0365$} & no \\
    & $n_u$ 1--10 & 292 & 0.3116 (91) & 0.3116 (91) & $+0.0000$ & \best{YES} \\
    & $n_u > 10$ & 344 & 0.2326 (80) & 0.2326 (80) & $+0.0000$ & \best{YES} \\
    \midrule
    \multicolumn{7}{l}{\emph{Second split, A2 vs C2}} \\
    & $n_u = 0$ & 892 & 0.1244 (111) & 0.1558 (139) & \best{$+0.0314$} & no \\
    & $n_u$ 1--10 & 322 & 0.2578 (83) & 0.2578 (83) & $+0.0000$ & \best{YES} \\
    & $n_u > 10$ & 217 & 0.2304 (50) & 0.2304 (50) & $+0.0000$ & \best{YES} \\
    \midrule
    \multicolumn{7}{l}{\emph{Canonical split under MMR $\lambda$0.7, A3 vs C3}} \\
    & $n_u = 0$ & 795 & 0.1472 (117) & 0.1824 (145) & \best{$+0.0352$} & no \\
    & $n_u$ 1--10 & 292 & 0.2945 (86) & 0.2945 (86) & $+0.0000$ & \best{YES} \\
    & $n_u > 10$ & 344 & 0.2151 (74) & 0.2151 (74) & $+0.0000$ & \best{YES} \\
    \bottomrule
  \end{tabular}}
\end{table}

\begin{figure}[ht]
  \centering
  \includegraphics[width=0.95\textwidth]{nexttrack_markov_gate}
  \caption{\textbf{Left:} the bucket decomposition on all three control/arm pairs
  (\texttt{ds1} canonical, \texttt{ds2} the \SI{64}{\percent}-cold second split,
  \texttt{ds1 $\lambda$0.7} the same canonical pair under an MMR re-rank). The
  entire effect sits in the zero-support bucket; in both warm buckets the arms'
  top-$k$ lists are byte-for-byte identical, which is why a global leg weight
  provably cannot express this intervention. \textbf{Right:} the same result on the
  recall $\times$ crown plane. MMR alone (blue) is a \emph{priced} lever on the
  champion --- it buys crown and loses recall with the CI below zero. The gate
  (green) buys \emph{both}. Composed, gate$+$MMR lands at champion recall
  ($\Delta$ $+\num{0.00140}$, straddling) with $\num{3.6}\times$ the champion's
  holisticness, so the gate changes MMR's tier on the champion from PRICED to FREE.}
  \label{fig:nexttrack-markov-gate}
\end{figure}

\emph{Interpretation.} \textbf{The defect is real, it belongs to the normalizer,
and it is worth \num{29} hits.} Scale-invariance is exactly what makes the z-blend
robust to legs on incommensurable scales, and exactly what destroys their
\emph{confidence}. A Markov leg that has never seen the current track emits a pure
artist/genre back-off; the normalizer re-inflates it to the same footprint as a GRU
that has actually modelled the prefix, and the blend averages the two as equals.
Turning that third of the vote off where it is uninformed converts \num{43} misses
into hits and \num{14} hits into misses --- net $+\num{29}$, $+\SI{9.6}{\percent}$
relative, on the best model on record. Recall@10 \num{0.23201} (\num{332} of
\num{1431}) is \textbf{the best recall figure in this document}, with the paired CI
above zero on two disjoint splits and MRR up on both. For scale: the champion's own
coronation margin over the prior crown was $+\num{0.0259}$ canonical and
$+\num{0.0426}$ on the second split; this gate adds $+\num{0.0203}$ /
$+\num{0.0196}$ \emph{on top of that}, at no additional training cost.

\textbf{The gate is also a crown lever, which was not predicted.} $\Delta H$ is
above zero in all three paired comparisons, and the facets say why: abstaining
kills the \emph{bigram-echo} component of the recommendation. Artist concentration
falls \num{0.69165}$\to$\num{0.59821} and adjacency \num{0.60971}$\to$\num{0.57477}
--- the list stops repeating one artist --- while \texttt{mood\_coh} \emph{rises}
\num{0.39822}$\to$\num{0.46226} and \texttt{music@10} rises
\num{0.45379}$\to$\num{0.49319}. All four crown factors move the right way at
once. This is the first lever on record that de-eagers the champion \emph{without}
paying in mood coherence, and it resolves the exact double-payment
Section~\ref{sec:mmr} identified as structural: there, de-eagering was a re-rank
fighting the model's own scores, so walking away from the seed artist also walked
the blend out of the prefix's mood neighbourhood. Here the de-eagering is the
removal of a vote that had no business being cast.

\textbf{Gate and MMR are additive, not substitutive, and their composition is the
headline for the crown board.} The gate-alone $\Delta$ ($+\num{0.02027}$) and the
gate-under-$\lambda\num{0.7}$ $\Delta$ ($+\num{0.01957}$) are
indistinguishable, and the buckets show why: MMR re-ranks the top of a scored list,
the gate changes which list gets scored, and they act on disjoint parts of the
pipeline. Standing alone, MMR $\lambda\num{0.7}$ on the champion is TIER-2 PRICED
($-\num{0.01817}$ recall, CI below zero). Add the gate and the recall cost is paid
back exactly: A3 is $+\num{0.00140}$ $[-\num{0.01188}, +\num{0.01468}]$, a
statistical tie with the champion, at $H$~\num{0.030980} --- $\num{3.6}\times$ the
champion's \num{0.008603}, comfortably above the mandatory \num{0.108} recall
floor. On the crown board that sits at \textbf{rank 2}, behind only the
$\lambda\num{0.7}$ single GRU, which scores its \num{0.034171} at recall
\num{0.11461}: A3 delivers \SI{96}{\percent} of the incumbent leader's crown at
$\num{1.86}\times$ its recall.

\textbf{The honest caveat, and why the out-of-sample structure redeems it.} The
memo screened ten gating variants on the test set, so its banked
$+\num{0.02027}$ was selection-optimistic and could not be evidence. Four things
redeem it here: the effect reproduces on a \emph{second, disjoint,
\SI{64}{\percent}-cold split that was not screened on}, within \num{0.0007} of the
canonical figure; it reproduces under a third condition ($\lambda\num{0.7}$) that
was not screened either; four mechanistically related variants in the memo all
landed at \numrange{0.019}{0.020}, so the screen was measuring a plateau rather
than selecting a lucky corner; and the warm buckets do not move by a single hit,
which is a \emph{structural} prediction of the mechanism rather than a fitted
outcome. One residual risk remains: the $n_u = 0$ \emph{boundary} itself was chosen
with test knowledge. The boundary-free version --- scaling the leg by
\texttt{trust} \emph{before} normalization instead of hard-abstaining --- is the top
follow-up, and it is strictly more general. Nor is any of this a fluke of one
split's cold rate: the second split is \SI{64}{\percent} cold, its zero-support
bucket grows to \SI{62.3}{\percent} of queries while the per-query gain shrinks
slightly, and the net $\Delta$ comes out nearly identical. The lever's value scales
with how much of the query distribution has no bigram evidence, which is the
majority on any split of this corpus.

\paragraph{Governance, and an infrastructure finding that nearly ate the result.}
Nothing was promoted, registered or tagged; by design boundary this campaign
carries a promotion \emph{recommendation} only, and the board row in
Table~\ref{tab:lb} is labelled accordingly. The server's deterministic
top-12-by-holisticness recompute nevertheless auto-promoted four of the batch's
runs, which is server behaviour rather than a campaign action and left the group
needing re-curation. Separately: the remote training worker was live with a
\textbf{stale checkout containing zero occurrences of the new flag}. Had any gated
arm been dispatched there it would have executed the \emph{ungated} path and
silently returned a null --- the registered ``remote checkout can silently diverge''
pitfall, in its worst form, because the failure is a wrong number rather than a
crash. It could not happen here only because remote dispatch is opt-in and every
run in this batch was posted without it (confirmed by an unclaimed marker on all
six). Until the checkout is freshened, any queued \texttt{seq-blend} run silently
ignores the gate.

\paragraph{Record amendments.} Two, both narrowing rather than overturning.
\textbf{(1)} The learned-combiner pitfall should forbid \emph{learned score-level
fusion} and not \emph{per-query evidence gating}; the prior oracle null and this
$+\num{0.020}$ win are both true and are reconciled by the bucket decomposition,
since a global static scalar cannot express a gate whose sign flips with $n_u$.
\textbf{(2)} The normalizer's erasure of per-leg \emph{confidence} is a
\textbf{general} defect, of which the Markov leg's zero-support back-off is one
instance. Any leg with a query-dependent notion of ``I have no evidence here'' is
mis-weighted the same way, and the content and GRU legs have not been audited for
it.
\emph{Source: \texttt{2026-07-27c-nexttrack-markov-gate-confirm.md}.}

\section{The eagerness regularizer at a calibrated margin: refuted, with a mechanism}
\label{sec:eagermargin}
\emph{Rationale.} Section~\ref{sec:holisticness} left \texttt{eager\_beta} --- the
only \emph{training-time} de-eagering lever in the repository --- on record as ``a
weak knob, not recommended''. That verdict deserved a re-test, because the probe
behind it used a margin of zero at a dose worth \SIrange{1}{3}{\percent} of the
loss. This campaign asked whether the lever works when the margin is placed in the
measured \emph{right tail} of the transition-cosine distribution at
$\numrange{15}{60}\times$ that dose. The stake was deployability: at the time,
\texttt{mmr\_lambda} entered only through the evaluation path and never reached
serving, so the training-time regularizer was the one anti-eager candidate that
could actually ship.

\paragraph{The design was written so that a null would still be informative.} The
pre-registered primary was the paired $\Delta$ on \texttt{artist\_conc@10} against
an in-batch control, with holisticness@10 explicitly \emph{demoted} to secondary
because at the operating point $\partial H/\partial\texttt{music@10} =
+\num{0.2062}$ against $\partial H/\partial\texttt{eager} = -\num{0.0378}$ --- a
$\num{5.46}\times$ ratio that would have scored a \emph{successful} de-eagering as
approximately zero. A relevance floor was pre-registered in two parts: absolute
\texttt{music@10} $\ge\num{0.4300}$ \emph{and} a paired $\Delta$\texttt{music@10}
lower bound above $-\num{0.030}$. Eight \texttt{seq-dualgru} runs on the canonical
split \dsid{seq-20260715-131139}, all remote, \SI{60}{\minute} of compute, none
failed. The bit-identity gate passed at exactly \num{170} hits with the
reconstructed concentration matching the on-record arm to \num{15} decimal places,
and a penalty-active gate was pre-registered as a silent-failure detector: it
\textbf{passed on three arms and FAILED --- reporting the penalty inert --- on
four}, which is the campaign's result.

\emph{Results.} Table~\ref{tab:eagermargin} is the batch and
Table~\ref{tab:eagermargin-paired} the paired statistics.
Figure~\ref{fig:nexttrack-eager-margin} shows the mechanism.

\begin{table}[ht]
  \centering
  \caption{The \texttt{eager\_beta} $\times$ margin batch, eight
  \texttt{seq-dualgru} runs on \dsid{seq-20260715-131139} ($n=\num{1431}$). Every
  run id is \protect\runid{run-20260730-} plus the value in the \emph{run} column.
  $m$ is the hinge margin, $\beta$ the penalty weight. H@10 is generation~3
  (Section~\ref{sec:power-generations}). ``best ep/ran'' is the best epoch over
  the epochs actually run; Recall@10 is \texttt{hits}$/\num{1431}$ throughout
  (C0 \num{0.118798}, A1 \num{0.053809}, A2 \num{0.013976}, B1 \num{0.114605}, and
  \num{0.118798} on every inert arm). \texttt{ild} is omitted for width and moves
  only on the arms that fire (\num{0.4506} on C0, \num{0.5572} on A1, \num{0.6762}
  on A2, \num{0.4530} on B1, \numrange{0.4504}{0.4506} on the inert arms).
  \texttt{mood}, mus@10, H@10 higher-better; \texttt{conc}, \texttt{adj}
  lower-better. The lowest-eagerness
  cells are bolded and belong to \emph{degenerate} arms, which is what the
  relevance floor exists to catch.
  \emph{Source: \texttt{2026-07-30-eager-margin-calibrated-refutation.md}.}}
  \label{tab:eagermargin}
  {\scriptsize\setlength{\tabcolsep}{2.8pt}\begin{tabular}{llllrlllllll}
    \toprule
    Arm & run & $m$ & $\beta$ & hits & conc & adj & mood & mus@10 & H@10 & ep & verdict \\
    \midrule
    C0 & \texttt{175402-cd006} & 0 & 0 & \best{170} & 0.370588 & 0.362963 & \best{0.479846} & \best{0.455347} & \best{0.023785} & 8/14 & CONTROL \\
    A1 & \texttt{181405-53088} & 0 & 3 & 77 & 0.240717 & 0.206010 & 0.346867 & 0.399997 & 0.013339 & 18/24 & DEGENERATE \\
    A2 & \texttt{181405-0bc1a} & 0 & 12 & 20 & \best{0.133535} & \best{0.107687} & 0.221756 & 0.358132 & 0.003569 & 37/40 & DEGENERATE \\
    B1 & \texttt{181405-84da4} & 0.40 & 9 & 164 & 0.367606 & 0.364151 & 0.463440 & 0.447386 & 0.021601 & 15/21 & no effect \\
    B2 & \texttt{175402-7c304} & 0.5843 & 19 & 170 & 0.370634 & 0.363103 & 0.479955 & 0.455210 & 0.023751 & 8/14 & INERT \\
    B3 & \texttt{184556-b08c1} & 0.70 & 35 & 170 & 0.370588 & 0.362963 & 0.479846 & 0.455347 & 0.023785 & 8/14 & INERT (bit-id.) \\
    B4 & \texttt{175402-0d310} & 0.8304 & 92 & 170 & 0.370588 & 0.362963 & 0.479846 & 0.455347 & 0.023785 & 8/14 & INERT (bit-id.) \\
    B5 & \texttt{184556-47d7b} & 0.5843 & 57 & 170 & 0.370743 & 0.363103 & 0.479955 & 0.455379 & 0.023771 & 8/14 & INERT \\
    \bottomrule
  \end{tabular}}
\end{table}

\begin{table}[ht]
  \centering
  \caption{Paired statistics against the in-batch control C0 (\num{2000}
  resamples, rng \num{1337}, $n=\num{1431}$). $n_{disc}$ counts sessions whose
  hit@10 changed. Measured half-widths on the \emph{active} arms
  (\numrange{0.0080}{0.0173}) match the design's predicted \num{0.0102}, so the
  instrument was correctly powered --- the axis failed anyway.
  \emph{Source: \texttt{2026-07-30-eager-margin-calibrated-refutation.md}.}}
  \label{tab:eagermargin-paired}
  {\footnotesize\setlength{\tabcolsep}{4.5pt}\begin{tabular}{lllrrr}
    \toprule
    Arm & $\Delta$\texttt{artist\_conc} [95\% CI] & $\Delta$\texttt{music@10} [95\% CI] & $\Delta H$ & $n_{disc}$ & mean $|\Delta$conc$|$ \\
    \midrule
    A1 & $-0.12987$ [$-0.14354$, $-0.11605$] & $-0.05535$ [$-0.06606$, $-0.04465$] & $-0.010445$ & 125 & 0.191304 \\
    A2 & $-0.23705$ [$-0.25493$, $-0.22026$] & $-0.09721$ [$-0.11107$, $-0.08333$] & $-0.020215$ & 166 & 0.282895 \\
    B1 & $-0.00298$ [$-0.01085$, $+0.00519$] & $-0.00796$ [$-0.01411$, $-0.00194$] & $-0.002183$ & 62 & 0.093734 \\
    B2 & $+0.00005$ [$-0.00028$, $+0.00039$] & $-0.00014$ [$-0.00037$, $+0.00000$] & $-0.000034$ & \best{0} & 0.000419 \\
    B3 & $+0.00000$ [$+0.00000$, $+0.00000$] & $+0.00000$ [$+0.00000$, $+0.00000$] & $+0.000000$ & \best{0} & \best{0.000000} \\
    B4 & $+0.00000$ [$+0.00000$, $+0.00000$] & $+0.00000$ [$+0.00000$, $+0.00000$] & $+0.000000$ & \best{0} & \best{0.000000} \\
    B5 & $+0.00016$ [$-0.00031$, $+0.00068$] & $+0.00003$ [$-0.00036$, $+0.00045$] & $-0.000014$ & \best{0} & 0.000901 \\
    \bottomrule
  \end{tabular}}
\end{table}

\begin{figure}[ht]
  \centering
  \includegraphics[width=0.95\textwidth]{nexttrack_eager_margin}
  \caption{\textbf{Left:} the campaign's central result. The hinge's realized
  penalty at initialization, reconstructed from epoch-1 loss inflation at shared
  initialization (identical seed, so identical init and batch order). The margins
  were calibrated from the \emph{item}$\to$\emph{item} transition cosine (median
  \num{0.262}, p75 \num{0.584}, p90 \num{0.830}), but the penalty acts on
  \emph{prediction}$\to$\emph{current}, whose ceiling over training sits between
  \num{0.5843} and \num{0.70}: at $m\ge\num{0.5843}$ the penalized quantity never
  enters the hinge's support and the loss trace is bit-identical to the control at
  all \num{14} epochs, at $\beta$ up to the schema cap. \textbf{Right:} where the
  knob \emph{does} fire, eagerness and relevance are the same mechanism. Pushing
  the prediction off the current track's latent also pushes it off the region where
  the true next track lives, so there is no window between real de-eagering and
  breaching the relevance floor. The dash-dotted line marks the free on-record
  comparator that no arm here matched inside the floor.}
  \label{fig:nexttrack-eager-margin}
\end{figure}

\emph{Interpretation.} \textbf{The refutation is expressive, not dosimetric, and
that is the whole finding: the calibration measured the wrong distribution.} The
margins were drawn from $\cos(x_t, x_{t+1})$ --- how far apart two consecutively
played tracks are --- but the hinge acts on $\cos(\mathrm{pred}_t, x_t)$, how close
the model's \emph{prediction} sits to the current track. Those distributions are
not comparable. Reconstructing the hinge from epoch-1 loss inflation puts the
ceiling of the penalized quantity between \num{0.5843} and \num{0.70}, so p75 and
p90 of the real-transition distribution lie \emph{outside the support of the
quantity being penalized}, and arms B3 and B4 reproduced the control's loss to
\num{16} significant digits at every one of \num{14} epochs. The design's hedge ---
that the realized penalty would \emph{exceed} its nominal value --- is empirically
inverted by orders of magnitude. The reason is structural: an InfoNCE-trained
prediction is a \emph{hedged near-average} over many plausible continuations, so its
absolute cosine to any single item is bounded low. \textbf{Eagerness is a ranking
property --- whether the current track's neighbours occupy the top slots --- and a
hinge on an absolute cosine cannot express it.}

\textbf{Where the knob does fire, eagerness and relevance turn out to be the same
mechanism.} At $m=0$ the lever works spectacularly on the facet --- A2 cuts
concentration \num{0.371}$\to$\num{0.134} and drives \texttt{ild} to \num{0.676}
--- while \texttt{music@10} collapses to \num{0.358} and recall to \num{20} of
\num{1431}. There is no separating window: \SI{27}{\percent} of real transitions
are negative-cosine, so pushing the prediction away from the current track's latent
also pushes it out of the region where the true next track lives, and
\texttt{mood\_coh} collapses too (\num{0.480}$\to$\num{0.222}). Already at
$\beta=3$ the absolute relevance floor is breached. The pre-registered failure mode
fired exactly as written.

\textbf{The prior is sharpened rather than retired, and the usable band is
identified.} The 2026-07-22 probe at $\beta \le \num{0.2}$, $m=0$ reported ``weak''
--- correct in band, wrong in interpretation. That setting sits just \emph{under}
the destruction threshold, and by $\beta=3$ the same setting is a relevance
destroyer. The entire usable dynamic range at $m=0$ therefore lies between
$\beta=\num{0.2}$ and $\beta=\num{3}$, a factor of \num{15}, unexplored, bounded
above by degeneration and below by no effect. The dose ladder run here was
\textbf{not underpowered; it was aimed past the target}.

\textbf{Two second-order results worth keeping.} First, demoting $H$ to secondary
was load-bearing: A2 has by far the best facets and the \emph{worst} $H$
(\num{0.003569}), because the grounding factor annihilates it exactly as the
$\num{5.46}\times$ derivative ratio predicted. Had $H$ been primary this campaign
would have read as a flat null instead of a mapped frontier. Second, the cheapest
real de-eagering on record remains an \emph{architecture} setting, not this knob:
the free on-record comparator at \texttt{fusion\_layers}${}=1$ gives
$\Delta$concentration $-\num{0.051}$ at a $-\num{0.016}$ recall cost, and no arm
here matched that inside the relevance floor --- the floor-respecting arms delivered
between $-\num{0.003}$ and $+\num{0.0002}$, and the arms that beat $-\num{0.051}$
paid $\numrange{4}{8}\times$ the recall.

\paragraph{Why the approved later phases did not run.} All three were gated on a
phase-A win, and there is none. Phase C's arms were additionally defined \emph{by
reference} to the winner's margin and $\beta$, so with no winner they are
undefined; applied literally to the arms that passed the floors, the selection rule
would have spent three runs transferring a \emph{provably inert} setting to another
family. The runner stopped rather than improvise, which is the correct outcome: the
next move is a re-aimed design, not the approved phase C. Concretely, re-specify
the penalty rather than re-tune it --- hinge on a \emph{ranking} quantity such as
$\cos(\mathrm{pred}, x_t) - \cos(\mathrm{pred}, x_{t+1})$, or an explicit
same-artist mass term over the top-$k$ --- and, by elimination, make the eval-time
re-rank reach the serving path, since with the training-time lever refuted the only
lever that works was at that moment one the serving path could not reproduce. That
last observation is what the next section acts on.
\emph{Source: \texttt{2026-07-30-eager-margin-calibrated-refutation.md}.}

\section{The walk surface: measuring what the product actually serves}
\label{sec:walk}
\emph{Rationale.} Every verdict up to this point --- including the crown board
of Table~\ref{tab:lb} and all three holisticness generations of
Section~\ref{sec:power-generations} --- is computed on a \textbf{one-shot
top-10}: the model sees one held-out prefix, ranks the catalog once, and is
scored against one target. But the artifact this lab ships,
\texttt{clients/infinite-playlist}, serves an \textbf{autoregressive walk}: it
picks a track, appends it to the session, re-scores, and repeats for as many
steps as the listener asks for. Those are different surfaces, and
\texttt{domain.toml} has said so in prose the whole time --- its
\texttt{[metrics]} block states the goal as ``mood-coherent, non-eager
\emph{sessions}, not guessing the literal next track.'' The stated intent and
the implemented metric had diverged. This section closes that gap: it defines
four walk-level metrics calibrated against the user's own listening, adds three
\emph{retrieval-time} controls, and runs a matched-sample head-to-head. It is
also the one place in this document where a documented refutation is
\emph{reversed} on a new objective, which is why it is reported at length rather
than folded into the campaign sections above.

\subsection{The surface mismatch, and why drift was invisible}
\label{sec:walk-mismatch}
Two mechanical facts explain why nothing above could have caught a drifting
journey.

\textbf{First, the anchor grows.} The \texttt{mood\_coh} facet of
holisticness@10 measures each candidate's cosine to the \emph{prefix} centroid
(\texttt{seq\_common.py:558}). On a one-shot evaluation the prefix is fixed and
that is a reasonable read of ``does this list stay on vibe.'' On a walk the
prefix \emph{is} the journey so far, so the reference point travels with the
model: a journey that leaves its starting region by a thousand individually
legal steps scores well \emph{at every step}, because each step is compared to a
centroid that has already moved most of the way there. \textbf{Drift is
invisible by construction} --- not under-weighted, but unmeasurable in this
formulation.

\textbf{Second, the composite is dominated by relevance.} Differentiating the
generation-3 composite at the champion's operating point gives
$\partial H/\partial\,\texttt{music@10} = +\num{0.206}$ against
$\partial H/\partial\,\texttt{eager} = -\num{0.038}$: holisticness@10 is
$\num{5.5}\times$ more sensitive to its relevance factor than to its eagerness
factor. So even the term that was \emph{meant} to price session quality is
mostly re-reporting graded relevance. This is the same asymmetry that
Section~\ref{sec:power-generations} observed from the outside, where the
grounding factor's ${\sim}0.3$ multiplier made every generation-2 value rank
below every generation-1 value; here it is quantified at the derivative.

\subsection{Four walk-level metrics, calibrated against the user's own sessions}
\label{sec:walk-metrics}
The new instrument (\texttt{tools/walk\_eval.py}) generates journeys from
held-out \emph{real} sessions --- a 5-track seed taken from the session, then 20
autoregressive steps, prefix items excluded, at most one track per artist unless
stated --- and scores them on four axes:

\begin{description}
  \item[1. Transition match (primary).] Do the walk's step-to-step distances
  look like this user's \emph{real} listening? This follows the playlist-lab
  finding directly: measure \emph{transitions}, not averages, and calibrate
  against the user's own sessions. Averaging a distance-to-centroid is the wrong
  shape --- it penalizes a legitimate genre move and rewards standing still.
  Reported as the 1-Wasserstein distance ($W_1$, lower better) between the
  journey's consecutive-pair cosine distribution and the real one, plus two
  legible reads: the median step cosine and the share of steps inside the real
  p10--p90 band.
  \item[2. Vibe and drift.] The \emph{seed-anchored} cosine: mean cosine of the
  generated tracks to the fixed seed-core centroid (\texttt{vibe}, higher
  better), and the end-quarter minus start-quarter difference (\texttt{drift},
  nearer zero better). The anchor is fixed at the seed and does \emph{not} grow,
  which is precisely the defect of \texttt{mood\_coh} that
  Section~\ref{sec:walk-mismatch} identifies.
  \item[3. Variety as credit.] Distinct artists and distinct genres in the
  journey, counted as a \emph{credit} rather than as the cost that
  \texttt{artist\_adj}/\texttt{artist\_conc} price. Cohesion is not one genre:
  different genres are welcome if they preserve the vibe, so a journey covering
  more of them at equal vibe is better, not worse.
  \item[4. Grounding.] Graded relevance (best sonic cosine, in the balanced
  musical-distance space of Section~\ref{sec:music}) between each generated
  track and what the user \emph{actually} played next in that session. This is
  the guard-rail: without it a pleasant unrelated drift wins, which is exactly
  the lesson the crown learned when \texttt{music\_rel} was added.
\end{description}

\paragraph{The calibration, and the single number that reframed the project.}
The reference distribution is the user's own consecutive-pair cosines over the
\num{5723} \emph{train} sessions of \dsid{seq-20260715-131139}:
\textbf{\num{73632} pairs, median \num{0.261}, p10--p90
$[-\num{0.192}, \num{0.830}]$, and \SI{27.0}{\percent} of real steps
\emph{negative}}. Real listening moves --- it changes direction more than a
quarter of the time. Against that, the engine the showcase actually shipped (a
single GRU as \texttt{gru.onnx} with a seed anchor of \num{0.4}) walks with a
median step of \textbf{\num{0.661}}, roughly $\num{2.5}\times$ tighter than the
user does (Figure~\ref{fig:nexttrack-walk-transitions}). The complaint that the
showcase felt ``flat despite lots of artists'' now has a measurement behind it:
many small, safe moves. The blend champion has the opposite failure --- its
median step, \num{0.317}, sits essentially on the real median, but its
\texttt{vibe} is only \num{0.187}, because its InfoNCE projection
(Section~\ref{sec:projection}) encodes ``what follows this,'' not ``what sounds
like this,'' so a mood anchor applied \emph{in} that projected space caps at
\num{0.44} while the same anchor in the raw PCA-192 space reaches \num{0.58}.
Table~\ref{tab:walk-arms} is the full first read.

\begin{figure}[ht]
  \centering
  \includegraphics[width=0.8\textwidth]{nexttrack_walk_transitions}
  \caption{The core argument of this section, in one figure. The user's own
  \num{73632} consecutive-transition cosines (PCA-192 item space, train sessions
  only) are a \emph{broad} distribution: median \num{0.261}, p10--p90
  $[-\num{0.192}, \num{0.830}]$, \SI{27.0}{\percent} of steps negative. The
  engine the showcase shipped walks at a median step of \num{0.661} --- about
  $\num{2.5}\times$ tighter than the listener it is imitating --- and no one-shot
  metric could see it, because holisticness@10 scores a top-10 and its
  \texttt{mood\_coh} anchors to a prefix centroid that grows with the walk. The
  blend champion (\num{0.317}) matches the real stride but cannot hold a vibe.
  Bin counts transcribed from the calibration block of
  \texttt{tools/walk\_eval.py}.}
  \label{fig:nexttrack-walk-transitions}
\end{figure}

\begin{table}[ht]
  \centering
  \caption{First walk-level read of the two shipped engines. \num{40} held-out
  sessions, 5-track seed, \num{20} autoregressive steps each; \texttt{cap} is the
  per-artist limit inside the journey. $W_1$ and $|$median step $-$ \num{0.261}$|$
  lower-better; \texttt{vibe}, in-band share, artists, genres and \texttt{ground}
  higher-better; \texttt{drift} nearer zero better. The last row is the reference
  the metrics are calibrated against (its in-band share is \SI{80}{\percent} by
  construction). Best cell per column bolded among the model arms.
  \emph{Source: \texttt{tools/walk\_eval.py}.}}
  \label{tab:walk-arms}
  {\footnotesize\setlength{\tabcolsep}{4pt}\begin{tabular}{lccccccccc}
    \toprule
    Arm & $W_1$ & med.\ step & in-band & vibe & drift & artists & genres & ground \\
    \midrule
    GRU $+$anchor \num{0.4}, cap 1 (\textbf{shipped}) & 0.320 & 0.661 & \SI{80.6}{\percent} & 0.609 & $-$0.089 & 20.0 & 9.93 & 0.345 \\
    GRU $+$anchor \num{0.4}, cap 3 & 0.341 & 0.681 & \SI{76.7}{\percent} & \best{0.610} & \best{$-$0.075} & 13.2 & 6.53 & \best{0.393} \\
    GRU no anchor, cap 1 & 0.263 & 0.574 & \SI{85.2}{\percent} & 0.302 & $-$0.277 & 20.0 & 8.78 & 0.276 \\
    Champion (\Rp+M+\Cp), cap 1 & \best{0.078} & \best{0.317} & \best{\SI{90.3}{\percent}} & 0.187 & $-$0.186 & 20.0 & \best{10.33} & 0.289 \\
    Champion (\Rp+M+\Cp), cap 3 & 0.147 & 0.452 & \SI{84.8}{\percent} & 0.193 & $-$0.224 & 11.1 & 5.33 & 0.306 \\
    \midrule
    \emph{the user's real sessions} & 0.000 & 0.261 & \SI{80.0}{\percent} & --- & --- & --- & --- & --- \\
    \bottomrule
  \end{tabular}}
\end{table}

\paragraph{$W_1$ is not pairable, and that mattered.} The distributional $W_1$ is
computed over a \emph{pooled} sample of steps, so it has no per-session value and
cannot enter a paired bootstrap --- the statistic every other verdict in this
document rests on. The first two sweeps therefore compared arms measured on
\emph{different} samples (20 sessions $\times$ 20 steps for the frontier, 12
$\times$ 15 for the dual-tower probe), and a vibe gap of \num{0.588} against
\num{0.624} is well inside what that mismatch alone can produce: nothing was
established. The head-to-head (Section~\ref{sec:walk-h2h}) therefore replaces
$W_1$ with \texttt{stride\_err} $=$ $|$median step $-$ \num{0.261}$|$, a
per-session quantity that pairs, at the cost of describing only the
distribution's centre.

\subsection{Three retrieval-time controls, and two refuted ones}
\label{sec:walk-controls}
None of the three controls below touches training. All are applied at
\emph{retrieval} time, on the scores the model already produces, which is what
makes them cheap to ship and cheap to re-tune.

\begin{description}
  \item[ANCHOR ($a$).] Add $a\cdot\cos(\text{candidate},\ \text{seed centroid})$
  to every candidate's score. It holds the geist by pulling the walk back toward
  the region the listener seeded it from. For the champion the term is
  z-normalized over the admissible candidates, because its base score is itself a
  z-blend of three legs and adding a raw cosine to a z-blend is scale-mismatched.
  \item[STRIDE ($s$).] Subtract $s\cdot\cos(\text{candidate},\ \text{previous
  track})$. It \emph{enlarges} each step without letting the walk leave the
  seed's neighbourhood, because the anchor still applies --- the two terms are
  orthogonal by design, one pulling toward a fixed point and one pushing off the
  last point. It is deliberately \emph{deterministic}: sampling would also loosen
  the stride, but the showcase ships permalinks, and ``same recipe $\Rightarrow$
  same journey'' is now a promise, so a stochastic decoder was ruled out rather
  than untried.
  \item[CAP.] A \emph{hard} per-artist limit applied to the ordering
  (\texttt{seq\_common.artist\_cap\_order}), after MMR, in both
  \texttt{eval\_from\_scores} and \texttt{rank\_topk}/\texttt{predict\_ranking},
  so evaluation and serving agree by construction. Best-effort rather than
  absolute: where a prefix's admissible candidate space is exhausted the
  remaining slots are backfilled with deferred items instead of returning fewer
  than $k$. Default \texttt{None}, so every pre-existing run and model is
  byte-identical; and because it needs no sonic vectors it is the only diversity
  lever here that ports to a thin client unchanged.
\end{description}

\paragraph{Why a cap and not a penalty --- the two refuted controls.} Both soft
anti-eager levers of Section~\ref{sec:holisticness} fail for one shared reason: a
soft penalty in score space loses to a steep relevance gradient.
\texttt{eager\_beta}, the training-time regularizer, is \emph{mis-specified}
(Section~\ref{sec:eagermargin}): it hinges on the absolute value of
$\cos(\text{prediction},\ \text{current track})$, which a hedged InfoNCE prediction never drives above ${\sim}0.6$,
whereas eagerness is a \emph{ranking} property --- whether same-artist items
occupy the top slots --- not a property of an absolute similarity.
\texttt{mmr\_lambda}, the eval-time re-rank, was additionally \emph{eval-only}
until the serving path was fixed, so the $\lambda$ frontier of
Section~\ref{sec:mmr} --- genuine, confirmed, and the first crown lever on record
--- had never been a deployable lever at all; and even fixed it measures inert on the pathological
case. On the real playlist \texttt{rock chabon ahre} the model ranks \num{59} of
its top \num{60} candidates as one artist and the next artist first appears at
\textbf{rank \num{75}}, while MMR's entire diversity budget there is worth
$\num{0.3}\times\num{0.26}\approx\num{0.078}$ of normalized score --- nowhere
near enough to pay for a 70-rank relevance drop. A cap does not compete on score
at all: it \emph{cannot be outscored}, the same property \texttt{seq\_extend}'s
\texttt{artist\_cooldown} already relied on. The showcase's own soft artist
penalty turned out to be effectively \textbf{binary} in the same way --- at
\num{0} it produced artist blocks with runs up to \num{13}, and at \emph{any}
value $\ge\num{0.25}$ every 30-step walk returned \num{30} distinct artists, max
run \num{1}, concentration \num{0.000} --- so the slider could not reach the
middle at all, and a hard cap exposed as a journey shape (Explore / Balanced /
Continue / Blocks) is what reaches it.

\paragraph{The cap on the one-shot surface, for the record.}
Table~\ref{tab:walk-cap} is the cap measured where the rest of this document
measures: the canonical split's one-shot top-10, paired on bit-identical cached
scores. It cuts \texttt{artist\_conc@10} from \num{0.692} to \num{0.123} at cap
3 and raises the generation-3 holisticness@10 from \num{0.008603} to
\num{0.033862}, at a cost of \num{73} of \num{1431} exact hits. The result that
was \emph{not} predicted is that \textbf{\texttt{music@10} goes up, CI above
zero, at every cap on both models tested} ($+\num{0.011}$ to $+\num{0.023}$):
the pre-registered failure mode --- de-eagering by degrading into generic filler
--- does not occur, because the items promoted from rank 75 and beyond are
\emph{more} musically apt to the true continuation than the same-artist tracks
they displace. The cap trades \emph{literal} relevance for \emph{graded}
relevance. On the champion plus \texttt{markov\_gate} (the evidence gate of
Section~\ref{sec:markovgate}, which suppresses the Markov leg where the current
item has no bigram support) the same cap 3 gives \texttt{artist\_conc}
\num{0.598}$\to$\num{0.113} and holisticness@10
\num{0.018699}$\to$\num{0.048063}, reproduced to every digit by three
registered confirming runs (\runid{run-20260730-235524-0ec1d-seq-blend} at cap
off, \runid{run-20260730-235524-93d38-seq-blend} at cap 3,
\runid{run-20260730-234635-a6252-seq-blend} at cap 5), and it \textbf{hardens on
the disjoint 0.64-cold split} \dsid{seq-20260718-211238} on two pre-registered
legs: $\Delta$\texttt{artist\_conc} $-\num{0.52699}$
$[-\num{0.54546}, -\num{0.50825}]$ CI below zero \emph{and}
$\Delta$\texttt{music@10} $+\num{0.02644}$
$[+\num{0.01919}, +\num{0.03368}]$ CI above zero, with the relevance gain
\emph{larger} on the colder window. Two cautions travel with it: exact recall is
the only regressing metric at every cap, and the canonical split's absolute
\texttt{music@10} $\ge\num{0.43}$ floor does \textbf{not} transfer --- the
second split's own control sits at \num{0.4295} on a legitimately colder window,
so floors must be derived per baseline and pre-registration must be on the
paired $\Delta$.

\begin{table}[ht]
  \centering
  \caption{The hard per-artist cap on the one-shot surface, on the deployed
  champion (\Rp+M+\Cp, ungated), canonical split \dsid{seq-20260715-131139},
  $n=\num{1431}$. Off-run measurement: scores are computed once per model and
  memoized, so every cap arm sees bit-identical score vectors and the pairing is
  exact rather than merely matched (\num{2000} resamples, rng \num{1337}).
  holisticness@10 is generation 3 (Section~\ref{sec:power-generations}).
  \texttt{artist\_conc}, \texttt{artist\_adj} lower-better; the rest
  higher-better. \emph{Source:
  \texttt{2026-07-30c-artist-cap-measurement.md}.}}
  \label{tab:walk-cap}
  {\footnotesize\setlength{\tabcolsep}{4.2pt}\begin{tabular}{lcccccccc}
    \toprule
    cap & R@10 & hits & MRR & artist\_conc & artist\_adj & ild & music@10 & H@10 (gen-3) \\
    \midrule
    off & \best{0.211740} & \best{303} & \best{0.121985} & 0.691653 & 0.609713 & 0.406374 & 0.453790 & 0.008603 \\
    5   & 0.178896 & 256 & 0.119693 & 0.252613 & 0.337247 & 0.563417 & 0.474291 & 0.027258 \\
    3   & 0.160727 & 230 & 0.116915 & 0.122975 & 0.215304 & 0.613641 & \best{0.476787} & 0.033862 \\
    2   & 0.138365 & 198 & 0.112313 & \best{0.061946} & \best{0.150384} & \best{0.639717} & 0.474317 & \best{0.036331} \\
    \bottomrule
  \end{tabular}}
\end{table}

\begin{figure}[ht]
  \centering
  \includegraphics[width=0.95\textwidth]{nexttrack_artist_cap}
  \caption{The cap, across both models measured.
  \textbf{Left:} artist concentration collapses,
  \num{0.692}$\to$\num{0.113} at cap 3 on the gated champion.
  \textbf{Middle:} the result that was \emph{not} predicted --- graded relevance
  \emph{rises}, with the paired CI above zero at every cap on both models, and it
  replicates on the disjoint \SI{64}{\percent}-cold split at $+\num{0.026}$. The
  pre-registered worry was that a cap would de-eager by degrading into generic
  filler; the opposite happens, because the items promoted from rank \num{75} and
  beyond are \emph{more} musically apt to the true continuation than the
  same-artist tracks they displace. \textbf{Right:} the price, on the recall
  $\times$ crown plane. The cap is \emph{priced}, not free: gate $+$ cap 3 costs
  \num{52} of \num{1431} exact hits against the deployed status quo, while landing
  a generation-3 holisticness@10 about \SI{40}{\percent} above the incumbent crown
  at $\num{1.5}\times$ its recall. Every arm clears the \num{155}-hit recall
  floor, so none is an instance of the ``maximize holisticness by predicting
  worse'' pathology.}
  \label{fig:nexttrack-artist-cap}
\end{figure}

\paragraph{The anchor $\times$ stride frontier.} With the two score-space
controls in hand the question becomes whether the \emph{conjunction} is
reachable at all: a real-listening stride \emph{and} a held vibe.
Table~\ref{tab:walk-frontier} sweeps \num{24} cells --- \num{16} on the GRU,
\num{8} on the champion, all at cap 1 --- against a target window
pre-registered before the sweep: median step inside the real band's middle
($\numrange{0.20}{0.38}$) \emph{and} \texttt{vibe} $\ge\num{0.50}$. Exactly two
cells land inside it, one per engine (Figure~\ref{fig:nexttrack-walk-frontier}).
The structure of the grid is the finding: the two controls are genuinely
orthogonal --- along a fixed anchor, raising the stride walks the median step
down monotonically while the vibe decays much more slowly --- so the window is
reachable only with both on, and neither shipped configuration was inside it.
The champion's anchor also has to be applied in the \emph{raw} PCA-192 space
rather than its own projected geometry, which is why the frontier reaches
\num{0.503} where the first probe capped at \num{0.44}.

\begin{table}[ht]
  \centering
  \caption{The anchor ($a$) $\times$ stride ($s$) frontier, \num{20} held-out
  sessions $\times$ \num{20} steps, cap 1, both engines. $W_1$ lower-better;
  ``med.\ step'' is compared against the user's real \num{0.261}; \texttt{vibe},
  genres, \texttt{ground} higher-better; \texttt{drift} nearer zero better. The
  window is median step $\in[\num{0.20},\num{0.38}]$ \emph{and} \texttt{vibe}
  $\ge\num{0.50}$, pre-registered before the sweep; the two rows that satisfy it
  are bolded. The shipped GRU configuration is $a=\num{0.4}$, $s=\num{0}$; the
  champion as deployed is $a=\num{0}$, $s=\num{0}$.
  \emph{Source: \texttt{tools/walk\_frontier.py}.}}
  \label{tab:walk-frontier}
  {\footnotesize\setlength{\tabcolsep}{5pt}\begin{tabular}{llcccccc}
    \toprule
    Engine & $a$ / $s$ & $W_1$ & med.\ step & vibe & drift & genres & ground \\
    \midrule
    GRU & 0.0 / 0.0 & 0.287 & 0.607 & 0.286 & $-$0.328 & 9.70 & 0.270 \\
    GRU & 0.0 / 0.3 & 0.273 & $-$0.018 & 0.200 & $-$0.157 & 8.95 & 0.291 \\
    GRU & 0.0 / 0.6 & 0.564 & $-$0.374 & 0.123 & $-$0.102 & 12.00 & 0.244 \\
    GRU & 0.0 / 1.0 & 0.702 & $-$0.557 & 0.084 & $-$0.036 & 15.80 & 0.224 \\
    GRU & 0.2 / 0.0 & 0.323 & 0.663 & 0.563 & $-$0.124 & 9.40 & 0.345 \\
    GRU & 0.2 / 0.3 & 0.186 & 0.094 & 0.426 & $-$0.098 & 11.05 & 0.335 \\
    GRU & 0.2 / 0.6 & 0.473 & $-$0.251 & 0.330 & $-$0.067 & 13.25 & 0.303 \\
    GRU & 0.2 / 1.0 & 0.675 & $-$0.522 & 0.185 & $-$0.077 & 15.05 & 0.253 \\
    GRU & 0.4 / 0.0 \emph{(shipped)} & 0.361 & 0.709 & 0.639 & $-$0.090 & 9.95 & 0.349 \\
    \best{GRU} & \best{0.4 / 0.3} & \best{0.104} & \best{0.297} & \best{0.588} & \best{$-$0.076} & \best{11.05} & \best{0.358} \\
    GRU & 0.4 / 0.6 & 0.356 & $-$0.127 & 0.437 & $-$0.044 & 12.55 & 0.303 \\
    GRU & 0.4 / 1.0 & 0.588 & $-$0.405 & 0.302 & $-$0.027 & 15.20 & 0.287 \\
    GRU & 0.8 / 0.0 & 0.364 & 0.713 & 0.678 & $-$0.067 & 11.00 & 0.352 \\
    GRU & 0.8 / 0.3 & 0.190 & 0.506 & 0.672 & $-$0.070 & 12.15 & 0.351 \\
    GRU & 0.8 / 0.6 & 0.122 & 0.176 & 0.605 & $-$0.038 & 12.60 & 0.358 \\
    GRU & 0.8 / 1.0 & 0.420 & $-$0.188 & 0.460 & $-$0.016 & 14.65 & 0.335 \\
    \midrule
    Champion & 0.0 / 0.0 \emph{(deployed)} & 0.062 & 0.318 & 0.185 & $-$0.203 & 11.00 & 0.275 \\
    Champion & 0.0 / 0.5 & 0.240 & $-$0.032 & 0.104 & $+$0.002 & 12.15 & 0.266 \\
    Champion & 1.0 / 0.0 & 0.239 & 0.577 & 0.551 & $-$0.138 & 8.95 & 0.327 \\
    \best{Champion} & \best{1.0 / 0.5} & \best{0.072} & \best{0.282} & \best{0.503} & \best{$-$0.130} & \best{9.55} & \best{0.324} \\
    Champion & 2.0 / 0.0 & 0.276 & 0.596 & 0.625 & $-$0.060 & 8.90 & 0.359 \\
    Champion & 2.0 / 0.5 & 0.170 & 0.472 & 0.615 & $-$0.105 & 9.45 & 0.338 \\
    Champion & 4.0 / 0.0 & 0.300 & 0.645 & 0.669 & $-$0.051 & 9.85 & 0.338 \\
    Champion & 4.0 / 0.5 & 0.236 & 0.544 & 0.664 & $-$0.085 & 9.70 & 0.340 \\
    \bottomrule
  \end{tabular}}
\end{table}

\begin{figure}[ht]
  \centering
  \includegraphics[width=0.8\textwidth]{nexttrack_walk_frontier}
  \caption{The two retrieval-time controls are orthogonal, and only their
  \emph{conjunction} reaches the target window (shaded: median step
  $\in[\num{0.20},\num{0.38}]$ and \texttt{vibe} $\ge\num{0.50}$). Each line
  joins a fixed anchor as the stride grows, which moves the journey leftward
  (larger steps) while the vibe decays slowly --- so the trade-off is real but
  gentle. Neither shipped configuration is inside the window: the GRU at
  $a=\num{0.4}$, $s=\num{0}$ holds a vibe of \num{0.639} but strides
  \num{0.709}, and the champion as deployed strides \num{0.318} at a vibe of
  \num{0.185}. \num{20} held-out sessions $\times$ \num{20} steps, cap 1
  (\texttt{tools/walk\_frontier.py}).}
  \label{fig:nexttrack-walk-frontier}
\end{figure}

\subsection{The matched-sample head-to-head}
\label{sec:walk-h2h}
\emph{Rationale.} The two sweeps above each found a cell inside the window, but
on different samples, and $W_1$ cannot be paired --- so the finalists had to be
re-run against one another on \emph{identical} material before anything could be
claimed. \texttt{tools/walk\_headtohead.py} walks five arms over the same
\num{40} held-out sessions from the same 5-track seeds, at cap 1, and reports
paired bootstrap differences (\num{2000} resamples, rng \num{1337}) against the
\textbf{shipped} configuration --- the single GRU at anchor \num{0.4}, stride
\num{0} --- on five per-session metrics: \texttt{vibe} (higher better),
\texttt{stride\_err} (lower better), \texttt{drift} (nearer zero),
\texttt{genres} (higher) and \texttt{ground} (higher).

\emph{Results.} Table~\ref{tab:walk-means} gives the means and
Table~\ref{tab:walk-paired} the paired differences;
Figure~\ref{fig:nexttrack-walk-headtohead} plots the latter one panel per
metric, each on its own scale.

\begin{table}[ht]
  \centering
  \caption{Walk-metric means over the \num{40} identical held-out sessions
  ($\num{20}$ steps, cap 1). \texttt{stride\_err} $=$
  $|$median step $-$ \num{0.261}$|$ (lower better); \texttt{vibe},
  \texttt{genres}, \texttt{ground} higher-better; \texttt{drift} nearer zero
  better. Best cell per column bolded.
  \emph{Source: \texttt{tools/walk\_headtohead.py}.}}
  \label{tab:walk-means}
  {\footnotesize\setlength{\tabcolsep}{6pt}\begin{tabular}{lccccc}
    \toprule
    Arm & vibe & stride\_err & drift & genres & ground \\
    \midrule
    GRU shipped, $a$\num{0.4} $s$\num{0.0} & 0.609 & 0.399 & $-$0.089 & 9.925 & 0.345 \\
    GRU tuned, $a$\num{0.4} $s$\num{0.3} & 0.568 & \best{0.129} & $-$0.080 & 10.450 & \best{0.357} \\
    dual latent/latent f0, $a$\num{0.8} $s$\num{0.4} & \best{0.638} & 0.157 & $-$0.063 & 10.975 & 0.351 \\
    \best{dual latent/cummean f0, $a$\num{0.8} $s$\num{0.5}} & 0.610 & 0.144 & \best{$-$0.053} & \best{11.425} & 0.351 \\
    Champion (\Rp+M+\Cp), $a$\num{1.0} $s$\num{0.5} & 0.476 & 0.135 & $-$0.131 & 9.600 & 0.334 \\
    \bottomrule
  \end{tabular}}
\end{table}

\begin{table}[ht]
  \centering
  \caption{Paired per-session differences against the \textbf{shipped} GRU
  ($a$\num{0.4} $s$\num{0.0}) over the same \num{40} sessions, \num{2000}
  bootstrap resamples, rng \num{1337}. ``Verdict'' reads the CI relative to the
  metric's better direction: \texttt{stride\_err} is lower-better, so a CI below
  zero is a \emph{win}; \texttt{drift} starts negative for every arm, so a
  positive $\Delta$ moves it toward zero. Wins bolded.
  \emph{Source: \texttt{tools/walk\_headtohead.py}.}}
  \label{tab:walk-paired}
  {\footnotesize\setlength{\tabcolsep}{6pt}\begin{tabular}{llcl}
    \toprule
    Arm & Metric & $\Delta$ [95\% CI] & Verdict \\
    \midrule
    GRU tuned $a$\num{0.4} $s$\num{0.3} & vibe & $-$0.0409 [$-$0.0602, $-$0.0221] & CI$<$0: \emph{costs vibe} \\
     & stride\_err & \best{$-$0.2695 [$-$0.3058, $-$0.2358]} & \best{CI$<$0: win} \\
     & drift & $+$0.0089 [$-$0.0255, $+$0.0420] & straddles \\
     & genres & $+$0.525 [$-$0.475, $+$1.500] & straddles \\
     & ground & $+$0.0116 [$-$0.0143, $+$0.0408] & straddles \\
    \midrule
    dual latent/latent f0 & vibe & \best{$+$0.0297 [$+$0.0146, $+$0.0468]} & \best{CI$>$0: win} \\
    $a$\num{0.8} $s$\num{0.4} & stride\_err & \best{$-$0.2423 [$-$0.2790, $-$0.2053]} & \best{CI$<$0: win} \\
     & drift & $+$0.0259 [$-$0.0011, $+$0.0521] & straddles \\
     & genres & \best{$+$1.050 [$+$0.250, $+$1.876]} & \best{CI$>$0: win} \\
     & ground & $+$0.0060 [$-$0.0172, $+$0.0314] & straddles \\
    \midrule
    \best{dual latent/cummean f0} & vibe & $+$0.0013 [$-$0.0151, $+$0.0180] & straddles: \emph{held} \\
    \best{$a$\num{0.8} $s$\num{0.5}} & stride\_err & \best{$-$0.2545 [$-$0.2962, $-$0.2127]} & \best{CI$<$0: win} \\
     & drift & \best{$+$0.0351 [$+$0.0071, $+$0.0627]} & \best{CI$>$0: win} \\
     & genres & \best{$+$1.500 [$+$0.500, $+$2.500]} & \best{CI$>$0: win} \\
     & ground & $+$0.0059 [$-$0.0163, $+$0.0285] & straddles \\
    \midrule
    Champion (\Rp+M+\Cp) & vibe & $-$0.1327 [$-$0.1629, $-$0.1041] & CI$<$0: \emph{costs vibe} \\
    $a$\num{1.0} $s$\num{0.5} & stride\_err & \best{$-$0.2638 [$-$0.3120, $-$0.2155]} & \best{CI$<$0: win} \\
     & drift & $-$0.0422 [$-$0.0983, $+$0.0093] & straddles \\
     & genres & $-$0.325 [$-$1.500, $+$0.876] & straddles \\
     & ground & $-$0.0115 [$-$0.0467, $+$0.0236] & straddles \\
    \bottomrule
  \end{tabular}}
\end{table}

\begin{figure}[ht]
  \centering
  \includegraphics[width=0.95\textwidth]{nexttrack_walk_headtohead}
  \caption{The matched-sample head-to-head, one panel per metric on its own
  scale --- deliberately not a dual axis, which would imply a comparability
  these five quantities do not have. Green is a CI entirely on the metric's
  better side, red entirely on the worse side, grey straddles zero. Every arm
  buys the real stride; only the dual-tower arms buy it without paying vibe, and
  \texttt{ground} straddles for \emph{every} arm.}
  \label{fig:nexttrack-walk-headtohead}
\end{figure}

\emph{Interpretation.} Read down the \texttt{stride\_err} column first: all four
tuned arms cut it by ${\sim}\num{0.25}$ with the CI below zero, which is simply
the stride control working --- the shipped engine's \num{0.399} error was
${\sim}\num{2.6}\times$ the next worst arm's. The interesting variation is what
each arm \emph{pays} for it. Pushing the single GRU to a real stride
\textbf{costs vibe} ($-\num{0.0409}$, CI below zero) and the champion pays far
more ($-\num{0.1327}$, CI below zero): for these two engines stride and vibe are
in tension, exactly as the frontier's gentle-but-real slope predicted. Both
dual-tower arms escape the trade-off. \texttt{latent/latent} at $a\num{0.8}$
$s\num{0.4}$ \emph{gains} vibe ($+\num{0.0297}$, CI above zero) and genres
($+\num{1.05}$, CI above zero) while matching the stride;
\texttt{latent/cummean} at $a\num{0.8}$ $s\num{0.5}$ holds vibe
($+\num{0.0013}$, straddling --- held, not traded) and is the only arm with
\emph{three} wins, adding \num{1.5} genres (CI above zero) and the only
significant drift improvement ($+\num{0.0351}$, CI above zero, i.e.\ the journey
ends closer to where it started). On the objective as stated ---
mood-holding, genre-exploring sequences --- \texttt{latent/cummean} is the
winner, and the single GRU that was actually shipped is not on the frontier at
all.

\subsection{The dual-tower reversal, and what the record already knew}
\label{sec:walk-reversal}
This is worth stating plainly, because it is the most consequential
methodological result in this document: \textbf{the dual-tower family was closed
by three campaigns on recall grounds and is the winner on the stated
objective.} Section~\ref{sec:dualtower} refuted it across the complete
$\{$\texttt{latent},\texttt{delta},\texttt{cummean}$\}^2$ view matrix; the
\texttt{cummean/cummean} arm was even flagged \emph{degenerate} for topping the
generation-1 holisticness at Recall@10 \num{0.066}. Every one of those judgements
was made on one-shot next-track accuracy, and none of them is wrong on that
surface --- the reversal is not a correction of the recall verdicts but a
demonstration that they answered a different question. Section~\ref{sec:power-stats}'s
power analysis makes the two readings compatible rather than contradictory: what
the dual-tower campaign established was the absence of a \emph{large} recall
gain, which is fully consistent with the family being better at something the
recall metric does not measure.

\textbf{The corroboration was already in the record, filed as a negative.} The
interpretability note at \texttt{predictors/seq\_dualgru.py:113} records that the
\texttt{cummean} tower ``added real linear directions (CKA \num{0.458}) but
\textbf{zero} next-item concepts tower A lacked.'' Read as a next-track result
that is dead weight, and it was written down as such. Read on the walk surface it
is the useful part: a tower whose content is orthogonal to next-item prediction
is exactly a tower carrying \emph{session state} rather than adjacency. Two
structural facts complete the explanation. The family retrieves in the
\emph{raw} item-latent space with no learned projection, so a raw-space mood
anchor is native to its geometry --- unlike the champion, whose InfoNCE
projection strips precisely the sound-similarity structure an anchor needs. And
one promoted arm reads the \emph{causal mean} of the prefix, which is
structurally a running-average signal.

\textbf{One sub-hypothesis is refuted.} The natural prediction from the above ---
that the \texttt{cummean} tower holds a vibe \emph{natively} --- is false: with
both controls off its vibe is \num{0.236}, statistically indistinguishable from
\texttt{latent/latent}'s \num{0.234}. The architecture does not supply the mood
hold on its own; it only becomes the best host for one once the anchor and stride
are applied. That is a weaker and more interesting claim than the one the
structural argument suggested, and it is the one the measurement supports.

\subsection{Deployment, and what is not established}
\label{sec:walk-limits}
The tuned dual-tower arm is now the \textbf{default engine} in the showcase,
exported as \texttt{dualgru.onnx} (\num{3.16}\,MB) with ONNX-versus-torch
parity of \num{1e-7} across every prefix length $T = 1\ldots40$ --- the
\texttt{cummean} cumulative sum was the trace risk, and it is the reason the
parity check swept $T$ rather than spot-checking it. The single GRU is kept
alongside for comparison, and the tuned anchor/stride live in the deployment
manifest rather than in code, so re-tuning is a re-bake and not a code edit. The
blend champion was removed from the showcase end-to-end: served from a Worker
its walk cost ${\sim}\SI{0.94}{\second}$ of CPU per batch of 8, roughly
$100\times$ the \SI{10}{\milli\second} CPU cap of the free plan the showcase is
deployed on, and batching cannot rescue a single step that already costs
${\sim}\SI{117}{\milli\second}$.

\paragraph{What this section does \emph{not} establish.} Four limits, stated
plainly. \textbf{(1)} \texttt{ground} straddles zero for \emph{every} arm
(Table~\ref{tab:walk-paired}), so no arm is established as better at predicting
the real continuation --- the wins are about journey \emph{shape} only, and the
grounding guard-rail confirms only that no arm degraded into pleasant
irrelevance. \textbf{(2)} One split, one seed, \num{40} sessions of \num{20}
steps; the walk metrics have had no second-split hardening of the kind
Section~\ref{sec:projection-hardening} applies to the recall crown, and the
target window itself was chosen from the first sweep's spread. \textbf{(3)}
Nothing here has been \emph{listened} to. Every claim is a geometric claim about
cosines in a PCA-192 space, and ``holds a mood'' is operationalized as
``stays near a seed centroid,'' which is a hypothesis about perception, not a
measurement of it. \textbf{(4)} The project's primary metric is
\textbf{unchanged}: holisticness@10 remains \texttt{[metrics].primary} and this
document's board remains denominated in Recall@10
(Section~\ref{sec:power-generations}). Whether the walk metrics should become
primary is an open decision, not a settled one --- and until it is settled, the
honest summary is that the walk surface is now \emph{measurable}, and on it one
refuted family wins.
\emph{Sources: \texttt{2026-07-30c-artist-cap-measurement.md},
\texttt{PROJECT-FACTS.md} (roll-up 2026-07-31b),
\texttt{tools/walk\_eval.py}, \texttt{tools/walk\_frontier.py},
\texttt{tools/walk\_dualgru.py}, \texttt{tools/walk\_headtohead.py}.}

\section{Stage-stack topology: depth refuted on both surfaces}
\label{sec:stagestack}
\emph{Rationale.} The user's question was direct: does an explicit \emph{stack} of
per-step MLP and recurrent stages beat a single GRU, on a corpus where five
consecutive architecture campaigns have returned nulls? Four arrangements were
specified --- \texttt{ann+gru+ann+gru}, \texttt{ann+gru+gru+ann},
\texttt{ann+gru+gru+gru}, \texttt{gru+gru+gru+gru} --- plus a fifth,
\texttt{gru+ann+gru+gru}, chosen by the user to test \emph{why} a leading MLP hurt
at $N=1$: Section~\ref{sec:preencoder} attributed that loss to half-wave
rectification destroying signed prefix information that a GRU's gated linear input
transform preserves, so putting the MLP on a \emph{gated hidden state} instead of
raw signed PCA should make the harm vanish. This is also the first campaign in this
arc to be judged on \emph{two} surfaces by design rather than in retrospect, and it
sets the standard the campaign after it inherits.

\paragraph{Design, and the equivalence proof that made a control possible.} Seven
runs of a new predictor, \texttt{seq-stack}, on the canonical split
\dsid{seq-20260715-131139} (\num{1431} test sessions), \SI{90}{\minute} of compute,
zero failures, zero contingency runs spent. Capacity was equalised at
${\sim}\num{1282000}$ parameters across the five 4-stage arms (mutual spread
\SI{0.445}{\percent}), and all seven counts were re-verified by instantiating the
model rather than trusted from the design. Everything but the stage string and the
widths is held: dropout \num{0.1}, InfoNCE, $\tau$~\num{0.07}, \num{40} epochs,
patience \num{6}, lr \num{0.001}, batch \num{128}, seed \num{1337}, MMR off, artist
cap off.

Bit-identity with the family's historical \num{176} hits is \emph{not available}
here and its absence is not a failure: the trainer re-seeds on entry, so handing it
a pre-built model offsets the dropout stream even when the initial weights match.
The campaign therefore gated on a \textbf{functional-equivalence proof} plus an
\textbf{in-batch control}, and every verdict is stated against that control. The
proof was run \emph{before} the registry edit and the server restart, so a failure
would have cost nothing: at \texttt{stages=["gru"]} the new model computes the same
function as the incumbent single GRU, with $\max|\Delta|$ \textbf{\texttt{0.00e+00}}
on both \texttt{forward} and \texttt{predict\_next} at prefix lengths
$T=1,2,5,17,40$ and \num{394944} parameters on both sides. The in-batch control
then returned \textbf{\num{178} of \num{1431}}, inside the pre-registered band
$[\num{155},\num{198}]$ and only \num{2} hits off the historical \num{176} ---
exactly the small offset the re-seeding argument predicts. Two further gates
mattered: metric completeness was green on all seven runs (no worker-image key
stripping of the kind Section~\ref{sec:mmr} documents), and the walk instrument
was verified against the shipped engine, whose $\Delta$\texttt{vibe}
($-\num{0.0033}$) and $\Delta$\texttt{stride\_err} ($+\num{0.0137}$) both straddle
zero against the freshly-trained control. The walk harness additionally reproduced
\emph{all ten} previously recorded means exactly and stayed bit-stable as arms were
added, which is what licenses pairing across separate invocations.

\emph{Results, the one-shot surface.} Table~\ref{tab:stagestack-s1}.

\begin{table}[ht]
  \centering
  \caption{The stage-stack batch on the one-shot surface, seven runs on
  \dsid{seq-20260715-131139} ($n_{\mathrm{test}}=\num{1431}$), sorted by
  Recall@10. Run ids are \protect\runid{run-20260801-} plus the chunk time and hex
  suffix, in full: \texttt{003536-eaf81} (C0), \texttt{003536-bcedf} (C1),
  \texttt{003646-07994} (B), \texttt{003646-8cc2d} (E), \texttt{003646-dbe68} (C),
  \texttt{003536-d612e} (D), \texttt{003646-6e3bf} (A). Stage codes: \texttt{a}
  $=$ per-step MLP, \texttt{g} $=$ GRU. MRR is truncated at \num{10} and
  recomputed from the stored predictions (the stored MRR is over the full ranking
  and is therefore not pairable). Every H@10 was recomputed from
  \texttt{predictions.json} under the live definition and matches the stored value
  exactly. The paired $\Delta$ is against the in-batch control C0 (\num{2000}
  resamples, rng \num{1337}). Truncated MRR@10 and \texttt{mood\_coh} are omitted
  for width and track recall closely: MRR@10 \num{0.04555} (C0), \num{0.04579}
  (C1), then \numrange{0.02080}{0.02894} across the five 4-stage arms;
  \texttt{mood\_coh} \num{0.4788} (C0), \num{0.4856} (C1), then
  \numrange{0.3821}{0.3964}. R@10, H@10, mus@10
  higher-better; \texttt{conc}, \texttt{adj} lower-better --- and the
  lowest-eagerness cells belong to \emph{refuted} arms, which is what the recall
  floor is for. \texttt{ild} rises monotonically as the arms get worse
  (\num{0.4549} on C0, \num{0.4424} on C1, then \numrange{0.5112}{0.5373} across
  the five 4-stage arms) and is omitted here for width.
  \emph{Source: \texttt{2026-08-01-stage-stack-topology.md}.}}
  \label{tab:stagestack-s1}
  {\scriptsize\setlength{\tabcolsep}{2.0pt}\begin{tabular}{lllrllllll}
    \toprule
    Arm & id & stages / width & params & R@10 (hits) & $\Delta$R@10 [95\% CI] & H@10 & conc & adj & mus@10 \\
    \midrule
    C0 \emph{ctrl} & \texttt{eaf81} & \texttt{gru} h256 & \num{394944} & \best{0.12439 (178)} & --- & \best{0.024378} & 0.3589 & 0.3482 & \best{0.4552} \\
    C1 \emph{cap.} & \texttt{bcedf} & \texttt{gru} h537 & \num{1280937} & 0.12159 (174) & $-0.00280$ [$-.0119$, $+.0056$] & 0.023418 & 0.3773 & 0.3653 & 0.4556 \\
    B & \texttt{07994} & \texttt{agga} h299 & \num{1281407} & 0.08595 (123) & $-0.03843$ [$-.0538$, $-.0231$] & 0.017558 & 0.3176 & 0.2639 & 0.4264 \\
    E & \texttt{8cc2d} & \texttt{gagg} h259 & \num{1277321} & 0.08246 (118) & $-0.04193$ [$-.0566$, $-.0280$] & 0.021053 & 0.2704 & 0.2376 & 0.4310 \\
    C & \texttt{dbe68} & \texttt{aggg} h256 & \num{1283008} & 0.07966 (114) & $-0.04472$ [$-.0601$, $-.0301$] & 0.021442 & \best{0.2626} & 0.2375 & 0.4324 \\
    D & \texttt{d612e} & \texttt{gggg} h229 & \num{1282821} & 0.07687 (110) & $-0.04752$ [$-.0622$, $-.0328$] & 0.018529 & 0.2644 & 0.2326 & 0.4222 \\
    A & \texttt{6e3bf} & \texttt{agag} h299 & \num{1281407} & 0.07058 (101) & $-0.05381$ [$-.0685$, $-.0391$] & 0.017374 & 0.2730 & \best{0.2335} & 0.4168 \\
    \bottomrule
  \end{tabular}}
\end{table}

No arm meets any win clause. Every 4-stage arm loses recall \emph{and}
holisticness by paired CI, and all five sit \textbf{below the mandatory
\num{0.108} absolute recall floor} (\numrange{0.0706}{0.0860}). Derived
per-baseline from the control's own \texttt{music@10} of \num{0.45519}, every
4-stage arm is \SIrange{5.0}{8.4}{\percent} below the relevance floor too, so none
is even relevance-neutral. The losses are $\numrange{2.5}{3.7}\times$ the measured
half-width: this is not a noise-band question.

\emph{Results, the walk surface.} Table~\ref{tab:stagestack-s2}, at the matched
primary cell $a\num{0.4}$/$s\num{0.3}$, paired against the in-batch control on the
walk metrics of Section~\ref{sec:walk-metrics}. \texttt{vibe} is always reported with
\texttt{stride\_err} beside it, at matched stride, per the standing caveat that a
vibe gain bought by taking shorter steps is not a gain.

\begin{table}[ht]
  \centering
  \caption{The stage-stack batch on the generated walk: \num{40} held-out sessions
  $\times$ \num{20} steps at artist cap 1, paired against the \emph{in-batch}
  control C0 (\num{2000} resamples). The refuted clause is
  ``$\Delta$\texttt{vibe} CI below zero \emph{and} $\Delta$\texttt{stride\_err}
  not CI below zero''. \texttt{vibe}, \texttt{genres}, \texttt{ground}
  higher-better; \texttt{stride\_err} lower-better; \texttt{drift} nearer zero
  better. CI bounds are quoted where the interval clears zero.
  \emph{Source: \texttt{2026-08-01-stage-stack-topology.md}.}}
  \label{tab:stagestack-s2}
  {\footnotesize\setlength{\tabcolsep}{2.6pt}\begin{tabular}{lccccclc}
    \toprule
    Arm & vibe & stride\_err & drift & genres & ground & $\Delta$vibe [95\% CI] & V2 \\
    \midrule
    GRU shipped $a$0.4 $s$0.0 & \best{0.609} & 0.399 & $-0.089$ & 9.925 & 0.345 & $+0.0376$ CI$>$0 & reference \\
    shipped engine $a$0.4 $s$0.3 & 0.568 & 0.129 & $-0.080$ & \best{10.450} & \best{0.357} & $-0.0033$ straddles & instrument gate \\
    \textbf{C0} $a$0.4 $s$0.3 & 0.571 & 0.116 & $-0.080$ & 10.325 & 0.344 & base & base \\
    C1 \texttt{gru} h537 & 0.561 & 0.122 & $-0.082$ & 10.625 & 0.359 & $-0.0100$ straddles & indistinguishable \\
    D \texttt{gggg} & 0.557 & 0.122 & $-0.071$ & 10.225 & 0.319 & $-0.0140$ [$-.0255$, $-.0022$] & REFUTED $\times2$ \\
    C \texttt{aggg} & 0.548 & \best{0.113} & $-0.064$ & 10.250 & 0.340 & $-0.0231$ [$-.0381$, $-.0084$] & REFUTED \\
    B \texttt{agga} & 0.545 & 0.126 & \best{$-0.061$} & 9.675 & 0.335 & $-0.0258$ [$-.0421$, $-.0108$] & REFUTED \\
    E \texttt{gagg} & 0.540 & 0.109 & $-0.084$ & 9.225 & 0.335 & $-0.0316$ [$-.0494$, $-.0144$] & REFUTED \\
    A \texttt{agag} & 0.534 & 0.117 & $-0.085$ & 9.125 & 0.336 & $-0.0375$ [$-.0559$, $-.0214$] & REFUTED \\
    \bottomrule
  \end{tabular}}
\end{table}

Every 4-stage arm fires the refuted clause. D fires it \emph{twice}, also failing
the grounding floor ($\Delta$\texttt{ground} $-\num{0.0250}$ [$-\num{0.0474}$,
$-\num{0.0048}$]), and A additionally loses genre variety ($-\num{1.200}$
[$-\num{2.125}$, $-\num{0.300}$]). No arm approaches a win clause. Two robustness
checks close the escape routes. A pre-registered \textbf{stride probe} over
$s\in\{\num{0.2},\num{0.3},\num{0.4}\}$ was run for the only two arms with any
directional movement toward a walk win, and \emph{both peak at $s\num{0.3}$ --- the
cell they were already measured in} (C: vibe/stride\_err \num{0.586}/\num{0.183},
\num{0.548}/\num{0.113}, \num{0.517}/\num{0.153}; E: \num{0.576}/\num{0.155},
\num{0.540}/\num{0.109}, \num{0.516}/\num{0.174}); at $s\num{0.4}$ both lose vibe
with the CI below zero while \texttt{stride\_err} gets \emph{worse}. And the single
best arm was re-measured at the \textbf{deployed} cell $a\num{0.8}$/$s\num{0.5}$,
where it is refuted again ($\Delta$\texttt{vibe} $-\num{0.0158}$
[$-\num{0.0267}$, $-\num{0.0059}$], $\Delta$\texttt{stride\_err} straddling).

\begin{figure}[ht]
  \centering
  \includegraphics[width=0.95\textwidth]{nexttrack_stage_stack}
  \caption{\textbf{Left:} all five 4-stage arms fall \emph{through} the mandatory
  absolute recall floor, while the $\num{3.24}\times$ capacity control ties the
  anchor --- so the damage is depth, not capacity. Error bars are the paired
  bootstrap CI against the in-batch control, mapped onto the hit scale.
  \textbf{Middle:} the campaign's most interesting number. Arm B reaches the
  \emph{best} validation InfoNCE in the batch and retrieves \num{55} fewer hits
  than the control: the objective the whole family trains on and top-10 retrieval
  quality come apart at depth, so any architecture search selecting on validation
  loss would pick the wrong arm. \textbf{Right:} the two-surface closure. The walk
  read at the matched primary cell refutes the same five arms that the one-shot
  surface refutes; the capacity control ties on both. The axis does not merely fail
  to close on one surface --- it closes on two.}
  \label{fig:nexttrack-stage-stack}
\end{figure}

\emph{Interpretation.} \textbf{The axis is closed, and it is the first null in this
arc to close on both surfaces.} All five 4-stage arms are refuted by paired CI on
one-shot recall \emph{and} on the walk with the instrument gate green, and the
capacity control is indistinguishable on both. This is the sixth consecutive
architecture null on this corpus and, unlike the five before it, it satisfies the
standing rule that a one-shot verdict says nothing about the surface the product
serves.

\textbf{The user's mechanism hypothesis is refuted, and that is the most reusable
fact here.} Arm E put the MLP on a \emph{gated hidden state} rather than on raw
signed PCA, which the rectifier story predicts should remove the harm. It does not:
E reaches \num{118} of \num{1431}, ties two other arrangements, beats only the
worst, and comes nowhere near the control's \num{178}. So half-wave rectification
of the \emph{raw input} is not the mechanism --- whatever is wrong survives moving
the rectifier downstream of a gate. One run bought a clean refutation of a
plausible and cheap-to-believe story.

\textbf{The damage is depth itself, and three controls close the alternatives
simultaneously.} It is not capacity: the $\num{3.24}\times$ control ties on both
surfaces. It is not regularization or early stopping: \emph{every} 4-stage arm
trains \emph{longer} than the control (\numrange{21}{40} epochs against \num{14}),
and two of five reach a \emph{lower} best validation loss, so the pre-authorized
dropout-count contingency failed both of its trigger clauses and was not spent ---
the four-fold dropout application is exonerated. And it is not the MLP: a leading
per-step MLP at depth 4 is recall-neutral and \emph{holisticness-positive} against
the pure-recurrent arrangement ($\Delta H$ $+\num{0.00291}$, CI above zero). What
remains is the depth of the recurrent chain, and its cost is large ---
\numrange{55}{77} hits of \num{1431}.

\textbf{Bracketing beats interleaving at byte-identical capacity.}
\texttt{ann+gru+gru+ann} beats \texttt{ann+gru+ann+gru} by $+\num{0.01537}$ recall
with the CI above zero at \emph{the same parameter count by formula} --- the
tightest-controlled topology comparison the campaign could make. Interposing a
per-step MLP \emph{between} two recurrences is worse than putting MLPs on the
outside, which is consistent with the depth reading: breaking the recurrent chain
in the middle costs more than adding depth at the ends. Two further contrasts are
flat: the number of recurrent stages (2 versus 3) straddles on recall, and the
rectifier's position buys nothing beyond beating the worst arrangement.

\textbf{A leading per-step MLP is \emph{not} harmful at depth 4, which narrows
rather than confirms the $N=1$ result.} Section~\ref{sec:preencoder}'s ``leading
MLP is harmful'' was measured at $N=1$ against a single GRU; at depth 4 the same
addition is free-to-slightly-good. The pre-encoder pitfall should be scoped to the
shallow case.

\textbf{Objective and retrieval dissociate at depth, and this is the sharpest
instance yet.} Arm B achieves the \emph{lowest} validation InfoNCE in the entire
batch (\num{6.5378} against the control's \num{6.5832}) while retrieving
\textbf{\num{55} fewer hits} (\num{123} against \num{178}). Every prior
architecture null on this corpus was consistent with ``the extra machinery does
nothing''; this one shows the extra machinery does something measurable to the loss
and that it is the \emph{wrong} something. The operational consequence is blunt:
select on Recall@10 and the walk, never on the training objective. One caveat rides
with it --- B ran to the \num{40}-epoch cap with patience never triggering, so its
recall is a \emph{lower bound}, and a longer re-run is the clean test of the
dissociation (the prediction being a better validation loss and no better recall).

\textbf{The eagerness facets improve monotonically as the arms get worse, and the
floors caught it.} Every 4-stage arm de-eagers substantially (concentration
\num{0.359}$\to$\numrange{0.263}{0.318}) and raises \texttt{ild}; under a naive
holisticness-primary read some would look interesting. All are killed by the
absolute recall floor and by the per-baseline relevance floor --- the same mechanism
that has now saved four campaigns. Note that the canonical \emph{absolute}
\num{0.43} relevance figure would have passed several of these arms, so the floor
must be derived per baseline.

\paragraph{Two by-products, and the governance line.} A previously closed $N=2$
relative --- a dual tower with a pre-MLP on one side, closed on a one-shot recall
tie --- was walk-evaluated for the first time here and is \textbf{refuted on the
walk} ($\Delta$\texttt{vibe} $-\num{0.0255}$ [$-\num{0.0432}$, $-\num{0.0096}$]),
its one real gain being less drift. That is the same signature a recurrent-cell
variant showed earlier --- better stride and drift, worse vibe --- which now looks
like a general property of adding pre-recurrence machinery rather than anything
specific to either arm. Second, the grounding metric again discriminated almost
nothing, separating one arm of nine, consistent with its recorded arithmetic
(between-arm range \num{0.0060} against a per-session standard deviation of
\num{0.1417}, needing ${\sim}\num{2140}$ sessions): it is a floor that catches
degeneracy, not a ranking metric, and it was used only that way. Nothing was
promoted as champion, nothing was registered as a crown claim and nothing was
deployed; the seven arms were promoted \emph{for measurement only}, each carrying
the verbatim note ``promoted to enable walk evaluation; not a crown claim'', which
flooded the best-models group and was handed to the curation agent.

\paragraph{Two process findings worth keeping.} First, the predictor-checksum gate
fired \emph{on a file the design did not name for copying} --- it listed one module
to copy and a shared module only to \emph{verify}, and the verification failed. The
divergence turned out to be inert for this batch, but that is a \emph{post-hoc}
judgement that required reading the diff; the rule that caught it is the one worth
keeping. \textbf{Verify the checksum of every file in the family's import closure,
not just the ones you copied.} Second, the approved launch body was rejected by the
approved registry block: the design's JSON sent six hyperparameters the block does
not declare, so the first posts failed outright. It was resolved by dropping the
six keys \emph{after proving the drop is a no-op} --- the predictor's defaults supply
exactly those values and the resolved hyperparameter dict is identical either way,
verified by direct comparison --- and the registry block was left exactly as
approved.
\emph{Source: \texttt{2026-08-01-stage-stack-topology.md}.}

\section{Parallel tower banks: the $N$-axis closed on both surfaces}
\label{sec:towerbank}
\emph{Rationale.} The walk reversal of Section~\ref{sec:walk-reversal} left this
corpus with exactly one measured architecture \emph{win}, and with a stated
mechanism for it: a dual tower whose second view is \texttt{cummean} holds a mood
because it carries a \textbf{fixed reference never overwritten by recurrence}. A
mechanism of that shape makes a testable prediction --- it should be
\emph{dose-responsive}. If one fixed-reference tower buys a held vibe plus
$+\num{1.5}$ genres, then $N$ towers, enough of them on the raw \texttt{latent}
view to hold the recall floor and the rest on \texttt{cummean}, should buy
\emph{more} of the same. This campaign runs that dose curve, and it is the
follow-up the reversal itself asked for.

The vehicle is a new predictor, \texttt{seq-bank}: $N$ \emph{parallel}
single-layer bare GRU towers over configurable causal views, their per-step
hidden states concatenated and read by \textbf{one linear head}. Its design
choice is to collapse tower count and view multiset into a \emph{single}
\texttt{views} hyperparameter, so $N$ is not a separate axis that can drift out of
sync with the views. What the predictor \emph{cannot} express is as important as
what it can: it declares exactly \num{15} hyperparameters and none of
\texttt{fusion\_layers}, \texttt{fusion\_hidden}, \texttt{pre\_hidden},
\texttt{pre\_layers}, \texttt{layers}, \texttt{num\_layers}, \texttt{arch},
\texttt{bidirectional}, \texttt{residual} or \texttt{eager\_beta} is among them.
Every mechanism this document has already refuted --- the MLP fusion head
(Section~\ref{sec:dualtower}), the per-step pre-encoder
(Section~\ref{sec:preencoder}), the soft anti-eager regularizer
(Section~\ref{sec:holisticness}) --- is therefore \emph{structurally unreachable}
rather than merely unused, and a \texttt{parse\_views} guard rejects any spec
without a \texttt{latent} tower, which puts the registered
\numrange{0.058}{0.087} \texttt{delta}/\texttt{cummean} pathologies out of reach
as well.

Two hypotheses were pre-registered. The primary one is the \emph{dose}
hypothesis above. The independent sub-hypothesis is the mechanism the dual-tower
scan had already established: if the raw \texttt{latent} view \emph{subsumes} its
siblings, then $N$ identical \texttt{latent} towers should buy nothing, whatever
$N$ is.

\paragraph{Design and the two surfaces.} Ten runs on the canonical split
\dsid{seq-20260715-131139} (\num{1431} test sessions, cold-item rate
\num{0.5507}), \SI{84.2}{\minute} of compute, zero failures. Everything but
\texttt{views} / \texttt{hidden} / \texttt{dropout} is held: \texttt{loss}
InfoNCE, $\tau=\num{0.07}$, \num{40} epochs, patience \num{6}, lr \num{0.001},
batch \num{128}, seed \num{1337}, MMR off, artist cap off. \textbf{Capacity is
matched at \num{3.00}$\times$ the h256 baseline} for every bank arm \emph{and}
for the single-GRU control C1 (h512), with the counts re-verified from the
artifacts rather than trusted from the design --- \texttt{seq\_bank} emits
\texttt{n\_params} into \texttt{metrics.json}, and the mutual spread across the
seven \num{3.00}$\times$ arms including C1 is \SI{0.405}{\percent}. The decision
rule was written on \emph{both} surfaces before launch, and both legs are
load-bearing:
\begin{description}
  \item[S1, one-shot recall@10.] A win requires a paired $\Delta$ above zero
  against \emph{both} the h256 anchor C0 \emph{and} the parameter-matched single
  GRU C1. Two mandatory floors apply: an absolute recall floor
  $\ge\num{0.108}$ (\num{155} of \num{1431}) and a per-baseline relevance floor
  (paired $\Delta$\,music@10 vs C0 must not fall below zero).
  \item[S2, the generated walk.] \num{40} held-out sessions $\times$ \num{20}
  steps at artist cap 1, paired against the \emph{in-batch} C0, on the walk
  metrics of Section~\ref{sec:walk-metrics}. The primary cell is anchor
  \num{0.4} / stride \num{0.3}; a saved definition additionally requires
  \textbf{the same sign at the deployed cell} $a\num{0.8}$ / $s\num{0.5}$.
\end{description}

\paragraph{The instrument was pinned before anything was measured.} All ten
pre-registered gates came back green, and three of them are worth recording
because they make the comparison unusually tight. A \emph{functional-equivalence}
gate showed \texttt{seq-bank} to be a strict generalization of both incumbents
rather than a new family: at \texttt{views="latent"} it reproduces
\texttt{SeqNextLatent} and at \texttt{views="latent+latent"} it reproduces the
\texttt{fusion\_layers}${}=0$ \texttt{DualTowerNextLatent}, with
$\max|\Delta|$ \textbf{\texttt{0.00e+00}} on both \texttt{forward} and
\texttt{predict\_next} at prefix lengths $T=1,2,5,17,40$ and identical parameter
counts on both sides; an extra training-stream check matched per-epoch train
\emph{and} validation loss to the last digit over three epochs and all ten
state-dict tensors at $\max|\Delta|$ \texttt{0.00e+00}. Both bit-identity gates then
hit exactly: C0 reproduced its on-record recall@10
\num{0.12299091544374564} $=$ \num{176} of \num{1431} (a \textbf{10th}
reproduction, with \texttt{music@10}, H@10 and MRR matching to the last digit)
and R2 reproduced the dual-tower \texttt{latent/latent} f0 arm at
\num{0.1187980433263452} $=$ \num{170} of \num{1431}, so the design's band
fallback was never needed. And the walk instrument was checked against the
shipped \texttt{gru.onnx}: $\Delta$\texttt{vibe} $+\num{0.0028}$ and
$\Delta$\texttt{stride\_err} $+\num{0.0005}$ both straddle zero, so S2 verdicts
against the in-batch C0 are not an artefact of the base.

\emph{Results, S1.} Table~\ref{tab:towerbank-s1} is the batch. \textbf{No arm is
bolded as a winner because there is none}: no arm achieves even the first leg of
the win rule, all nine non-anchor arms are ties against C0, and \emph{five of the
eight bank arms lose to the parameter-matched single GRU C1 with the CI below
zero}. Both floors pass for every arm --- the worst is V1 at \num{162} hits,
clear of the \num{155} floor, and the most negative relevance $\Delta$ is V2's
$-\num{0.006109}$, straddling.

\begin{table}[ht]
  \centering
  \caption{The tower-bank batch on S1, ten runs on \dsid{seq-20260715-131139}
  (\num{1431} test sessions), sorted by hits. The batch was launched in four
  chunks; every run id is \protect\runid{run-20260802-} plus the chunk time and
  the hex suffix in the id column, in full: \texttt{182050-91b54},
  \texttt{182050-2e6f2}, \texttt{183404-91e31},
  \texttt{183404-62604}, \texttt{183404-67185},
  \texttt{190415-1d662}, \texttt{190415-90866},
  \texttt{190415-b5928}, \texttt{192627-17502},
  \texttt{192627-836b1}. Paired $\Delta$ is a per-session paired
  bootstrap (\num{2000} resamples, rng \num{1337}) computed from the stored
  per-row \texttt{predictions.json}; a CI straddling zero is a \emph{tie}, and
  \textsuperscript{$\ast$} marks a $\Delta$ vs C1 whose CI lies entirely
  \emph{below} zero. d\num{0.1}\,/\,d\num{0.0} is the dropout of the
  byte-identical D2\,/\,R2 pair. Capacity: the seven arms C1, B2, B3, B4, V1, V2
  and V3 span \numrange{2.996}{3.007}$\times$ the h256 baseline (mutual spread
  \SI{0.405}{\percent}), D2 and R2 sit at exactly \num{2.000}$\times$ and C0 at
  \num{1.000}$\times$. View codes:
  lat\,$=$\,\texttt{latent}, cum\,$=$\,\texttt{cummean},
  del\,$=$\,\texttt{delta}. H@10 is the \emph{generation-3} holisticness
  (Section~\ref{sec:power-generations}), in raw next-track units. R@10, H@10 and
  mus@10 higher-better; a\_conc is lower-better but display-only and is
  deliberately \emph{not} bolded --- no arm's eagerness is a claim in this
  campaign. Best cell per metric bolded; C1 and B4 tie exactly on R@10 and both
  are bolded.
  \emph{Source: \texttt{2026-08-02-tower-bank-parallel-towers.md}.}}
  \label{tab:towerbank-s1}
  {\scriptsize\setlength{\tabcolsep}{1.9pt}\begin{tabular}{lllrlllccc}
    \toprule
    Arm & id & views / width & params & R@10 (raw) & $\Delta$ vs C0 [95\% CI] & $\Delta$ vs C1 & H@10 & mus@10 & a\_conc \\
    \midrule
    C1 \emph{ctrl} & \texttt{91e31} & 1$\times$GRU h512 & \num{1182912} & \best{0.12998 (186)} & $+0.0070$ [$-.002$, $+.016$] & --- & 0.02380 & 0.4581 & 0.3825 \\
    B4 & \texttt{1d662} & lat$\times$4 h211 & \num{1187700} & \best{0.12998 (186)} & $+0.0070$ [$-.002$, $+.017$] & $+0.0000$ & 0.02446 & \best{0.4594} & 0.3716 \\
    D2 & \texttt{62604} & lat$\times$2 h256 d\num{0.1} & \num{789696} & 0.12369 (177) & $+0.0007$ [$-.008$, $+.008$] & $-0.0063$ & 0.02426 & 0.4567 & 0.3698 \\
    C0 \emph{base} & \texttt{91b54} & 1$\times$GRU h256 & \num{394944} & 0.12299 (176) & --- & $-0.0070$ & \best{0.02454} & 0.4558 & 0.3559 \\
    B3 & \texttt{90866} & lat$\times$3 h256 & \num{1184448} & 0.12020 (172) & $-0.0028$ [$-.010$, $+.006$] & $\mathbf{-0.0098}$\textsuperscript{$\ast$} & 0.02358 & 0.4554 & 0.3764 \\
    R2 & \texttt{2e6f2} & lat$\times$2 h256 d\num{0.0} & \num{789696} & 0.11880 (170) & $-0.0042$ [$-.013$, $+.004$] & $\mathbf{-0.0112}$\textsuperscript{$\ast$} & 0.02378 & 0.4554 & 0.3706 \\
    B2 & \texttt{67185} & lat$\times$2 h334 & \num{1186560} & 0.11880 (170) & $-0.0042$ [$-.013$, $+.003$] & $\mathbf{-0.0112}$\textsuperscript{$\ast$} & 0.02336 & 0.4554 & 0.3799 \\
    V3 & \texttt{17502} & lat$+$cum$+$cum h256 & \num{1184448} & 0.11880 (170) & $-0.0042$ [$-.016$, $+.006$] & $-0.0112$ & 0.02320 & 0.4497 & 0.3581 \\
    V2 & \texttt{836b1} & lat$+$cum$+$del h256 & \num{1184448} & 0.11600 (166) & $-0.0070$ [$-.020$, $+.005$] & $\mathbf{-0.0140}$\textsuperscript{$\ast$} & 0.02245 & 0.4497 & 0.3547 \\
    V1 & \texttt{b5928} & lat$\times$2$+$cum h256 & \num{1184448} & 0.11321 (162) & $-0.0098$ [$-.022$, $+.001$] & $\mathbf{-0.0168}$\textsuperscript{$\ast$} & 0.02299 & 0.4500 & 0.3572 \\
    \bottomrule
  \end{tabular}}
\end{table}

\paragraph{The sharpest number in the batch is a zero.} B4 --- four independent
h211 \texttt{latent} towers, \num{1187700} parameters --- and C1 --- one h512
GRU, \num{1182912} parameters --- return the \textbf{identical \num{186} hits of
\num{1431}}, a paired $\Delta$ of precisely $\num{0.000000}$
$[-\num{0.00839},+\num{0.00839}]$. At a matched ${\sim}\num{3.00}\times$ budget,
\emph{topology and width are interchangeable on this corpus}: whatever four
towers with different random initializations discover, a single wider recurrence
already had. Figure~\ref{fig:nexttrack-tower-bank} (left) plots the $N$-axis at
matched capacity, and the line through $\num{186}/\num{170}/\num{172}/\num{186}$
is flat inside the split's own measured resolution
(Section~\ref{sec:power-stats}).

This is also the cleanest demonstration on record of \emph{why the
parameter-matched control leg is mandatory}. Read against the h256 anchor alone,
B4's $+\num{10}$ hits would have looked like a topology gain and the five arms
that lose to C1 would all have read as ties --- the anchor cannot distinguish
``the bank helped'' from ``the extra parameters helped, and the bank spent them
worse than width would have''. The C1 leg does exactly that, and it is what turns
a batch of nine ties into a closure. Note too what the flatness costs the
\emph{other} direction: C1's $+\num{10}$ hits over C0 is itself not significant
($+\num{0.006988}$ $[-\num{0.00210},+\num{0.01607}]$), the \textbf{seventh}
parameter-matched width tie on record, so the campaign cannot claim h512 as an
improvement either despite it holding the highest single-GRU hit count measured
(\num{186} against h454's \num{179}). The single-GRU width points on PCA-192 now
read h256\,$\to$\,\num{176}, h297\,$\to$\,\num{162}, h425\,$\to$\,\num{173},
h454\,$\to$\,\num{179}, h512\,$\to$\,\num{186}, h537\,$\to$\,\num{174} ---
non-monotone, spanning \num{0.0168}, with no trend.

\emph{Results, S2.} Table~\ref{tab:towerbank-s2} is the walk. Two arms --- B2
(\texttt{latent}$+$\texttt{latent} h334) and V3
(\texttt{latent}$+$\texttt{cummean}$+$\texttt{cummean}) --- reached a genuine
\textbf{TRADE} at the primary cell, each buying a drift improvement with the CI
above zero while holding vibe. Both then \textbf{flip} at the deployed cell: the
drift gain evaporates into a straddle and \texttt{vibe} --- the non-negotiable
line --- goes below zero for both. Under the pre-agreed rule, \textbf{no
definition was saved}.

\begin{table}[ht]
  \centering
  \caption{The tower-bank batch on S2, the generated walk: \num{40} held-out
  sessions $\times$ \num{20} steps at artist cap 1, paired against the
  \emph{in-batch} C0 (\num{2000} resamples). The upper block is the primary cell
  $a\num{0.4}$ / $s\num{0.3}$; the lower block re-measures the two TRADE arms at
  the \emph{deployed} cell $a\num{0.8}$ / $s\num{0.5}$, where the save clause
  requires the same sign. \texttt{vibe}, $\Delta$drift, $\Delta$genres and
  $\Delta$ground higher-better; \texttt{stride\_err} lower-better. CI bounds are
  quoted only where the interval clears zero, as in the source report, and
  \textsuperscript{$\ast$} marks every such interval --- a \emph{gain} on
  $\Delta$drift, a \emph{refutation} on $\Delta$vibe and $\Delta$genres.
  \emph{Source: \texttt{2026-08-02-tower-bank-parallel-towers.md}.}}
  \label{tab:towerbank-s2}
  {\scriptsize\setlength{\tabcolsep}{2.6pt}\begin{tabular}{lccllllc}
    \toprule
    Arm & vibe & stride\_err & $\Delta$vibe & $\Delta$stride\_err & $\Delta$drift & $\Delta$genres & $\Delta$ground \\
    \midrule
    \multicolumn{8}{l}{\emph{Primary cell $a\num{0.4}$ / $s\num{0.3}$}} \\
    C0 (base) & 0.566 & 0.130 & --- & --- & --- & --- & --- \\
    \textbf{B2} \emph{TRADE} & 0.567 & 0.139 & $+0.0017$ & $+0.0099$ & $\mathbf{+0.0339}$ [$+.007$, $+.063$]\textsuperscript{$\ast$} & $-0.225$ & $+0.0044$ \\
    \textbf{V3} \emph{TRADE} & 0.556 & 0.124 & $-0.0094$ & $-0.0045$ & $\mathbf{+0.0286}$ [$+.001$, $+.059$]\textsuperscript{$\ast$} & $-0.450$ & $-0.0024$ \\
    C1 & 0.563 & 0.130 & $-0.0020$ & $+0.0015$ & $+0.0194$ & $-0.375$ & $+0.0008$ \\
    D2 & 0.559 & 0.129 & $-0.0065$ & $+0.0005$ & $+0.0104$ & $-0.750$ & $+0.0003$ \\
    B4 & 0.566 & 0.139 & $+0.0006$ & $+0.0103$ & $+0.0164$ & $-0.775$ & $+0.0004$ \\
    B3 & 0.557 & 0.135 & $-0.0078$ & $+0.0060$ & $+0.0146$ & $\mathbf{-1.025}$\textsuperscript{$\ast$} & $+0.0154$ \\
    V2 \emph{refuted} & 0.551 & 0.127 & $\mathbf{-0.0144}$ [$-.029$, $-.001$]\textsuperscript{$\ast$} & $-0.0019$ & $+0.0259$ & $-0.025$ & $-0.0148$ \\
    V1 \emph{refuted} & 0.548 & 0.119 & $\mathbf{-0.0174}$ [$-.035$, $-.001$]\textsuperscript{$\ast$} & $-0.0099$ & $+0.0227$ & $\mathbf{-1.725}$\textsuperscript{$\ast$} & $+0.0017$ \\
    \midrule
    \multicolumn{8}{l}{\emph{Deployed cell $a\num{0.8}$ / $s\num{0.5}$ --- both TRADEs collapse}} \\
    B2 \emph{refuted} & --- & --- & $\mathbf{-0.0064}$ [$-.012$, $-.001$]\textsuperscript{$\ast$} & $-0.0107$ & $-0.0052$ & $-0.175$ & $+0.0006$ \\
    V3 \emph{refuted} & --- & --- & $\mathbf{-0.0112}$ [$-.017$, $-.006$]\textsuperscript{$\ast$} & $-0.0074$ & $+0.0043$ & $-0.400$ & $-0.0039$ \\
    \bottomrule
  \end{tabular}}
\end{table}

\begin{figure}[ht]
  \centering
  \includegraphics[width=0.95\textwidth]{nexttrack_tower_bank}
  \caption{The parallel-tower-bank axis, closed on both surfaces.
  \textbf{Left:} the $N$-axis at matched capacity (${\sim}\num{3.00}\times$ the
  h256 baseline). Error bars are the paired bootstrap CI against C0 mapped onto
  the hit scale; the grey band is the split's measured resolution
  ($\pm\num{0.0124}$ recall@10 $=$ $\pm\num{17.7}$ hits,
  Figure~\ref{fig:nexttrack-power}), and every arm sits inside it. Four h211
  towers (B4) and one h512 GRU (C1) return the same \num{186} hits, paired
  $\Delta$ $\num{0.000000}$. \textbf{Middle:} the two arms that earned a walk
  TRADE at the primary cell both flip at the deployed cell --- the drift gain
  that bought the TRADE decays to a straddle while \texttt{vibe} goes below zero.
  Grey markers are point estimates whose CI straddles (the source report quotes
  bounds only for intervals clearing zero). \textbf{Right:} the campaign's one
  positive result, and it is not about topology --- at byte-identical parameters
  (\num{789696} both, \texttt{model.pt} \num{3162277} bytes both) dropout
  \num{0.1} beats dropout \num{0.0}, CI above zero.}
  \label{fig:nexttrack-tower-bank}
\end{figure}

\emph{Interpretation.} \textbf{The dose hypothesis is refuted, and it fails in
the direction opposite to the prediction.} Adding fixed-reference
\texttt{cummean} towers did not buy more of what $N=2$ bought; it monotonically
destroyed the walk. Ranked by $\Delta$\texttt{vibe} at the primary cell the
ordering is B4/B2 (all-\texttt{latent}, ${\approx}0$) $>$ B3 $>$ V3 (one
\texttt{cummean}, $-\num{0.0094}$) $>$ V2 ($-\num{0.0144}$, CI below zero) $>$
V1 ($-\num{0.0174}$, CI below zero), and the same ordering holds on S1, where the
three \texttt{cummean}-bearing arms occupy the \emph{bottom three} places
(\num{162}, \num{166}, \num{170} hits) and are the only arms whose
$\Delta$\,music@10 point estimates fall past $-\num{0.005}$. The irony is worth
stating plainly: \textbf{the tower added in order to hold a mood is the tower
that made the walk's mood measure worse} --- the only two arms whose
\texttt{vibe} is refuted outright are the two that carry the most
\texttt{cummean}. Whatever the deployed $N=2$ \texttt{latent}/\texttt{cummean}
engine is doing, ``more \texttt{cummean} is more mood-holding'' is not it, and
the structural reading that motivated the design does not survive being dosed.

\textbf{The sub-hypothesis is confirmed and extends cleanly from $N=2$ to
$N=4$.} B2 (\num{170}), B3 (\num{172}) and B4 (\num{186}) all straddle C0 and
show no trend distinguishable from width noise. This is the same mechanism the
dual-tower scan found (Section~\ref{sec:dualtower}) --- the raw \texttt{latent}
view subsumes its siblings --- now measured at three tower counts instead of two,
which is what closing an axis by dose rather than by a single point requires.

\textbf{The one positive finding is a by-product of the null, and it lands on the
deployed engine.} D2 and R2 differ in \emph{one} hyperparameter: parameters are
byte-identical (\num{789696} both, \texttt{model.pt} \num{3162277} bytes both),
only \texttt{dropout} moves, and $\Delta$recall@10 is $+\num{0.004892}$
$[+\num{0.000699},+\num{0.009783}]$, \textbf{CI above zero} --- dropout
\num{0.1} \emph{beats} dropout \num{0.0} (\num{177} against \num{170} hits).
Read carefully this does not overturn the record's three-times-confirmed ``f0
beats f1'' result; it does something more awkward. Dropout enters
\texttt{seq\_dualgru} only through a pre-MLP, an inter-layer gap or the fusion
MLP, and at \texttt{fusion\_layers}${}=0$ with one layer and no pre-encoder
\emph{all three are absent} --- the f0 path constructs \textbf{zero} dropout
modules. So every f0 arm on record, \textbf{including the deployed showcase
engine} (Section~\ref{sec:walk-limits}), trained under-regularized and was
leaving ${\approx}+\num{0.005}$ recall on the table; and the three recorded
f0-vs-f1 gaps ($+\num{0.016}$ / $+\num{0.020}$ / $+\num{0.0266}$) were measured
against an f1 that \emph{had} regularization its f0 comparator lacked, so they
are not like-for-like and the true topology advantage of f0 is if anything
\emph{larger} than recorded. Wherever those three figures are quoted as clean
topology measurements they should be re-stated as not like-regularized. The
cheapest actionable follow-up on the whole board falls out of this: retrain the
deployed engine at \texttt{dropout}${}=\num{0.1}$ and re-read both surfaces ---
a free improvement, not an architecture bet.

\textbf{The validation-loss dissociation gets a fourth confirmation, now as a
rank correlation} (Figure~\ref{fig:nexttrack-bank-valloss}). Across the ten arms,
Spearman(best validation InfoNCE, hits) $=-\num{0.2954}$, $p=\num{0.407}$ --- no
usable predictive power, and if anything the wrong sign. The sharp instances are
brutal: B4 and B2 differ by \num{0.0003} in validation loss and by \num{16}
hits; V1 holds the fourth-best validation loss in the batch and the \emph{worst}
hit count; C0 holds the \emph{worst} validation loss and beats six of the eight
bank arms; and R2 reaches a better validation loss than D2 while retrieving seven
fewer hits, which makes the dissociation and the dropout finding two faces of one
coin. The decision rule never read validation loss, and this is why.

\begin{figure}[ht]
  \centering
  \includegraphics[width=0.8\textwidth]{nexttrack_bank_valloss}
  \caption{The training objective does not rank retrieval --- fourth
  confirmation, and its sharpest instance yet. Ten arms of the tower-bank batch:
  best validation InfoNCE against hits of \num{1431}. Spearman
  $=-\num{0.2954}$, $p=\num{0.407}$. B4 and B2 are \num{0.0003} apart in
  validation loss and \num{16} hits apart; C0, the worst fit in the batch, beats
  six of the eight banks. Every arm early-stopped (epoch \num{14} seven times,
  then \num{21} and \num{24} twice), so no arm's recall is a lower bound imposed
  by the epoch cap.}
  \label{fig:nexttrack-bank-valloss}
\end{figure}

\textbf{Both walk TRADEs were cell-local, which is a methodological result in
its own right.} B2 and V3 each earned a real $\Delta$drift above zero at
$a\num{0.4}$/$s\num{0.3}$ with vibe held, and both flipped to a vibe below zero
at $a\num{0.8}$/$s\num{0.5}$ with the drift gain gone. Had the design not
pre-registered the secondary-cell requirement, this campaign would have reported
two positive walk results and saved two definitions. \textbf{The secondary-cell
clause is the clause that did the work here}, and it should stay mandatory in
every future walk design. The pre-registered stride probe points the same way: at
$a\num{0.4}$ over $s\in\{\num{0.2},\num{0.3},\num{0.4}\}$ all three
\texttt{cummean} arms peak at $s\num{0.3}$, the cell they were already measured
in --- the third campaign in a row in which re-tuning the stride rescues nothing
--- and at $s\num{0.2}$ their $\Delta$\texttt{vibe} turns
positive-straddling only by taking visibly shorter steps
($\Delta$\texttt{stride\_err} above zero for all three) and shedding genre
variety ($-\num{1.90}$ / $-\num{1.78}$ / $-\num{1.00}$, all below zero). That is
precisely the trade the walk win clause exists to reject.

\textbf{The failure signature is the opposite of the preceding architecture
campaign's, and that is diagnostic.} No arm here was refuted by a floor: the
recall floor and the relevance floor both pass for all ten runs, and eight of the
nine non-anchor arms sit at \emph{higher} artist concentration than the anchor
(\numrange{0.3547}{0.3799} against \num{0.3559}), i.e.\ the arms did not
de-eager. The floors were therefore not the operative mechanism --- these arms
are not degenerate, they are simply \textbf{flat}. That is a different failure
mode from the immediately preceding depth-in-series campaign
(Section~\ref{sec:stagestack}), in which all five arms fell \emph{through} the
recall floor, and it is why this axis closes as a null rather than as a harm.

\paragraph{The scope of the closure, precisely.} What closes is
``$N\ge3$ parallel bare-GRU tower banks at \texttt{fusion\_layers}${}=0$'', and
only that, to the standard Section~\ref{sec:stagestack} set for a closure on this
corpus: both surfaces,
walk instrument verified. It does \textbf{not} re-close the dual-tower family,
which stays reopened (Section~\ref{sec:walk-reversal}) and remains the deployed
engine. Nothing was promoted as champion, no definition was saved, and no
dataset was built; the eight measurement arms were promoted for measurement only,
each carrying the verbatim note ``promoted to enable walk evaluation; not a crown
claim'', and the resulting flooding of the best-models group (to \num{15}
entries) was disclosed up front and handed to the curation agent. Zero of the
three pre-authorized contingency runs fired, and the epoch trace says why: the
dropout-confound extension needed R2 to early-stop later than D2 and both stopped
at epoch \num{14}; the epoch-cap contingency needed an arm to reach the \num{40}
epoch cap and every arm early-stopped; the width-bracket contingency needed C1 to
lose, and it did not.

\paragraph{One limitation is inherited rather than introduced.} This campaign's
premise --- the dose mechanism --- was read off a finding that has \emph{no
campaign report of its own}: the walk-surface reopening of the dual-tower family
(Section~\ref{sec:walk-reversal}) exists only as a \texttt{PROJECT-FACTS.md}
roll-up, a git commit, and the transcription in this document. The
harness-reproduction gate re-derived all ten of that reopening's recorded means
exactly and re-derived its deltas, so the \emph{measurement} is reproducible and
the instrument has not moved; what is missing is the primary write-up, and it
should be backfilled from the commit and the facts entries.
\emph{Source: \texttt{2026-08-02-tower-bank-parallel-towers.md}.}

\section{Conclusions}
\label{sec:conclusions}
Twenty-one campaigns and one research memo (Jul--Aug 2026) on the next-track
retrieval task converge on a clear
story. (1) A session sequence model, even at its best (GRU-infonce), only
\emph{ties} the first-order Markov co-listening bar on Recall@10 and loses on
MRR: predicting a next-item latent and retrieving by one cosine hop is a harder
route to the exact next track than a bigram that directly memorizes adjacency.
(2) The training \emph{loss} is the dominant single-model lever (InfoNCE
$\approx 2.3\times$ cosine); network \emph{topology} --- depth, width, residual,
direction --- is a spent lever, and no architecture beats the bar. (3) The
breakthrough is \emph{rank fusion}: an R+M z-blend of the GRU and Markov legs
beats the bar by $\sim$$+61\%$ Recall@10 (three-seed confirmed), because the two
legs rank \emph{different} correct tracks. The win is about signal
\emph{diversity}, not model count --- adding a redundant neural leg does not
help, and a learned XGB stacker \emph{underperforms} the fixed z-blend, because
its pointwise training does not capture the ranking objective. (4) An
item-representation sweep found latent \emph{dimension}, not the compressor, to
be the lever: linear PCA beats the nonlinear autoencoder at every dim $\ge 64$
(the acoustic-preservation hypothesis is refuted), the curve peaks at PCA-192 and
turns over by 256, and PCA-128 is the confirmed single-model base. (5) On the
rich PCA space the one genuinely additive new signal --- \emph{content-kNN},
which only \emph{tied} the champion on the AE-64 space --- becomes a win: a
three-leg R+M+C z-blend on PCA-192 is the current best-on-record row at Recall@10
\num{0.1859}, MRR \num{0.1025} (paired $\Delta$ over the prior champion
$+\num{0.0133}$, CI above zero, confirmed across five reproductions). (6) Two
honest caveats bound the crown: the reproduction confirm had zero spread because
\texttt{seq-nexttrack} is deterministic on this CPU (so it verified determinism,
not seed-robustness), and the paired-$\Delta$ lower bound clears zero by only
$+\num{0.0007}$ on a single split --- hardening it needs a second leak-free test
split. (7) Every ranking predictor here is train-only and off-server, so the
crown is a best-on-record \emph{row}, not a promoted, app-exportable model; the
prior champion R+M on AE-64 remains the productization anchor, and registering a
three-leg blend with a \texttt{predict} subcommand is the standing top follow-up.
(8) A subsequent literature-fit campaign then \emph{superseded} the crown with the
first genuinely-additive new signal since content: a supervised \emph{learned
content projection} (\Rp+M+\Cp), an InfoNCE metric map of the PCA-192 item space
fit toward next-track adjacency in which both the GRU and content legs retrieve, at
Recall@10 \num{0.2117}, MRR \num{0.1220} (paired $\Delta$ over R+M+C $+\num{0.0259}$,
CI$>$0). It resolves both R+M+C caveats --- it is seed-robust across three seeds and
cold-reaching (its content leg scores cold targets as well as warm) --- and it was
\emph{registered, promoted and servable} (\texttt{blend-gru-markov-content-proj},
best-models rank 1), closing the productization gap; an objective / rank / $\tau$
scan confirms InfoNCE $\gg$ BPR and holds the default, and the crown then hardened
on a second disjoint 0.64-cold split (paired $\Delta$ $+\num{0.0426}$, CI$>$0, the
margin \emph{wider} on the colder split). Every borrowed literature-fit method
(session-kNN, rank-fusion / learned combiners, GRU-side loss tricks, Markov
smoothing) closed --- the only additive gain came from new signal, not smarter
fusion. (9) An item-vector representation exploration then confirmed \emph{why} the
projection is the champion: the equal-weighted PCA-192 geometry is a real handicap
for a raw model (six of eight whitened / acoustic-dropped / metric spaces lift the
frozen GRU, paired $\Delta$ CI$>$0), but the gain is fully \emph{absorbed} by the
learned projection (no space beats the champion on readout A) --- whitening and the
projection are the same fix, and the champion has already pulled it. (10) Finally, a
continuous ``sounds-alike'' metric, \emph{music@10}, was introduced and back-filled
across all 53 runs: it correlates with Recall@10 ($r=\num{0.609}$) but crowns a
different winner (the content-regressing \texttt{seq-nexttrack} lands the
musically-closest misses, best \num{0.482}), corroborating the exact-recall findings
without overturning them; it is display-only and the primary metric stays Recall@10. (11) Finally, a
measurement campaign added a display-only \emph{session-holisticness} suite
(artist\_adj@10 / mood\_coh@10 / ild@10) and two anti-eager levers. It found the
eagerness symptom \emph{tracks} Recall@10 inside the blend family --- the champion
is the most artist-eager model (artist\_adj@10 \num{0.610}) --- but the eagerness
is sourced from the \emph{Markov} bigram leg, not the RNN (single GRU / LSTM
artist\_adj@10 $\approx$ \num{0.345}), and the GRU and LSTM are statistically
tied on both recall and eagerness (the ``GRU wins because more eager'' hypothesis
is refuted). Of the two levers the eval-time MMR re-rank is the strong knob
($\lambda{=}\num{0.9}$ near-free: recall tie, ild@10 \num{0.406}$\to$\num{0.436};
$\lambda\approx\num{0.7}$ de-eagers for $\sim$2 recall points; music@10 even
rises as the list diversifies) while the training-time $\beta$ regularizer is weak
($\beta$ \numrange{0}{0.2} moves artist\_adj@10 only $-\num{0.0074}$); nothing
was promoted, and $\lambda\approx\num{0.9}$ is flagged for a future user
holisticness decision (Section~\ref{sec:holisticness}).
(12) Two final architecture campaigns then closed the last open architecture
axes, and both are nulls --- but nulls with mechanism. A jointly-trained
\emph{dual-tower} GRU, crossing the complete
$\{$\texttt{latent},\texttt{delta},\texttt{cummean}$\}^2$ view matrix against a
\SI{0.098}{\percent} parameter-matched single tower, fails to beat that tower on
any of eight arms (Section~\ref{sec:dualtower}). This is the \emph{fourth}
independent learned-combiner null and the one that matters most, because it
answers the standing rebuttal to the earlier three (``of course a bolted-on
combiner fails --- train the fusion end-to-end''): trained end-to-end at the
representation level, it still fails. The mechanism is that the raw
\texttt{latent} view already \emph{subsumes} the others --- a GRU over raw
latents integrates movement and running mean internally, and the one genuinely
complementary pair is \texttt{delta}$+$\texttt{cummean}, which beats both its
parents yet is hopeless in absolute terms. Two by-products are durable: the MLP
fusion head is actively \emph{harmful} (removing it wins in every regime
measured, three times over, Table~\ref{tab:fusion-depth}), and single-GRU width
is now closed on PCA-192 as it already was on AE-64. A learned per-step
\emph{pre-encoder} then fails harder (Section~\ref{sec:preencoder}): the
decisive single-tower ablation \emph{loses} outright ($-\num{0.0140}$, CI below
zero, \num{156} hits against \num{176}), the dose curve is flat to twice the
input dimension, and placement is irrelevant. Because a \emph{linear}
pre-encoder is provably expressivity-neutral here, the rectifier is the only
candidate cause --- and a counterfactual interpretability read confirmed it,
measuring the affine map neutral at ${\sim}0.8\sigma$ and the ReLU costing
${\sim}5.9\sigma$, with \texttt{frac\_exact\_zero} \num{0.5327} matching
\texttt{frac\_negative\_coords} \num{0.5327} to four digits: half-wave
rectification of a zero-centred latent deletes exactly its negative half. Across
five architecture nulls the pattern is now hard to dismiss --- on this
single-user, \SI{55}{\percent}-cold corpus, adding representation machinery
\emph{around} a single GRU over a good latent space does not help, while every
real gain came from improving the latent space itself or from the fixed z-blend.
(13) Finally, two facts about the \emph{instrument} rather than any model
reframe the whole architecture arc (Section~\ref{sec:power}). The nulls are
\textbf{power-limited}: at $n=\num{1431}$ the paired-bootstrap CI half-width is
\num{0.0124} on Recall@10 and ${\sim}\num{0.0018}$ on holisticness@10, while a
whole architecture family spans only \num{0.028} and \num{0.0031} --- so of
eight adjacent arms roughly two tiers are separable on recall and one on
holisticness, and resolving a \num{0.005} difference would need
${\sim}\num{8800}$ sessions. The honest reading is therefore that \emph{large
harms were detected reliably while moderate wins were never detectable at all},
which is a different claim from ``these architectures do not work'' and yields a
prioritization rule: further architecture arms are a poor use of runs on this
split whatever they return, because the axis's whole dynamic range is barely
twice the measurement error. And \emph{holisticness@10} exists in three
mutually non-comparable code generations (3-factor; 4-factor grounded;
grounded plus artist-concentration), whose differences dwarf every model
difference and under all of which the servable champion ranks \emph{last} ---
so this synthesis keeps its board denominated in Recall@10 and tags every
holisticness column with its generation (Table~\ref{tab:hgen}).
(14) Finally, a campaign about the \emph{surface} rather than the model finds
that the board above had been scoring the wrong one
(Section~\ref{sec:walk}). Everything else in this document is a one-shot top-10,
while the shipped product serves an autoregressive \emph{walk} --- and
\texttt{holisticness@10}'s \texttt{mood\_coh} anchors to the \emph{prefix}
centroid, which on a walk grows with the journey, so a run of individually legal
steps that ends in the opposite region of the space scores well at every step:
drift was invisible by construction, and the composite is $\num{5.5}\times$ more
sensitive to relevance than to eagerness besides. Calibrated against the user's
own \num{73632} consecutive transitions (median step cosine \num{0.261},
\SI{27.0}{\percent} negative), the engine the showcase shipped walks
$\num{2.5}\times$ too tight (median step \num{0.661}) --- ``flat despite lots of
artists'' was many small safe moves --- while the blend champion matches the real
stride but cannot hold a mood, because its InfoNCE projection encodes what
\emph{follows} rather than what \emph{sounds alike}. Three deterministic
retrieval-time controls fix the reachability: a seed anchor, a stride term
(deterministic on purpose --- sampling would break the permalink guarantee), and
a \emph{hard} per-artist cap, which succeeds exactly where the two soft
anti-eager levers failed because it cannot be outscored (on the pathological
playlist \num{59} of the top \num{60} candidates are one artist and the next
artist sits at rank \num{75}, against an MMR diversity budget worth
${\approx}\num{0.078}$ of normalized score; \texttt{eager\_beta} is mis-specified
on an absolute cosine when eagerness is a \emph{ranking} property). The cap
hardens on two disjoint splits and, unexpectedly, \emph{raises} \texttt{music@10}
at every setting: it trades literal for graded relevance. And on \num{40}
identical held-out sessions under paired bootstrap the \textbf{dual-tower family,
closed by three campaigns on recall grounds, is the winner on the stated
objective} --- \texttt{latent/cummean} at anchor \num{0.8} stride \num{0.5} takes
stride error $-\num{0.2545}$, genres $+\num{1.50}$ and drift $+\num{0.0351}$
with every CI clear of zero while \emph{holding} vibe, where the single GRU
($-\num{0.0409}$) and the champion ($-\num{0.1327}$) both pay vibe for the same
stride. The corroboration was already in the record as a \emph{negative}
(\texttt{seq\_dualgru.py:113}: the \texttt{cummean} tower ``added real linear
directions (CKA \num{0.458}) but zero next-item concepts tower A lacked'') ---
dead weight for next-track prediction is exactly session state for mood-holding,
though the stronger reading is refuted: \texttt{cummean} does \emph{not} hold a
vibe natively (\num{0.236} against \num{0.234}), it only hosts one once the
controls are applied. It is deployed as the showcase's default engine; and
\texttt{ground} straddles zero for every arm, so what is established is journey
\emph{shape}, not next-track truth, on one split, one seed and \num{40}
sessions, with nothing yet listened to and the primary metric unchanged.
(15) A closing architecture campaign then tested the reversal's own stated
mechanism --- ``holding a mood needs a fixed reference never overwritten by
recurrence'' --- as a \emph{dose}, and closed the axis on both surfaces
(Section~\ref{sec:towerbank}). A bank of $N$ parallel bare-GRU towers under one
linear head, matched at \num{3.00}$\times$ the h256 baseline and dosed from
$N=2$ to $N=4$, reaches no win on either surface: every arm \emph{ties} the
baseline on one-shot recall while \textbf{five of eight lose to a
parameter-matched single GRU} with the CI below zero, and the two arms that earn a
walk trade at the primary cell both flip to \texttt{vibe} below zero at the
deployed cell, so no definition was saved. The decisive datapoint is a zero: four
h211 towers (\num{1187700} params) and one h512 GRU (\num{1182912} params)
return the \emph{identical} \num{186} hits of \num{1431}, paired $\Delta$
\num{0.000000} --- \emph{topology and width are interchangeable at this budget},
which is also the cleanest demonstration on record of why the parameter-matched
control leg is mandatory: against the anchor alone, five losing arms would have
read as ties. The dose hypothesis is refuted \emph{ironically} --- the
\texttt{cummean} towers added in order to hold a mood are the worst arms in the
batch on both surfaces and the only two whose walk \texttt{vibe} is refuted
outright --- and the one CI-clear result in the whole batch is a by-product about
regularization, not topology: at byte-identical parameters dropout \num{0.1}
beats dropout \num{0.0} by $+\num{0.0049}$ recall@10. That is consequential
because the \texttt{fusion\_layers}${}=0$ path constructs \emph{zero} dropout
modules, so every f0 arm on record, \textbf{including the deployed showcase
engine}, trained under-regularized, the three recorded f0-vs-f1 gaps are not
like-regularized, and the cheapest improvement on the board is a single retrain.

(16) Three score-time campaigns then produced the arc's first genuine
\emph{holisticness} lever and its largest \emph{recall} gain, neither of them a
model change. The eval-time MMR $\lambda$ frontier (Section~\ref{sec:mmr}) is
confirmed natively on both anchors: every one of eight $\lambda<1$ arms lifts
holisticness@10 with the CI above zero, three single-GRU arms clear the prior
all-time bar \num{0.024872} at a recall tie (best \num{0.030355} at
$\lambda\num{0.7}$, pool~\num{50}), and the gain is legitimate rather than
composite-gaming because \emph{three} of the four crown factors --- diversity,
eagerness and \texttt{music@10} --- improve together. Its pre-registered mechanism
question resolves unanimously to the responsive hypothesis (8 of 8 arms,
$|\Delta\texttt{artist\_conc}|/|\Delta\texttt{artist\_adj}|$ between \num{1.59} and
\num{3.06}), so same-artist tracks \emph{are} near-neighbours in the content
metric, sonic MMR \emph{is} a de-eagering lever, and the explicit artist-penalty
sequel is cancelled. Two by-products are durable: the binding facet
\emph{switches} from concentration to adjacency as $\lambda$ falls, which locates
the end of the lever; and the shared-weights pairing is $\numrange{3}{33}\times$
sharper than the architecture-family noise floor (half-widths from
\num{0.000382}), so eval-time levers must budget their minimum detectable effect
from the pairing, not from Section~\ref{sec:power-stats}'s figure.
(17) The largest single recall gain in the whole arc is then \texttt{markov\_gate}
(Section~\ref{sec:markovgate}): make the Markov leg \emph{abstain} where the last
prefix item has zero train-bigram support --- \SI{55.6}{\percent} of queries --- and
Recall@10 rises to \num{0.23201} (\num{332} of \num{1431}), paired
$+\num{0.02027}$ $[+\num{0.01048}, +\num{0.03075}]$ canonically and
$+\num{0.01957}$ on a disjoint \SI{64}{\percent}-cold split, MRR up on both. The
defect is the blend's \emph{scale-invariant} normalizer, which re-inflates an
evidence-free artist/genre back-off to a full one-third vote; the fix has zero
learned parameters and nine lines. The bucket decomposition is the argument: the
whole effect sits in the zero-support bucket while both warm buckets are
\textbf{byte-for-byte identical} on three independent control/arm pairs, which
mechanically proves no global static leg weight could express it (dropping the leg
globally is a $-\num{0.0070}$ tie, the two buckets cancelling) and therefore
\emph{narrows} the four-times-confirmed learned-combiner pitfall to \emph{learned
score-level fusion} rather than per-query evidence gating. Unpredicted, it is also
a crown lever --- all four factors move the right way at once, and it is the first
lever that de-eagers the champion \emph{without} paying mood coherence, because it
removes a vote that had no business being cast rather than fighting the model's own
scores. Composed with MMR $\lambda\num{0.7}$ it is recall-neutral against the
champion ($+\num{0.00140}$, straddling) at holisticness \num{0.030980},
$\num{3.6}\times$ the champion's. Nothing was promoted: this is a recommendation,
and the one residual concern is that the zero-support \emph{boundary} was chosen
with test knowledge.
(18) Against those two positives, the training-time anti-eager lever is refuted
with a mechanism, and the refutation is expressive rather than dosimetric
(Section~\ref{sec:eagermargin}). The margins were calibrated from the
item$\to$item transition cosine (p75 \num{0.584}, p90 \num{0.830}) but the hinge
acts on prediction$\to$current, whose ceiling over training lies between
\num{0.5843} and \num{0.70} --- so the calibrated margins fall \emph{outside the
support of the penalized quantity} and four arms are provably inert, reproducing
the control's loss to \num{16} significant digits at every epoch at $\beta$ up to
the schema cap. Where the knob fires at margin zero it de-eagers spectacularly and
destroys relevance in the same motion, because an InfoNCE prediction is a hedged
near-average and eagerness is a \emph{ranking} property that an absolute cosine
cannot express. The prior ``weak knob'' reading is sharpened rather than retired:
the entire usable range at margin zero lies in an unexplored factor of \num{15}
between $\beta\num{0.2}$ and $\beta\num{3}$.
(19) Two campaigns about \emph{instruments} rather than models bracket these.
A purpose-built rollout metric (Section~\ref{sec:rollout}) --- fade per legible
axis, calibrated on \num{2135} real train sessions, with the anti-gaming term
\emph{inside} the metric rather than bolted on as an external recall floor ---
returns a sixth architecture null in the deployment regime, which is the strongest
form available because the standing rebuttal to the earlier five was that Recall@10
cannot see session quality. Its durable finding is that \emph{per-seed} variance
dominates model identity ($\pm\num{0.3}$ per seed against a model-level difference
of $+\num{0.004}$), so budget $\ge\num{40}$ seeds for any rollout comparison and
prefer serving an ensemble; and its three retracted positives are a catalogue of
how a generation metric goes wrong (a cross-split comparison read as architecture,
mood-extreme seed selection inflating effects ${\sim}\num{2.5}\times$, and an
anti-gaming term normalized by band width that scored a flatline \emph{above} a
real generator). A research memo (Section~\ref{sec:memo}) is reported at a lower
register --- nothing in it was run on the server, and its verified and provisional
numbers are kept partitioned --- but its convergent thesis, that every proposal
surviving adversarial screening is an \emph{evidence-conditioning} change rather
than a capacity change, is what points at (17), and its graveyard is cheaper to
reuse than to rediscover.
(20) Finally, a stage-stack depth campaign (Section~\ref{sec:stagestack}) closes
the \emph{serial} branch of the architecture axis to the two-surface standard the
tower-bank campaign then inherits. Five arrangements of per-step MLP and recurrent
stages, capacity-equalised at ${\sim}\num{1282000}$ parameters, all fall
\emph{through} the mandatory recall floor (\numrange{0.0706}{0.0860} against
\num{0.108}) and are refuted again on the walk at the matched cell, while a
$\num{3.24}\times$ capacity control ties on both surfaces --- so the damage is depth
itself, not capacity, not regularization (every deep arm trains \emph{longer} and
two reach a \emph{lower} validation loss) and not the MLP (a leading per-step MLP
is recall-neutral and holisticness-positive at depth 4, which \emph{narrows} the
$N=1$ pre-encoder result to the shallow case). Two facts outlive the null:
\emph{bracketing beats interleaving} at byte-identical capacity
($+\num{0.01537}$, CI above zero), and the training objective \textbf{dissociates}
from retrieval --- the arm with the best validation InfoNCE in the batch retrieves
\num{55} fewer hits than the control, so any architecture search selecting on
validation loss on this family will select the wrong arm.

\section{Source reports}
This synthesis consolidates twenty-two next-track source documents from
\texttt{experiments/} --- twenty-one campaign reports and one research memo, the
latter flagged as such because nothing in it was run on the server:
\begin{itemize}
  \item \texttt{2026-07-13-nexttrack-sequence-sweep.md} --- the founding
    next-track campaign: a GRU/LSTM sequence sweep against the first-order Markov
    bar; the best GRU-infonce ties on R@10 but loses on MRR
    (Section~\ref{sec:sweep}).
  \item \texttt{2026-07-13-nexttrack-topology-scan.md} --- the depth / width /
    residual / bidirectional topology scan; architecture is a spent lever, all
    arms cap below the Markov bar (Section~\ref{sec:topology}).
  \item \texttt{2026-07-14-nexttrack-blend-and-ensemble.md} --- the
    GRU$\times$Markov R+M z-blend beats the bar $+61\%$ R@10 (the first
    next-track champion); no richer ensemble improves on it and the learned XGB
    stacker underperforms (Section~\ref{sec:blend}).
  \item \texttt{2026-07-14-nexttrack-graded-relevance-and-new-legs.md} ---
    first-class graded-relevance metrics (artist@k / genre@k / artist\_mrr) and
    new leak-free base learners; content-kNN is the only additive signal, tying
    the champion on AE-64 (Section~\ref{sec:graded}).
  \item \texttt{2026-07-15-nexttrack-compression-representation-bakeoff.md} ---
    the item-representation compression bake-off (PCA vs autoencoder $\times$
    latent dim 32/64/128) under a frozen GRU readout; PCA-128 is the single-GRU
    winner but the win does not carry into the blend (Section~\ref{sec:compression}).
  \item \texttt{2026-07-15b-nexttrack-compression-dim-extension.md} --- the
    PCA/AE dimension sweep extended to 192/256; the curve peaks at PCA-192 and
    turns over, so PCA-128 stays the Phase-2 base (Section~\ref{sec:dim}).
  \item \texttt{2026-07-15c-nexttrack-phase2-model-sweep.md} --- the model-family
    sweep on the rich PCA spaces; the fixed z-blend beats the learned stacker,
    LSTM is the best single model, and R+M+C on PCA-192 is crowned the new
    best-on-record row (Section~\ref{sec:phase2}).
  \item \texttt{2026-07-18-nexttrack-literature-fit-campaign.md} --- a
    bibliography-driven exploration; a supervised learned content projection
    (\Rp+M+\Cp) is the new best-on-record row (R@10 \num{0.2117}, paired $\Delta$
    over R+M+C $+\num{0.0259}$ CI$>$0, seed-robust, cold-reaching), and every
    borrowed method (session-kNN, rank fusion, learned combiner, Markov smoothing,
    GRU importance negatives) closes (Section~\ref{sec:projection}).
  \item \texttt{2026-07-18b-nexttrack-projection-registration-and-scan.md} --- the
    projection reproduced on-server, registered $+$ promoted $+$ servable as
    \texttt{blend-gru-markov-content-proj} (best-models rank 1); an objective / rank
    / $\tau$ scan (InfoNCE $\gg$ BPR; default held) and second-split hardening on a
    disjoint 0.64-cold split (paired $\Delta$ $+\num{0.0426}$ CI$>$0)
    (Section~\ref{sec:projection-scan}).
  \item \texttt{2026-07-18c-nexttrack-representation-exploration.md} --- an
    item-vector representation audit; the equal-weighted PCA-192 geometry handicaps
    the raw GRU (six of eight alt spaces lift it CI$>$0) but the gain is absorbed by
    the projection (no space beats the champion on readout A)
    (Section~\ref{sec:representation}).
  \item \texttt{2026-07-20-musical-distance-metric.md} --- the continuous
    ``sounds-alike'' music@10 metric and its 53-run offline backfill; it correlates
    with Recall@10 ($r=\num{0.609}$) but crowns a different winner, the
    content-regressing \texttt{seq-nexttrack} (Section~\ref{sec:music}).
  \item \texttt{2026-07-22-session-holisticness-and-antieager-levers.md} ---
    a display-only session-holisticness metric suite (artist\_adj@10 /
    mood\_coh@10 / ild@10) and two anti-eager levers (eval-time MMR
    \texttt{mmr\_lambda}, training-time \texttt{eager\_beta}); the champion is
    the most artist-eager model, the eagerness is sourced from the Markov leg,
    GRU\,$\approx$\,LSTM, and MMR $\lambda$ is the strong knob
    (Section~\ref{sec:holisticness}).
  \item \texttt{2026-07-25-nexttrack-dual-tower-fusion-scan.md} --- the
    dual-tower view-pairing $\times$ fusion-depth scan against a
    \SI{0.098}{\percent} parameter-matched single tower; \textbf{REFUTED} (the
    fourth learned-combiner null, and the one that closes the ``train the fusion
    end-to-end'' rebuttal), with the raw-latent view shown to subsume
    \texttt{delta}/\texttt{cummean}, the MLP fusion head shown actively harmful,
    single-GRU width closed on PCA-192, and holisticness/recall shown
    anti-correlated across the batch (Section~\ref{sec:dualtower}).
  \item \texttt{2026-07-26-nexttrack-pre-encoder-scan.md} --- the learned
    per-step pre-encoder scan (presence $\times$ placement $\times$ dose) with
    the decisive single-tower ablation; \textbf{REFUTED and actively harmful}
    ($-\num{0.0140}$ Recall@10, CI below zero), the rectifier identified
    analytically and then measured (affine map neutral at ${\sim}0.8\sigma$, ReLU
    ${\sim}5.9\sigma$, \texttt{frac\_exact\_zero} $=$
    \texttt{frac\_negative\_coords} $=$ \num{0.5327}), plus the measured power
    analysis and the three holisticness generations
    (Sections~\ref{sec:preencoder} and~\ref{sec:power}).
  \item \texttt{2026-07-26-nexttrack-mmr-lambda-frontier.md} --- the native
    eval-time MMR $\lambda$ frontier on both anchors; \textbf{CONFIRMED} as the
    first genuine holisticness lever on record (three TIER-1 ``free'' single-GRU
    arms above the prior all-time bar at a recall tie, best $H$ \num{0.030355}),
    with the pre-registered mechanism question resolved unanimously to the
    responsive hypothesis (8 of 8 arms, ratio \numrange{1.59}{3.06}), the binding
    facet shown to switch as $\lambda$ falls, the candidate-pool axis closed as a
    fine-tuning knob, and the shared-weights pairing measured
    $\numrange{3}{33}\times$ sharper than the architecture-family noise floor
    (Section~\ref{sec:mmr}).
  \item \texttt{2026-07-27-rollout-hold-metric.md} --- \texttt{hold@}$N$, a
    generation metric for the autoregressive rollout regime (per-axis fade
    calibrated on \num{2135} real train sessions, with the anti-gaming term inside
    the metric), plus five sweeps and \num{240} journeys; \textbf{REFUTED} as a
    model discriminator --- a sixth architecture null, obtained in the deployment
    regime --- with per-seed variance shown to dominate model identity and three
    retracted positives documented as a catalogue of generation-metric pitfalls
    (Section~\ref{sec:rollout}).
  \item \texttt{2026-07-27b-nexttrack-model-topology-research-memo.md} ---
    \textbf{a research memo, not a campaign report}: nothing was run on the
    server, nothing promoted, nothing registered, and its numbers come from
    offline screens over frozen checkpoints. Consolidated at a lower register,
    with its own verified/provisional partition preserved and its test-set
    screening caveat stated; its reusable content is the convergent thesis that
    every surviving proposal is an \emph{evidence-conditioning} change rather than
    a capacity change, the mechanical explanation of the learned-combiner anomaly,
    and the graveyard (Section~\ref{sec:memo}).
  \item \texttt{2026-07-27c-nexttrack-markov-gate-confirm.md} --- the native
    two-split confirm of \texttt{markov\_gate}; \textbf{CONFIRMED} and the largest
    single recall gain in the arc (R@10 \num{0.23201} $=$ \num{332} of \num{1431},
    paired $+\num{0.02027}$ canonically and $+\num{0.01957}$ on a disjoint
    \SI{64}{\percent}-cold split, MRR up on both), with the warm evidence buckets
    \emph{byte-for-byte identical}, an unpredicted crown gain in which all four
    factors move the right way at once, a first recall-neutral crown arm when
    composed with MMR, and the learned-combiner pitfall narrowed accordingly
    (Section~\ref{sec:markovgate}).
  \item \texttt{2026-07-30-eager-margin-calibrated-refutation.md} --- the
    training-time eagerness regularizer at margins calibrated from the measured
    transition tail; \textbf{REFUTED across the whole axis, with a mechanism}: the
    margins were drawn from the item$\to$item cosine while the hinge acts on
    prediction$\to$current, whose ceiling sits between \num{0.5843} and \num{0.70},
    so four arms are provably inert (loss traces bit-identical at $\beta$ up to the
    schema cap) and the arms that fire destroy relevance before they help
    (Section~\ref{sec:eagermargin}).
  \item \texttt{2026-07-30c-artist-cap-measurement.md} --- the \emph{hard
    per-artist cap}, the first de-eagering lever that works, and it works because
    it cannot be outscored: \texttt{artist\_conc@10}
    \num{0.692}$\to$\num{0.123} with \texttt{music@10} \emph{rising} (CI above
    zero) at every setting, hardened on a second disjoint 0.64-cold split,
    together with the refutations of the two soft levers it replaces
    (Section~\ref{sec:walk-controls}). The walk-level material of
    Section~\ref{sec:walk} --- the surface mismatch, the four walk metrics, the
    anchor $\times$ stride frontier, the matched-sample head-to-head and the
    dual-tower reversal --- is transcribed from the \texttt{PROJECT-FACTS.md}
    roll-up of 2026-07-31b and from the four measurement tools it cites
    (\texttt{tools/walk\_eval.py}, \texttt{walk\_frontier.py},
    \texttt{walk\_dualgru.py}, \texttt{walk\_headtohead.py}), whose module
    docstrings are the authority for how each metric is computed; that material
    has no campaign report of its own yet.
  \item \texttt{2026-08-01-stage-stack-topology.md} --- the stage-stack
    topology campaign (arrangements of per-step MLP $\times$ recurrence at
    equalised capacity, via the new \texttt{seq-stack} predictor);
    \textbf{REFUTED on both surfaces} and the first closure on this corpus to that
    standard: all five 4-stage arms fall \emph{through} the recall floor and are
    refuted again on the walk, the $\num{3.24}\times$ capacity control ties on
    both, the user's rectifier-position hypothesis is refuted, bracketing beats
    interleaving at byte-identical capacity, and the best-fitting arm in the batch
    retrieves \num{55} fewer hits than the control
    (Section~\ref{sec:stagestack}).
  \item \texttt{2026-08-02-tower-bank-parallel-towers.md} --- the
    parallel-tower-bank dose campaign ($N$ bare-GRU towers under one linear head,
    via the new \texttt{seq-bank} predictor), \textbf{CLOSED on both surfaces}:
    no arm wins, five of eight lose to a parameter-matched single GRU, four h211
    towers and one h512 GRU return the identical \num{186} hits (paired $\Delta$
    \num{0.000000}), the \texttt{cummean} dose makes the walk strictly worse,
    both walk trades are cell-local, and the batch's one CI-clear result is that
    dropout \num{0.1} beats dropout \num{0.0} at byte-identical parameters ---
    which implies every \texttt{fusion\_layers}${}=0$ arm on record, the deployed
    engine included, trained under-regularized (Section~\ref{sec:towerbank}).
\end{itemize}

\end{document}
